\bookpart{视频处理}{Video Processing}
\label{part:video-processing}

\BookPartGuide{%
连续画面增加了运动与历史信息，也要求结果在时间上保持正确对应。对象身份、活动区间和有效邻帧都需要结合信息到达条件判断，不能只从单帧质量推断视频结果可靠。

视频感知、跟踪与检索章研究属性证据的积累、轨迹和掩码的维护、活动区间以及片段匹配。视频复原、增强与压缩章研究跨帧信息的选择、对齐与传播，进一步说明空间和时间采样扩展，以及依赖参考帧的视频编码和解码。两章共同核对内容、时间连续性、等待范围与资源成本。}

\chapter{视频感知、跟踪与检索}
\label{chap:vision-video-recognition}

路口录像中的一辆汽车在十帧里出现，交通统计应计作一辆还是十辆？厨房录像中有人在杯子旁伸手，究竟是拿起还是放下？过了半小时，有人问“刚才那只杯子放到哪里了”，系统又该返回哪一段画面？\ref{chap:vision-image-recognition}章的单帧模型可以回答“这一帧里有什么、在哪里”，却回答不了这三个问题。车辆计数需要知道不同帧中的框是否属于同一个实例，拿放判断需要知道状态变化的先后，回答查询则需要把问题对应到过去某个可以回看的时间范围。

本章的学习目标是：给定一段视频和一项任务，能够说明怎样把不同时刻的观测组织成时间证据，选用合适的模型形成对象、身份、事件和检索结果，并检查这些结果实际依赖了哪些帧。路线沿着证据的依赖关系展开。最简单的时间证据是同一对象的多次观测，所以先讨论怎样选择和汇总它们；汇总的前提是这些观测确实属于同一对象，于是接着讨论跨帧对应与跟踪状态的维护；有了延续的对象状态，才能判断对象之间发生了什么活动、何时开始和结束；前面得到的帧、轨迹和片段表示又可以作为索引，供查询驱动的检索使用；最后按输出、错误链和历史依赖三个方面检验这些判断。

阅读本章需要两类准备。一类是第一部分的经典基础：\ref{chap:vision-video-foundations}章的时间采样、块匹配与经典光流、镜头分割与背景建模，\ref{chap:vision-image-features}章的局部特征、HOG、模板匹配与视觉词袋。它们给出了不经学习也能工作的做法，本章在每个任务处先回顾这些做法在什么条件下有效、在哪里失效，再讲学习方法改变了哪一步。另一类是学习模型：\ref{chap:vision-image-recognition}章的分类、检测、分割与图像检索，\ref{chap:vision-multimodal-models}章的图文对齐接口，以及\ref{sec:vision-ivm-video-understanding}节的视频主干（TimeSformer、VideoMAE、V-JEPA 2.1）。主干的结构与预训练在那里讲解，本章只讲它们在具体任务中的用法。SiamRPN++、RAFT、ByteTrack、XMem、TSN、I3D、ActionFormer、MS-TCN、CLIP4Clip等专用于某项视频任务的模型在本章完整讲解，并给出逐层结构表或配置表；表中数字取自各自的论文，论文没有给出的项目明确标出。以下汽车、拿杯和敲门序列中的分数、坐标与代价均为教学构造。

\section{视频感知、跟踪与检索的问题与挑战}
\label{sec:vision-video-recognition-challenges}
\label{sec:vision-video-information}

视频感知、跟踪与检索不是一个任务，而是一组输出形式不同的任务。它们的输入都是按时间排列的帧，有时附带首帧给定的框或掩码、逐帧的检测结果、一句查询或一个视频库；输出则从一个活动类别，到每帧一个框的轨迹、每帧一张掩码、一组带起止时间的事件，再到按相关程度排序的片段列表。在选择任何模型之前，需要先回答两个问题：任务要返回什么，以及模型在作出判断时能看到什么。前者决定标注、损失和评价对象，后者决定哪些帧可以作为证据。两个问题彼此独立：同一种输出可以在离线或在线条件下完成，同一种观测条件也可以服务多种输出。本节先规定输入、输出与评价指标，再讨论使这些任务变难的五类挑战；能看到什么，作为其中一类挑战单独讨论。

视频用按时间排列的帧$\bm{X}_1,\ldots,\bm{X}_n$表示，帧的拍摄时刻为$\tau_1,\ldots,\tau_n$。等帧率时，相邻时刻差为$1/f$秒，其中$f$是每秒帧数；可变帧率、丢帧或抽样后，应保存实际时间戳。帧编号相差十，只有在时间采样已知时才能换算成持续时间，\ref{sec:vision-videofound-frame-rate}节讨论了帧率不足时的时间混叠。

\label{sec:vision-video-outputs}
同一段拿杯视频可以产生多种标注。“有人拿杯子”是给定范围内的活动标签；每帧的人物框及共同身份组成框轨迹；杯子的逐帧像素区域组成掩码轨迹；从手接触杯子到杯子离开桌面的区间，则描述某种约定下的动作范围；“刚才拿杯子的那段”是一次查询的检索结果。这些输出分别回答“是什么”“是不是同一个”“何时发生”和“在哪里能找到”，标注成本与评价方式也各不相同。训练时只有视频标签，就不能把未提供的逐帧边界当作监督；评价时框位置正确，也不能替代身份或活动正确。

\term{目标跟踪}（object tracking）在连续观测中延续选定对象的位置或区域，使后续结果仍指向同一个实例。给定首帧汽车框，即使不知道它的车型，也可以尝试跟踪；自动发现所有车辆则还需要检测及新目标管理，这称为\term{多目标跟踪}（multi-object tracking, MOT）。跟踪返回的身份通常是序列内编号，不是现实世界中的车牌或人物姓名。

输出的时间范围也要事先规定。裁剪片段通常已经把目标活动置于有限范围，未裁剪视频还含等待、背景和多次重复。若训练只见到以动作为中心的短片段，使用时直接对整小时录像分类，输入与输出条件都发生了变化。镜头切换还可能使图像坐标失去连续性；跨镜头重新认出同一辆车需要额外身份依据。视频时间上的位移描述也不能直接当作带米和秒单位的三维运动。

表\ref{tab:vision-video-task-metrics}汇总了本章各任务的输入、输出与主要指标。表中每个指标都在下文给出定义。同一名称在不同基准中的细节可能不同，比较结果时必须说明所用协议。

\begin{table*}[tbp]
  \centering
  \small
  \begin{tabularx}{\linewidth}{@{}lXXX@{}}
    \toprule
    任务 & 输入 & 输出 & 主要指标\\
    \midrule
    对象与属性的时序识别 & 帧或对象裁剪序列 & 类别、属性及状态序列 & 准确率、前$k$准确率\\
    单目标跟踪 & 视频与首帧框 & 每帧一个框 & 成功率曲线下面积、中心距离精度\\
    多目标跟踪 & 视频与逐帧检测 & 带身份编号的框轨迹 & MOTA、IDF1、HOTA\\
    视频目标分割 & 视频与首帧掩码或提示 & 每帧的目标掩码 & $\mathcal{J}$、$\mathcal{F}$及其均值\\
    光流与点轨迹 & 两帧或整段视频 & 二维位移场或三维点 & 端点误差\\
    片段活动分类 & 已裁剪的片段 & 活动类别 & 准确率、前$k$准确率\\
    时间动作定位 & 未裁剪视频的片段特征 & 带类别的区间集合 & 多个tIoU门限下的mAP\\
    动作分割 & 未裁剪视频的片段特征 & 逐帧步骤标签 & 帧准确率、编辑分数、F1@$k$\\
    视频检索 & 查询与视频库 & 候选的排序列表 & Recall@$K$、中位排名\\
    未完成事件的预判 & 持续到达的帧与查询 & 带时刻的响应 & Timeliness-F1\\
    \bottomrule
  \end{tabularx}
  \caption{本章各任务的输入、输出与主要评价指标。指标还依赖观察协议：抽帧方式、初始化、可见的未来范围和人工纠正预算不同，同名指标也不能直接比较。}
  \label{tab:vision-video-task-metrics}
\end{table*}

视频级识别沿用\ref{sec:vision-recognition-challenges}节的准确率与前$k$准确率，只是样本单位变成一段视频或一个片段。同一视频读取多个片段时，还要说明片段预测怎样汇总成视频预测，否则两个“视频准确率”可能比较的是不同的读取协议。单目标跟踪逐帧比较预测框$\widehat{\bm{b}}_t$与真实框$\bm{b}_t$。按式\eqref{eq:vision-iou}计算第$t$帧的交并比$o_t$，对门限$\theta\in[0,1]$定义成功率$S(\theta)=(1/n)\sum_{t=1}^{n}\bm{1}[o_t>\theta]$，即重叠超过门限的帧所占的比例，$\bm{1}[\cdot]$在条件成立时取1。让$\theta$从0扫到1得到成功率曲线，曲线下面积等于各门限上成功率的平均；SiamRPN++在OTB2015上报告的就是这一面积\scite{vision-siamrpn2019}。另一种口径只看中心：预测中心与真实中心的距离不超过$\delta$像素的帧所占比例称为精度，KCF报告$\delta=20$像素处的值\scite{vision-kcf2015}。成功率同时检查位置与尺度，中心精度不检查尺度。VOT系列基准还在跟丢后重新初始化，并计算期望平均重叠（expected average overlap, EAO），其协议细节不在此展开。

多目标跟踪既要框对，也要身份对，\ref{sec:vision-videofound-identity}节已经引入身份切换与轨迹碎片，并给出了MOTA与IDF1，这里用本章的记号重述，再补上HOTA。\term{CLEAR MOT指标}（CLEAR MOT metrics）逐帧建立预测与真实对象之间的对应，并优先沿用上一帧的对应，以便发现身份切换\scite{vision-bernardin2008clear}。令$\operatorname{FN}_t$、$\operatorname{FP}_t$和$\operatorname{IDSW}_t$分别为第$t$帧的漏检数、误报数和身份切换数，$\operatorname{GT}_t$为该帧的真实对象数，多目标跟踪准确度为
\begin{equation}
  \operatorname{MOTA}=1-\frac{\sum_t(\operatorname{FN}_t+\operatorname{FP}_t+\operatorname{IDSW}_t)}{\sum_t \operatorname{GT}_t}\text{。}
  \label{eq:vision-video-mota}
\end{equation}
三类错误直接相加，误报总数又可以超过真实对象总数，所以MOTA可以为负。漏检和误报的数量通常远多于身份切换，ByteTrack的作者因此指出MOTA更多反映检测质量\scite{vision-bytetrack2022}。身份F1（IDF1）改为在整段视频上作全局的一对一对应：每条真实轨迹至多与一条预测轨迹配对，再按满足身份及位置条件的观测数计算\scite{vision-idf1-2016}
\begin{equation}
  \operatorname{IDF1}
  =\frac{2\,\operatorname{IDTP}}
  {2\,\operatorname{IDTP}+\operatorname{IDFP}+\operatorname{IDFN}}\text{，}
  \label{eq:vision-video-idf1}
\end{equation}
其中$\operatorname{IDTP}$为身份正确的匹配观测数，$\operatorname{IDFP}$和$\operatorname{IDFN}$分别是身份误报和漏报的观测数。

\term{高阶跟踪准确度}（higher order tracking accuracy, HOTA）把检测与关联拆成两个可以分别读取的分量\scite{vision-hota2021}。先固定定位门限$\alpha$，逐帧一对一匹配相似度不低于$\alpha$的预测与真实检测，得到真阳性集合TP、漏检集合FN和误报集合FP。对一个真阳性$c$，再看整段视频：与$c$具有相同预测编号和相同真实编号的真阳性构成TPA$(c)$；真实编号相同、但预测编号不同或漏检的检测构成FNA$(c)$；预测编号相同、但真实编号不同或为误报的检测构成FPA$(c)$。于是
\begin{equation}
  \begin{gathered}
    A(c)=\frac{|\mathrm{TPA}(c)|}{|\mathrm{TPA}(c)|+|\mathrm{FNA}(c)|+|\mathrm{FPA}(c)|},\\
    \mathrm{DetA}_\alpha=\frac{|\mathrm{TP}|}{|\mathrm{TP}|+|\mathrm{FN}|+|\mathrm{FP}|},\qquad
    \mathrm{AssA}_\alpha=\frac{1}{|\mathrm{TP}|}\sum_{c\in\mathrm{TP}}A(c)\text{，}
  \end{gathered}
  \label{eq:vision-video-hota-parts}
\end{equation}
并取$\mathrm{HOTA}_\alpha=\sqrt{\mathrm{DetA}_\alpha\cdot\mathrm{AssA}_\alpha}$。$A(c)$回答“与$c$同属一条真实轨迹的观测里，有多大比例被分到了同一个预测编号”，DetA是检测层面的交并比，AssA是关联层面的交并比。最后在$\alpha=0.05,0.10,\ldots,0.95$这19个门限上对$\mathrm{HOTA}_\alpha$取平均。

用一个算例比较三者。某辆车出现在十帧中，十个框都与真值完全重合，但跟踪器在第7帧把编号从$a$换成了$b$。按式\eqref{eq:vision-video-mota}，没有漏检和误报，只有一次身份切换，MOTA为$1-1/10=0.9$。按式\eqref{eq:vision-video-idf1}，全局一对一对应只能让真实轨迹配上$a$或$b$之一，选较长的$a$，得到$\operatorname{IDTP}=6$、$\operatorname{IDFP}=4$、$\operatorname{IDFN}=4$，IDF1为$12/20=0.6$。按式\eqref{eq:vision-video-hota-parts}，各门限下DetA都为1；前六帧每个真阳性的TPA为6、FNA为4，$A=0.6$，后四帧的$A=4/10=0.4$，AssA为$(6\times0.6+4\times0.4)/10=0.52$，HOTA约为$\sqrt{0.52}\approx0.72$。同一个错误在三个指标中的分量相差很大：MOTA只把它记作一次切换，IDF1惩罚整段被分走的观测，HOTA把检测与关联分开后再合并。报告时应同时给出三者，以及身份切换次数和轨迹中断次数，这两个概念见\ref{sec:vision-videofound-identity}节。

视频目标分割的指标沿用DAVIS基准的两项度量\scite{vision-davis2017}。区域相似度$\mathcal{J}$就是预测掩码与真实掩码的交并比；轮廓精度$\mathcal{F}$先在两张掩码的边界像素之间作二分匹配，得到边界精确率$P_c$与召回率$R_c$，再取$\mathcal{F}=2P_cR_c/(P_c+R_c)$。DAVIS 2017先在每个对象的各帧上平均、再在所有对象上平均，分别得到两项的均值，最后取二者的平均作为$\mathcal{J}\&\mathcal{F}$。沿车身外扩两个像素的掩码，$\mathcal{J}$仍然很高，边界匹配却会明显变差；两项合用，才能同时检查面积和轮廓。光流常用\term{端点误差}（endpoint error, EPE）衡量：在有效位置集合$V$上取$(1/|V|)\sum_{\bm{p}\in V}\lVert\widehat{\bm{u}}(\bm{p})-\bm{u}^*(\bm{p})\rVert_2$，即预测位移与真实位移之差的平均欧氏长度，单位是像素。

事件区间用\term{时间交并比}（temporal intersection over union, tIoU）衡量范围的吻合。对两个正长度区间$[a,b]$和$[c,d]$，其中$a<b$、$c<d$，交集长度不能小于零，并集长度等于两段长度之和减去交集，故
\begin{equation}
  \operatorname{tIoU}
  =\frac{\max\{0,\min(b,d)-\max(a,c)\}}
  {(b-a)+(d-c)-\max\{0,\min(b,d)-\max(a,c)\}}\text{。}
  \label{eq:vision-video-tiou}
\end{equation}
$[9,17]$与$[10,18]$相交7秒、并集9秒，得分为$7/9$。时间动作定位沿用\ref{sec:vision-recognition-challenges}节检测AP的计算，只把框的IoU换成tIoU：在给定门限下，按得分从高到低把预测区间与同类、尚未匹配的真实区间配对，计算每一类的AP，再对类别平均得到mAP。THUMOS14通常报告tIoU从0.3到0.7、间隔0.1的五个门限及其平均，ActivityNet-1.3则用0.5到0.95、间隔0.05的十个门限\scite{vision-actionformer2022}。动作分割同时报告三项：帧准确率是逐帧标签正确的比例；编辑分数先把预测和真实各自压缩成片段的类别序列（如“拿、倒、放”），计算两序列的编辑距离$e$，常用实现报告$(1-e/\max(\ell_{\mathrm{pred}},\ell_{\mathrm{gt}}))\times100$，其中$\ell$是片段数；F1@$k$把与同类真实片段的tIoU不低于$k$的预测片段算作命中，MS-TCN报告$k$取10\%、25\%和50\%时的F1\scite{vision-mstcn2019}。帧准确率奖励长片段内部的正确，编辑分数与F1惩罚碎片化和顺序错误。检索沿用\ref{sec:vision-recognition-challenges}节的Recall@$K$；文字查视频时每个查询通常只有一个正确视频，Recall@$K$就是正确视频排进前$K$名的查询所占比例，中位排名则是正确视频名次的中位数，越小越好。未完成事件的响应还要检查时机，Timeliness-F1在\ref{sec:vision-video-proactive-response}节定义。

有了输入、输出和指标，就可以看清这些任务为什么难。下面五个小节各讨论一类挑战。它们不只属于某一个任务，而是在跟踪、事件判断和检索中以不同形式反复出现。

\subsection{运动与相机运动带来的对应困难}
\label{sec:vision-video-challenge-motion}

路口录像中，汽车每帧向右移动约30像素，相机同时向左摇动，路面上的所有像素也一起向右移动约10像素。要把第一帧车门上的一个像素与第二帧中的同一像素对应起来，至少有三处困难。其一，亮度未必守恒：车驶入树荫，同一表面点的颜色变暗。其二，局部信息不足：车门是一大片均匀的白色，小窗口内处处相像，只看窗口无法确定它沿门面滑动了多少，这就是\ref{sec:vision-videofound-motion}节的孔径问题。其三，位移的大小跨度很大：远处行人每帧只动一两个像素，近处汽车动几十个像素，相机运动又给每个像素叠加了一份与物体无关的位移。对应一旦出错，后面的轨迹汇总、身份关联和事件判断都会建立在错误的“同一对象”之上。

经典方法从两个方向处理这一困难。\ref{sec:vision-videofound-block-matching}节的块匹配在搜索窗口内逐个比较图像块，用绝对差之和选出最相似的位移；它不需要线性化，却只能找到搜索范围内的整数位移，均匀区域还会给出许多同样好的候选。\ref{sec:vision-videofound-horn-schunck}节的Horn--Schunck方法与\ref{sec:vision-videofound-lucas-kanade}节的Lucas--Kanade方法从亮度守恒出发，在小位移下把它线性化为光流约束方程。前者用全局平滑项补足孔径问题缺少的约束；后者假设小窗口内位移相同，用最小二乘求解，并借助图像金字塔处理较大的位移。相机运动则可以先用一个全局单应近似、再从光流中扣除，改进的密集轨迹正是这样做的\scite{vision-idt2013}。这些方法的假设很明确：亮度守恒，位移小或可由金字塔分解，运动场平滑或局部恒定。遮挡边界、快速移动的小物体和光照变化恰好会同时违反它们。

学习方法保留了“比较、搜索、迭代修正”的结构，但把比较对象从原始亮度换成学到的特征，把手工设定的平滑项换成从数据学到的更新规则。RAFT在学到的特征上计算所有像素对的相似度，不再依赖线性化，也不再逐级细化；D4RT更进一步，在共同的三维坐标中回答“某个表面点在某时刻位于何处”，从而把相机运动与物体运动分开。\ref{sec:vision-video-correspondence}节按像素、对象和表面点三种粒度回答这一挑战。

\subsection{遮挡与外观变化带来的身份维护困难}
\label{sec:vision-video-challenge-identity}

两辆外观相近的白色轿车A、B在路口交错，A被驶过的公交车挡住约两秒，再出现时已经转过半个弯，车身由侧面变成了背面。跟踪器要在每帧回答“当前这个框延续哪条轨迹”，而这里同时出现了三种情形：观测暂时消失；目标外观变化很大，与首帧模板已不相像；附近恰好有一个与之相像的干扰者。只看外观，会在A、B交错时换号；只看运动，会在A转弯时跟丢；遮挡期间若继续用新观测更新目标模型，还可能把公交车的外观写进A的模型。

经典方法给出了几种不同的目标模型。\ref{sec:vision-features-template-matching}节的模板匹配用归一化互相关比较像素，能抵消整体的亮度变化，却对形变和视角无能为力。均值漂移跟踪改用颜色直方图描述目标，不要求像素逐点对齐，因此能容忍非刚性形变，但颜色相近的干扰者会把它拉走。相关滤波在首帧及后续各帧上学习一个判别式的滤波器，并按固定比例在线更新，能适应缓慢的外观变化，却会在遮挡期间逐渐吸收错误的内容。多目标场景中，卡尔曼滤波预测每条轨迹的位置，匈牙利算法求一对一分配；相机静止时，还可以用\ref{sec:vision-videofound-background}节的背景建模先分出前景区域。这些做法的共同局限是：目标模型的形式由人选定，遮挡期间是否更新也只能靠手设的门限判断。

学习方法从三个方面改变了这一局面。孪生网络在大量视频帧对上学习“两块图像是否属于同一目标”的相似度，不再只依赖颜色或亮度；检测器驱动的关联用检测分数和轨迹历史共同决定低分框的去留；记忆式分割把多个时刻的外观保存为可查询的键值，并显式判断目标是否出现在当前帧。\ref{sec:vision-video-initialization}节比较单目标跟踪的三种目标模型，\ref{sec:vision-video-association}节讨论多目标关联，\ref{sec:vision-video-state-update}节与\ref{sec:vision-video-recovery}节讨论状态的写入与遮挡后的找回。

\subsection{时间尺度跨度带来的事件判断困难}
\label{sec:vision-video-challenge-timescale}

厨房录像中，敲门不到一秒，拿杯约三秒，倒水十几秒，从取杯到洗杯的完整流程则持续几分钟。同一段视频里，事件的持续时间可以相差两个数量级；拿起与放下又只在顺序上不同，单看任何一帧都分不清。模型读取的时间范围太短，就看不到“杯子离开桌面”这一结果；读取范围长而采样稀疏，又会漏掉不到一秒的敲门。事件的起止本身也有歧义：手刚伸出时是否已经开始拿杯，不同标注者可能相差半秒。

经典方法先把视频切成可处理的单元，再在单元内描述运动。\ref{sec:vision-videofound-shot-boundary}节的镜头边界检测与\ref{sec:vision-videofound-keyframes}节的关键帧选择把长视频分成镜头和代表帧，供浏览与索引使用。对活动识别，手工特征路线沿稠密轨迹提取外观与运动直方图，再把整段视频编码为一个固定长度的向量，交给线性分类器\scite{vision-idt2013}；区间定位则在多种窗口长度上滑动这样的分类器。这条路线能描述几帧到十几帧内的局部运动，却要为每种窗口长度单独计算，编码时又把所有局部特征混在一起，丢掉了先后顺序。

学习方法分别从时间覆盖、局部运动和多尺度三个方向回应。TSN以分段稀疏采样覆盖整段视频，I3D用三维卷积直接从连续帧学习局部运动，ActionFormer在多尺度的时间特征金字塔上让不同层负责不同长度的事件，MS-TCN用逐层加倍的空洞率扩大时间感受范围，并逐阶段修正碎片化的预测。\ref{sec:vision-video-activities}节按输出类型讨论这些模型。

\subsection{在线因果与延迟约束带来的响应困难}
\label{sec:vision-video-availability}

能看到什么，由两个条件共同决定。信息到达方式区分离线与在线：离线处理读取已经给定的完整记录，在线处理在时刻$\tau$只能使用已经到达的帧。处理范围区分局部片段与长时历史：前者围绕有限的时间窗口形成证据，后者需要选择、压缩和索引过去的观测。这两个维度可以组合，离线算法可以只看五帧，在线算法也可以保存数小时的历史。若允许等待$\delta$秒后输出，应把新增的后续帧和等待时间一起记录，具体范围见图\ref{fig:vision-observation-window}。

在线条件带来的困难是具体的。路口计数系统以每秒25帧运行，每帧只有40毫秒的处理预算；厨房助手要在手碰到热水壶之前发出提醒，而“碰到”这一结果此时还没有被拍到。离线分析可以读完整段再回头修订，在线系统发出的判断却无法用后来的画面改写。许多离线效果最好的模型恰恰依赖未来帧：区间定位读取事件两侧的上下文，非因果的时间卷积在每个位置都读取后续时刻，离线分割可以从后往前再传播一遍。

经典方法大多天然是因果的递推。卡尔曼滤波只用过去的观测更新状态，\ref{sec:vision-videofound-background}节的背景模型按固定学习率逐帧更新，相关滤波每帧只做几次快速傅里叶变换，原文报告的速度达到每秒669帧\scite{vision-mosse2010}。它们每帧的计算量固定，代价是模型容量有限，无法表达复杂的外观与事件。学习方法要在线使用，就必须改写读取范围：跟踪与分割模型逐帧处理并维护流式记忆，事件模型要区分“预测的结束时刻”与“已确认的结束时刻”，响应模块还要显式决定何时发言。\ref{sec:vision-video-tracking}节的跟踪模型都按帧顺序处理，\ref{sec:vision-video-proactive-response}节讨论未完成事件的预判，\ref{sec:vision-video-output-protocol}节要求在评价中记录每个输出当时可见的范围。

\subsection{长视频规模带来的记忆与检索困难}
\label{sec:vision-video-challenge-memory}

一小时、每秒30帧的录像有108000帧。若像SAM 2那样为每帧保存一张$256\times64\times64$的图像特征，按每个数4字节计，每帧约4.2 MB，一小时约450 GB；这只是说明规模的教学估算，SAM 2自身正是为避免这种增长而只保留少量记忆帧。保存每帧全部表示的存储随时间线性增长，而查询又可能指向其中任意一刻。

“刚才把钥匙放在哪里”是Ego4D第一视角经历查询所关注的问题类型，需要把问题对应到过去的画面\scite{vision-ego4d2022}。这类查询要求系统知道历史中哪些细节值得保存，第一视角的遮挡、相机运动和可见范围又会影响证据是否存在。部分样本有其他传感信息，使用时应逐样本确认，不能统一假设视频都有完整的音频或视线记录。MovieChat用长短期记忆压缩密集的视觉信息，区分全局问题与针对某一时刻的问题\scite{vision-moviechat2024}。压缩后若只剩“有人开车”，后来就难以从这句话恢复车牌字符。Flash-VStream进一步面向持续到达的视频维护概括信息与细节记忆，响应成本要连同帧率、历史长度、硬件和提问时刻一起解释\scite{vision-flash-vstream2025}。这些系统所用的视觉语言连接已在\ref{chap:vision-multimodal-models}章建立。历史中存在过某个画面，不等于回答时仍能访问它。

经典方法用结构化的索引控制规模。镜头分割与关键帧把一小时录像压缩成几百幅代表帧；\ref{sec:vision-features-bovw}节的视觉词袋把局部特征量化成视觉词，用倒排表记录每个词出现在哪些帧，查询时只访问含有相同词的帧。代价是代表帧之间的短暂事件可能根本没有进入索引，词频统计又丢掉了先后顺序。学习方法提供了三种取舍：XMem把常被读取的历史压缩成少量原型，存储有界但会丢掉细节；CLIP4Clip为每段视频预先算好一个向量，检索很快但时间细节被平均掉；TRACE不预先压缩，而是按问题回到原始画面中取证，细节保留得最多，每次查询的代价也最高。\ref{sec:vision-video-state-update}节和\ref{sec:vision-video-retrieval}节分别讨论跟踪中的记忆与查询驱动的检索。

五类挑战按证据的依赖顺序在后续各节中得到回答。\ref{sec:vision-video-object-recognition}节处理最简单的时间证据，即同一对象的多次观测怎样汇总；汇总的前提已经是对应可靠，长视频中证据分散的问题也在这里第一次出现。\ref{sec:vision-video-tracking}节正面回答对应与身份两类挑战：先比较单目标跟踪的三种目标模型，再按像素、对象和表面点三种粒度建立跨帧对应，最后讨论记忆式的状态更新与遮挡后的找回，在线因果约束贯穿其中。\ref{sec:vision-video-activities}节在对象状态之上回答时间尺度的挑战，并在未完成事件的预判中再次面对在线约束。\ref{sec:vision-video-retrieval}节把前面形成的帧、轨迹和片段表示作为索引，回答长视频的记忆与检索挑战。每一处都先回顾经典路线的做法与失效条件，再讲学习模型改变了哪一步；最后一节把这些输出放在同一观察协议下检验。

\section{对象与属性的时序识别}
\label{sec:vision-video-object-recognition}

本节要解决的问题是：给定一段视频中某个对象的多次观测，判断它的类别和属性，例如这辆车是轿车还是货车、车身是什么颜色、车灯是否点亮。输入是若干帧或对象裁剪，输出是指定对象在指定时间范围内的类别或属性，变化属性还要输出随时间变化的状态序列。难点在于证据分散：一辆车穿过画面时，车身、车牌和车灯的最佳观测可能出现在不同帧，其余帧可能模糊、被遮挡或混入另一辆车。直接把所有帧平均起来，会把这些坏观测也加进来。

解决这个问题需要依次完成三步。先决定读取哪些帧和区域，因为没有被采到的证据，后续任何汇总都无法恢复；再决定怎样把多次观测合成一个判断，这一步的关键选择是在得分、对象轨迹还是密集特征的层次上合并；最后，对会随时间变化的属性，还要决定何时改变输出状态，而不是只给一个整段平均值。三步的顺序对应证据从“采到”到“合并”再到“随时间解释”的流程。

\subsection{观测证据的选择}
\label{sec:vision-video-evidence-selection}

若相机以每秒30帧记录，一辆车只在0.3秒内经过，即只出现在约9帧中。用每秒一帧的等间隔抽样，可能一帧也采不到它；把抽样起点推迟半秒，结果又可能改变。提高采样率能增加覆盖，同时增加检测和编码开销。一种折中是先用较粗的时间索引找出运动或对象出现的范围，再在这些范围内加密读取。这样的筛选仍要检查漏选，不能只评价已经找到的片段。

读取的空间范围同样需要选择。整帧表示适合利用场景上下文，对象裁剪更集中于车身或服饰。裁剪之前若没有可信的框，识别器可能读到路面或另一辆车；相机运动和出画面还会使固定裁剪失效。保存帧号、框坐标和清晰度，才能在识别出错时查明模型看到了什么。车牌需要的分辨率与车型不同，不能仅以“已采到一帧”判断证据充分。

\subsection{跨帧证据的汇总}
\label{sec:vision-video-aggregation}

选好观测之后，同一辆车有了$T$次观测，需要合成一个判断。这一步的核心难点是：每次观测的质量不同，观测之间是否属于同一对象也未必确定，而不同任务需要保留的信息量又不一样。对这个难点有三种回答，它们的差别在于“在哪个层次上合并”。

第一种在得分层次合并：先让单帧分类器为每次观测输出类别得分，再对这些得分求加权平均。它不需要额外训练，也不要求跨帧身份，只要求调用者事先保证这些观测来自同一对象或同一给定范围；代价是空间位置和时间顺序在合并时被丢弃。第二种在对象轨迹上合并：借助跟踪器给出的身份，为每个身份维护一份随时间更新的属性统计。它能在线更新、能区分同一画面中的多辆车，但完全依赖身份关联是否可靠。第三种在特征层次上合并：不先压成类别得分，而是保留每个时空网格位置的特征向量，再由任务预测头读取。它能回答局部深度、短时预期等得分汇总回答不了的问题，代价是需要预训练的视频主干和另外训练的预测头。下面三小节分别展开这三种做法，第三种以V-JEPA 2.1的特征为例。

\subsubsection{基于逐帧得分的汇总}
\label{sec:vision-video-weighted-scores}

设已经确认若干裁剪属于同一辆车，第$t$次观测对类别$c$给出得分$s_{t,c}$。这些得分可以是分类器的原始输出，也可以是归一化后的类别概率；具体选择应在实验中固定。以下用归一化输出演示算术，但并不假定这些数值已经校准。

平均汇总让每次观测同等参与，最大值保留最强的单次响应，可靠性加权则让清晰、完整的观测有更大影响。对非负可靠性$r_t$，只要$\sum_t r_t>0$，加权汇总可写为
\begin{equation}
  \bar{s}_c=\frac{\sum_t r_t s_{t,c}}{\sum_t r_t},
  \qquad
  \widehat{c}=\argmax_c\bar{s}_c\text{。}
  \label{eq:vision-video-weighted-scores}
\end{equation}
分母使权重和为1，避免同一组相对权重整体放大后改变得分尺度。没有可靠观测时应返回未知或重新取帧，不能把分母为零的情形悄然当作一组零分。

考虑清晰、局部遮挡和严重模糊的三帧。它们对“轿车、货车”的输出分别为$(0.80,0.20)$、$(0.60,0.40)$和$(0.05,0.95)$，可靠性设为$1,0.5,0.1$。均值为$(0.4833,0.5167)$，最大值为$(0.80,0.95)$，二者都选择货车。加权轿车分数却为
\[
  \frac{1\times0.80+0.5\times0.60+0.1\times0.05}{1.6}
  =0.690625\text{，}
\]
货车分数为$0.309375$，结果改为轿车。这说明加权可以限制模糊误判的影响，但并不证明加权一定改善识别：如果可靠性规则把唯一有效的帧误判为不可靠，结果也会变差。

最大值还可能由一次异常高分主导；均值则可能稀释一帧中短暂可见的车牌或标志。人工设定的清晰度规则与通过标注学习的权重是两种不同方案。学习方案要交代训练目标、标注来源和验证划分，不能在测试结果上反复挑选有利的权重。

\subsubsection{基于对象轨迹的汇总}
\label{sec:vision-video-track-evidence}

逐帧得分汇总要求调用者事先给出“哪些观测属于同一辆车”。当画面中有多辆车、观测又持续到达时，这个前提需要由跟踪器提供。跟踪器在每帧返回身份及位置后，区域识别器就能为每个身份单独积累证据。只在关联与可见性可信时，才把当前得分写入该身份的历史。积累方式有两种：滑动窗口保留最近$k$次观测，窗外观测完全移除；指数衰减则逐步降低旧证据权重，不必保存全部历史分数。

令$\bm{s}_t$为当前类别得分向量，$r_t$同时表达观测与关联的可靠性，$0<\rho<1$是衰减因子。维护分子$\bm{a}_t$和总权重$b_t$：
\begin{equation}
  \bm{a}_t=\rho\bm{a}_{t-1}+r_t\bm{s}_t,\qquad
  b_t=\rho b_{t-1}+r_t,\qquad
  \bar{\bm{s}}_t=\bm{a}_t/b_t\text{。}
  \label{eq:vision-video-decay}
\end{equation}
初始时二者都为零，仅在$b_t>0$时输出。展开递推可得$\bm{a}_t=\sum_{k\leq t}\rho^{t-k}r_k\bm{s}_k$，所以$\bar{\bm{s}}_t$正是以$\rho^{t-k}r_k$为权重的式\eqref{eq:vision-video-weighted-scores}，越旧的观测权重越小。无新观测时可令$r_t=0$；这时分子、分母同比衰减，旧判断不会凭空改变，但应单独标出证据已经陈旧。采样间隔不等时，可按实际时间差设置$\rho=\exp(-\Delta\tau/\kappa)$，其中$\kappa$是以秒计的衰减时间尺度。

以车身颜色为例，过去的红色得分为0.8，新一次可靠观测为0.6，令旧总权重为1、$\rho=0.5$、新权重为1，更新后红色分数为$(0.5\times0.8+0.6)/1.5=2/3$。更新的是统计状态，分类器参数并未重新训练。若遮挡后错误接入另一辆红车，历史积累可能反而加强身份错误；\ref{sec:vision-video-tracking}节的关联可信性决定这里能否安全汇总。

\subsubsection{基于密集时空表示的汇总}
\label{sec:vision-video-dense-state}

一辆车的类别已经确定，并不表示系统还保留了车轮的位置、局部深度或车灯即将发生的变化。前两种汇总都把许多位置压成少数得分；若希望分别判断这些局部状态，就需要先留下带空间位置和时刻的特征，再让任务预测头读取它们。这里的密集表示是每个时空网格位置都有一个特征向量，而不是每个位置已经有正确的语义或几何答案。

V-JEPA 2.1提供了这样的特征，其结构与训练见\ref{sec:vision-ivm-vjepa}节。这里只回顾用法所需的结论：编码器把视频切成时空小块，为每个小块输出一个$d$维向量；预训练让预测器从可见小块估计教师编码器在遮蔽位置和可见位置上的特征，并在多个中间层施加这一约束，因此各位置的向量保留了局部信息，而不只是替全局类别服务\scite{vision-vjepa21-2026}。预训练不使用深度或动作标签，也没有把特征解码回像素的步骤。

使用时，先按预训练时的规则组织输入。设输入为$F$帧、每帧$H\times W$，时空小块大小为$2\times16\times16$，编码器输出$\bm{Z}\in\mathbb{R}^{(F/2)\times(H/16)\times(W/16)\times d}$。例如教学设定16帧、$256\times256$的片段，得到$8\times16\times16=2048$个位置向量。帧数、帧间隔和分辨率应与所用权重的训练阶段一致；论文主训练使用16帧片段，后续较高分辨率阶段使用更长片段，换用某个发布权重时要核对这些条件。推理只需编码器，不运行遮蔽与教师分支。

读取这些特征的方式取决于输出类型，原论文的迁移协议都冻结编码器，只训练新增的任务模块。对整段类别，可以用一个可学习查询向量$\bm{u}$对全部位置作注意力汇总，$\alpha_i=\operatorname{softmax}_i(\bm{u}^\top\bm{W}_K\bm{z}_i/\sqrt{d})$，$\bm{g}=\sum_i\alpha_i\bm{W}_V\bm{z}_i$，再接线性分类层，以类别交叉熵训练$\bm{u}$、$\bm{W}_K$、$\bm{W}_V$和分类层。与均值汇总相比，注意力权重让模型自己选择车身、手部等相关位置。对局部深度，论文用最终层网格特征训练线性预测层，即每个位置独立计算$\widehat{y}_i=\bm{w}^\top\bm{z}_i+b$，再把网格上采样回输入分辨率，以有效像素上的深度误差训练$\bm{w},b$；线性层能读出深度，说明这部分几何信息已经以近似线性的形式存在于特征中。对短时交互预期，模块同时读取最后一帧的高分辨率图像和截至该帧的短视频，把视频特征汇总到最后一帧坐标，与图像特征融合，再预测下一交互对象的框、动词及距接触时间。对高层动作预期，则把截至动作开始前一秒的编码特征与预测器给出的未来表示合并，训练动作、动词和名词分类头。训练中，未来帧只用来提供标签，不能作为测试输入；预测器给出的未来表示也不是真实的未来观测。

这种用法有三类常见失败。其一是输入条件不匹配，例如用4帧低分辨率片段读取按16帧训练的权重，局部特征可能失去意义。其二是冻结特征本身没有保留任务需要的信息，线性头再怎么训练也读不出来，此时需要更强的任务模块或微调编码器。其三是不同输出之间的取舍：原论文的消融显示，单独增加可见位置损失可改善密集任务却降低某些分类成绩，多层监督再改变这一取舍。我们可以固定一段拿杯视频，同时检查整段类别是否为拿杯、手附近网格的深度是否合理、只看伸手阶段时是否预期到后续接触。类别正确不能替手的深度通过；预期头答对，也没有生成尚未拍到的未来帧。

三种汇总在同一辆车的识别任务上的差别可以归纳为表\ref{tab:vision-video-aggregation-compare}。从逐帧得分到对象轨迹，改变的是“哪些观测可以合并”由谁决定：前者由调用者给定范围，后者由跟踪器的身份决定，换来了多目标和在线更新，代价是身份错误会直接污染统计。从对象轨迹到密集表示，改变的是合并时保留的信息量：不再压成类别得分，换来了局部状态和预期输出，代价是需要大规模预训练的主干、另行训练的任务头和更高的计算与存储。

\begin{table}[htbp]
  \centering
  \small
  \begin{tabularx}{\linewidth}{@{}lXXX@{}}
    \toprule
    比较项 & 基于逐帧得分 & 基于对象轨迹 & 基于密集时空表示\\
    \midrule
    合并对象 & 类别得分 & 每个身份的得分统计 & 每个时空位置的特征\\
    前提条件 & 调用者给定同一对象或范围 & 跟踪器给出可信身份 & 预训练视频主干\\
    保留信息 & 范围内的类别判断 & 随时间更新的属性 & 局部几何、运动与预期线索\\
    额外训练 & 无，或学习可靠性权重 & 无 & 任务预测头，必要时微调主干\\
    典型失败 & 模糊帧主导或短暂证据被稀释 & 身份交换后错误积累 & 输入条件不匹配，特征缺少任务信息\\
    \bottomrule
  \end{tabularx}
  \caption{三种跨帧汇总方式的对照。比较固定为同一辆车或同一段视频的识别任务；“额外训练”指在已有单帧分类器或预训练主干之外还需要学习的参数。}
  \label{tab:vision-video-aggregation-compare}
\end{table}

\subsection{属性状态的连续判定}
\label{sec:vision-video-continuous-attributes}

汇总回答的是“在这段范围内是什么”，但有些属性本身会变化。车身颜色在短录像中相对稳定，车灯是否点亮却会改变。对变化属性，应输出“哪辆车在何时处于什么状态”，不能用整条轨迹的均值替代状态序列。一个可核验的规则是：仅在连续若干次可靠观测都超过阈值时改变状态，并在长时间遮挡后回到未知。阈值和持续长度需要依据事件标注与误报、漏报代价选择。

例如灯光分数连续为$0.2,0.8,0.3,0.9,0.9$，设阈值为0.7、需要连续两次支持，则第三次出现的低分打断先前的支持，到第五次才确认点亮。这个规则能过滤短跳变，也会增加确认延迟，并可能漏掉真实的短闪烁。若任务关心的是转向灯闪烁频率，就应保留变化而不是强行平滑。时间识别的目标由用途决定，并非输出越稳定越好。

\section{对象状态与空间对应的维护}
\label{sec:vision-video-tracking}

上一节的轨迹汇总假定跟踪器已经告诉我们“这一帧的框延续上一帧的哪辆车”。本节就来回答这个问题。输入是视频帧，再加上首帧给定的目标框、点或掩码，或者逐帧检测器的输出；输出是每个对象在每一帧的状态，包括位置（框、掩码或三维点）、身份编号，以及当前是否可见。难点有四类：对象外观随姿态和光照变化，遮挡使观测暂时消失，相似对象相互靠近时容易交换身份，相机运动又让所有位置一起移动。

维护对象状态是一个按时间反复执行的流程。第一步建立初始状态，明确要延续哪个目标；第二步在新帧中建立与已有状态的对应，即决定当前哪部分观测延续哪个历史对象；第三步用对应结果更新状态，并决定把什么写入历史供以后读取；第四步处理对应失败的情形，即遮挡期间如何保持状态、重新出现时如何接回。每一步都有不同的模型实现，后文依次展开。

\subsection{初始状态的建立}
\label{sec:vision-video-initialization}

初始状态有两种来源。单目标跟踪由用户在首帧给出框、点或掩码，明确接下来延续哪个目标，跟踪器不需要知道它的类别。多目标跟踪则由检测器在每帧发现候选，给新出现且足够可信的候选建立编号；新身份的建立规则与关联规则联系紧密，在\ref{sec:vision-video-association}节随SORT和ByteTrack一起说明。相同的跟踪算法，在人工准确初始化与自动检测初始化下接收到的证据不同，比较时应记录初始化方式。

单目标跟踪的初始状态，就是由首帧给定的框建立的目标模型。此后每一帧都在上一位置附近搜索与这个模型最相符的位置，所以目标模型决定了跟踪器能容忍什么变化、会被什么干扰。三条路线的差别正在于目标模型是什么、怎样与新帧比较。均值漂移跟踪把目标描述成核加权的颜色直方图，用直方图之间的相似度定位，并沿相似度上升的方向迭代，不需要训练。相关滤波把目标描述成一个线性滤波器：它在首帧图像块的所有循环移位上以闭式解训练，使目标中心处的响应最高，并在后续帧上按比例更新。SiamRPN++把目标描述成由深层网络提取的特征，比较方式在大量视频上离线学好，与搜索区域逐通道相关后直接预测候选框，跟踪时不再更新。三者从颜色统计、判别式滤波到学习的表示，目标模型能区分的细节逐步增加，计算与训练数据的需求也依次增加。

\subsubsection{基于核加权直方图的目标模型}
\label{sec:vision-video-meanshift}

设首帧中被跟踪的是一个奔跑的人。四肢不断改变相对位置，若像\ref{sec:vision-features-template-matching}节那样把首帧框内的像素当作模板逐点比较，几帧之后模板就与目标对不齐了；在搜索区域的每个位置都计算一次相似度，代价又随搜索范围的面积增长。Comaniciu等人的做法是放弃像素的空间排列，只保留颜色的统计：目标表示为颜色直方图，靠近中心的像素权重更大；跟踪变成在新帧中寻找直方图与目标最接近的位置，而这个位置可以用均值漂移迭代找到\scite{vision-meanshift-tracking2003}。直方图对形变不敏感；核加权又使相似度随位置平滑变化，于是可以沿上升方向迭代，不必穷举所有位置。

先定义目标模型。把首帧目标区域的坐标归一化，使区域落在单位圆内，区域内像素位置为$\bm{x}_1^*,\ldots,\bm{x}_{n}^*$。函数$b(\bm{x})\in\{1,\ldots,m\}$给出像素$\bm{x}$的颜色所落入的直方图格，$\delta[\cdot]$在自变量为零时取1、否则取0。$k$是单调递减的核轮廓，Epanechnikov轮廓$k(r)=1-r$（$0\leq r\leq1$，其余为0）是原文推荐的选择。目标模型为
\begin{equation}
  \widehat{q}_u=C\sum_{i=1}^{n}k\bigl(\lVert\bm{x}_i^*\rVert_2^2\bigr)\,\delta\bigl[b(\bm{x}_i^*)-u\bigr],
  \qquad u=1,\ldots,m\text{，}
  \label{eq:vision-video-meanshift-target}
\end{equation}
$C$使$\sum_u\widehat{q}_u=1$。在新帧中以$\bm{y}$为中心、带宽为$h$的候选区域内，像素位置为$\bm{x}_1,\ldots,\bm{x}_{n_h}$，候选直方图按同样方式统计：
$\widehat{p}_u(\bm{y})=C_h\sum_{i=1}^{n_h}k(\lVert(\bm{y}-\bm{x}_i)/h\rVert_2^2)\,\delta[b(\bm{x}_i)-u]$。两者的相似度用\term{Bhattacharyya系数}（Bhattacharyya coefficient）衡量：
\begin{equation}
  \rho(\bm{y})=\sum_{u=1}^{m}\sqrt{\widehat{p}_u(\bm{y})\,\widehat{q}_u}\text{。}
  \label{eq:vision-video-bhattacharyya}
\end{equation}
它可以理解为两个单位向量$(\sqrt{\widehat{p}_1},\ldots,\sqrt{\widehat{p}_m})$与$(\sqrt{\widehat{q}_1},\ldots,\sqrt{\widehat{q}_m})$夹角的余弦：两个直方图相同时取1，没有共同的非零格时取0。原文用$d(\bm{y})=\sqrt{1-\rho(\bm{y})}$作为距离，最小化距离等价于最大化$\rho$。

直接对$\rho$求最大需要知道它怎样随$\bm{y}$变化。从上一帧位置$\bm{y}_0$出发，假设新位置附近的候选直方图与$\widehat{p}_u(\bm{y}_0)$相差不大，并且所有$\widehat{p}_u(\bm{y}_0)>0$。对每一项用一阶近似$\sqrt{p}\approx\sqrt{p_0}+(p-p_0)/(2\sqrt{p_0})$，得到
\[
  \begin{aligned}
  \rho(\bm{y})
  &\approx\sum_u\sqrt{\widehat{q}_u}\Bigl[\sqrt{\widehat{p}_u(\bm{y}_0)}+\frac{\widehat{p}_u(\bm{y})-\widehat{p}_u(\bm{y}_0)}{2\sqrt{\widehat{p}_u(\bm{y}_0)}}\Bigr]\\
  &=\frac12\sum_u\sqrt{\widehat{p}_u(\bm{y}_0)\widehat{q}_u}
   +\frac12\sum_u\widehat{p}_u(\bm{y})\sqrt{\frac{\widehat{q}_u}{\widehat{p}_u(\bm{y}_0)}}\text{。}
  \end{aligned}
\]
第二个等号把方括号展开：第一项给出$\sum_u\sqrt{\widehat{p}_u(\bm{y}_0)\widehat{q}_u}$，减去的那部分恰好是它的一半。第一项与$\bm{y}$无关。把$\widehat{p}_u(\bm{y})$的定义代入第二项，并交换对$u$与对$i$的求和，$\delta$只在$u=b(\bm{x}_i)$时非零，于是
\begin{equation}
  \rho(\bm{y})\approx\text{常数}+\frac{C_h}{2}\sum_{i=1}^{n_h}w_i\,k\Bigl(\Bigl\lVert\frac{\bm{y}-\bm{x}_i}{h}\Bigr\rVert_2^2\Bigr),
  \qquad
  w_i=\sqrt{\frac{\widehat{q}_{b(\bm{x}_i)}}{\widehat{p}_{b(\bm{x}_i)}(\bm{y}_0)}}\text{。}
  \label{eq:vision-video-meanshift-weights}
\end{equation}
第二项正是以$w_i$为样本权重的核密度估计。权重的含义很直接：某种颜色在候选中占得比目标模型少，这种颜色的像素就得到大于1的权重，窗口被拉向它们；占得多则权重小于1。核密度估计的局部最大点可以用均值漂移求得：令$g(r)=-k'(r)$，把窗口中心移到
\begin{equation}
  \bm{y}_1=\frac{\sum_i\bm{x}_iw_i\,g\bigl(\lVert(\bm{y}_0-\bm{x}_i)/h\rVert_2^2\bigr)}{\sum_iw_i\,g\bigl(\lVert(\bm{y}_0-\bm{x}_i)/h\rVert_2^2\bigr)}\text{。}
  \label{eq:vision-video-meanshift-update}
\end{equation}
这一步沿核密度估计的梯度方向上升，步长由加权平均自动决定；均值漂移的梯度推导可参考\BookChapterRef{chap:clustering}{卷四的聚类章}。对Epanechnikov轮廓，$g$在窗口内是常数，式\eqref{eq:vision-video-meanshift-update}化为像素位置的加权平均。算法\ref{alg:vision-video-meanshift}给出一帧内的定位过程。原文还有一步保护：若新位置的$\rho$反而下降，就把$\bm{y}_1$移到与$\bm{y}_0$的中点再检查；作者报告实践中很少触发，实际实现省略了这一步，只在迭代结束后计算一次$\rho$。

\begin{algorithm}[基于核加权直方图的均值漂移定位]
  \label{alg:vision-video-meanshift}
  \begin{algorithmic}[1]
    \Require 目标模型$\{\widehat{q}_u\}_{u=1}^{m}$，上一帧位置$\bm{y}_0$，带宽$h$，停止门限$\epsilon$，最大迭代次数$N_{\max}$
    \Ensure 当前帧的目标位置
    \For{$j=1,\ldots,N_{\max}$}
      \State 在$\bm{y}_0$处按核加权统计候选直方图$\{\widehat{p}_u(\bm{y}_0)\}_{u=1}^{m}$
      \State 对窗口内每个像素计算$w_i=\sqrt{\widehat{q}_{b(\bm{x}_i)}/\widehat{p}_{b(\bm{x}_i)}(\bm{y}_0)}$
      \State 按式\eqref{eq:vision-video-meanshift-update}计算$\bm{y}_1$
      \If{$\lVert\bm{y}_1-\bm{y}_0\rVert_2<\epsilon$}
        \State \Return $\bm{y}_1$
      \EndIf
      \State $\bm{y}_0\gets\bm{y}_1$
    \EndFor
    \State \Return $\bm{y}_0$
  \end{algorithmic}
\end{algorithm}

用一个一维教学算例走一遍。一行五个像素位于$0,1,2,3,4$，颜色依次为灰、灰、灰、红、红；目标模型中红占0.6、灰占0.4。上一帧位置$y_0=2$，带宽$h=2.5$，Epanechnikov轮廓给五个像素的核权重为$1-(\text{距离}/h)^2$，即$0.36,0.84,1,0.84,0.36$。候选直方图中红色权重为$0.84+0.36=1.2$，灰色为$2.2$，归一化后为$(0.353,0.647)$，此时$\rho\approx\sqrt{0.353\times0.6}+\sqrt{0.647\times0.4}\approx0.969$。红色在候选中偏少，权重为$\sqrt{0.6/0.353}\approx1.304$；灰色偏多，权重为$\sqrt{0.4/0.647}\approx0.786$。$g$为常数，新位置为
\[
  y_1=\frac{(0+1+2)\times0.786+(3+4)\times1.304}{3\times0.786+2\times1.304}
  \approx\frac{11.486}{4.966}\approx2.31\text{。}
\]
窗口向红色像素移动了约0.31个像素，图\ref{fig:vision-meanshift-window}标出了各像素的核权重、颜色权重和窗口中心的移动。在$y_1$处重新统计，红色比例升到0.442，$\rho$升到约0.987；再迭代下去，位移会逐步减小，直到小于门限。

\begin{figure}[htbp]
  \centering
  % 均值漂移一维教学算例；宽约90mm。五个像素位于0至4，颜色灰灰灰红红；目标模型红0.6、灰0.4；y0=2，h=2.5，Epanechnikov轮廓。
\begin{tikzpicture}[x=1.45cm,y=1cm,font=\figurefont\small,
  arr/.style={-{Stealth[length=2mm]},BookInk,line width=0.9pt}]
  % 像素
  \foreach \x/\col in {0/BookMuted!30,1/BookMuted!30,2/BookMuted!30,3/FigureRed!70,4/FigureRed!70}{
    \fill[\col] (\x-0.4,1.6) rectangle (\x+0.4,2.4);
    \draw[BookInk,line width=0.5pt] (\x-0.4,1.6) rectangle (\x+0.4,2.4);
    \node[below] at (\x,1.55) {$\x$};
  }
  \node[anchor=east] at (-0.6,2.0) {像素};
  % 核权重
  \node[anchor=east] at (-0.6,0.75) {核权重};
  \foreach \x/\k in {0/0.36,1/0.84,2/1.00,3/0.84,4/0.36}{
    \node at (\x,0.75) {$\k$};
  }
  % 像素权重
  \node[anchor=east] at (-0.6,0.15) {$w_i$};
  \foreach \x/\w in {0/0.786,1/0.786,2/0.786,3/1.304,4/1.304}{
    \node at (\x,0.15) {$\w$};
  }
  % 窗口中心与移动
  \draw[BookTeal,line width=1.1pt] (2,2.55) -- (2,3.05);
  \node[above] at (2,3.05) {$y_0=2$};
  \draw[BookBlue,line width=1.1pt] (2.31,2.55) -- (2.31,3.05);
  \draw[arr] (2.02,2.8) -- (2.29,2.8);
  \node[anchor=west] at (2.4,3.2) {$y_1\approx2.31$};
  % 带宽范围
  \draw[BookTeal,line width=0.7pt,{Bar[width=2mm]}-{Bar[width=2mm]}] (-0.5,-0.45) -- (4.5,-0.45);
  \node[below] at (2,-0.5) {带宽$h=2.5$覆盖的窗口};
\end{tikzpicture}

  \caption{均值漂移一次迭代的教学算例。五个像素中灰色三个、红色两个，目标模型中红占0.6、灰占0.4。核权重决定候选直方图，颜色权重$w_i$由目标与候选的比例给出；红色在候选中偏少，权重大于1，窗口中心因此从2移到约2.31。数值为教学构造。}
  \label{fig:vision-meanshift-window}
\end{figure}

这一做法的有效条件也由推导给出。一阶近似要求相邻帧之间直方图变化不大，原文的说法是：只要新帧中目标中心仍落在上一帧目标模型覆盖的区域内，Bhattacharyya系数的局部最大就是可靠的定位依据。原文实验在RGB空间用$16\times16\times16$个颜色格，Football序列中每帧平均迭代4.19次，最大迭代次数通常设为20。尺度变化由带宽处理：分别以$h$、$1.1h$和$0.9h$运行定位，保留$\rho$最大者，再按$h_{\mathrm{new}}=0.1h_{\mathrm{opt}}+0.9h_{\mathrm{prev}}$平滑，避免尺度跳变。失效情形同样清楚：目标移出上一窗口时，梯度上升找不到它；颜色相近的干扰者或同色背景会让直方图失去区分力；直方图不记录颜色的空间排列，红衣在上、蓝裤在下与上下颠倒得到同一个模型。均值漂移跟踪留下的思想是“在上一位置附近、用平滑的相似度做局部搜索”，后面的相关滤波和孪生网络都保留了这一结构，只是换掉了目标模型和比较方式。

\subsubsection{基于相关滤波的目标模型}
\label{sec:vision-video-correlation-filter}

颜色直方图丢掉了空间排列，\ref{sec:vision-features-template-matching}节的模板匹配保留了排列，却只从一个样本出发，没有学过怎样与背景区分。一个更好的模板应当满足判别要求：与目标中心对齐时响应高，偏离目标、落到背景上时响应低。把首帧图像块的每一个平移版本都当作训练样本，目标中心对应高响应、其他位置对应低响应，就可以学出这样一个线性滤波器。难点在计算：一个$64\times64$的图像块有4096个平移版本，在空间域逐个求解代价很高。\term{相关滤波}（correlation filter）的做法是把图像块视为周期延拓，于是所有平移都是循环移位，训练和检测都可以在频域逐元素完成。

需要用到的频域结论是相关定理。卷一给出了一维离散循环卷积定理：长度为$N$的两个序列，循环卷积的DFT等于各自DFT的逐元素乘积，可参考\BookChapterRef{chap:signal-analysis}{卷一的信号章}。互相关可以由它推出。按卷一的方向约定，循环互相关为$g[n]=\sum_{r}f[(n+r)\bmod N]\,\overline{h[r]}$，即把模板$h$平移到位置$n$后与$f$作内积。对它作DFT，并令$a=(n+r)\bmod N$：
\[
  \begin{aligned}
  \widehat{g}[k]&=\sum_n\sum_r f[(n+r)\bmod N]\,\overline{h[r]}\,\omega_N^{kn}\\
  &=\Bigl(\sum_a f[a]\,\omega_N^{ka}\Bigr)\Bigl(\sum_r\overline{h[r]}\,\omega_N^{-kr}\Bigr)
  =\widehat{f}[k]\,\overline{\widehat{h}[k]}\text{。}
  \end{aligned}
\]
第二个等号用到$\omega_N^{kn}=\omega_N^{ka}\omega_N^{-kr}$以及$\omega_N^{kN}=1$；第三个等号用到$|\omega_N|=1$，因而$\overline{\omega_N^{kr}}=\omega_N^{-kr}$。所以$\widehat{g}=\widehat{f}\odot\widehat{h}^*$，上标$*$表示逐元素取共轭。二维情形按行、列各做一次一维变换即可，二维DFT见\ref{sec:vision-transforms-dft2d}节。在$P$个像素的图像块上，一次相关只需两次正变换和一次逆变换，代价为$O(P\log P)$，而得到的是所有循环位移上的响应。

MOSSE（minimum output sum of squared error）滤波器直接在频域训练\scite{vision-mosse2010}。给定训练图像块$f_i$及期望输出$g_i$，$g_i$是以目标中心为峰、标准差为2像素的二维高斯。要求滤波器的输出尽量接近期望输出：
\begin{equation}
  \min_{\widehat{h}^*}\ \sum_i\bigl|\widehat{f}_i\odot\widehat{h}^*-\widehat{g}_i\bigr|^2\text{。}
  \label{eq:vision-video-mosse-objective}
\end{equation}
由于频域运算逐元素进行，每个频率可以单独求解。固定一个频率，记$z=\widehat{h}^*$在该频率的值，$A=\sum_i|\widehat{f}_i|^2$，$B=\sum_i\widehat{g}_i\overline{\widehat{f}_i}$。展开$|\widehat{f}_iz-\widehat{g}_i|^2=|\widehat{f}_i|^2|z|^2-2\operatorname{Re}(z\,\widehat{f}_i\overline{\widehat{g}_i})+|\widehat{g}_i|^2$并对$i$求和，得到$A|z|^2-2\operatorname{Re}(z\overline{B})+\text{常数}$。它等于$A|z-B/A|^2$加上与$z$无关的项，因此最小点为$z=B/A$，即
\begin{equation}
  \widehat{h}^*=\frac{\sum_i\widehat{g}_i\odot\widehat{f}_i^*}{\sum_i\widehat{f}_i\odot\widehat{f}_i^*}\text{，}
  \label{eq:vision-video-mosse}
\end{equation}
分数线表示逐元素相除。分子是输入与期望输出的互相关，分母是输入的能量谱。某个频率上所有训练块的能量都很小时，分母接近零，滤波器在该频率被放大；原文在分母上加一个小的正则量$\epsilon$，并报告初始化时$\epsilon$取0.1左右效果较好。

MOSSE的完整流程有四个环节。预处理先对像素取对数以缓解低对比度，再归一化为均值0、范数1，最后乘余弦窗，使边缘像素逐渐降为零，减轻周期延拓在边界处造成的伪影。初始化时对首帧跟踪窗口作8次随机仿射扰动，得到8个训练块。在线更新把分子与分母分别做指数滑动平均：$A_t=\eta\,\widehat{g}_t\odot\widehat{f}_t^*+(1-\eta)A_{t-1}$，$B_t=\eta\,\widehat{f}_t\odot\widehat{f}_t^*+(1-\eta)B_{t-1}$，滤波器取$A_t/B_t$，原文取$\eta=0.125$。失败检测用峰值旁瓣比：响应图的最大值减去峰值周围$11\times11$窗口以外区域的均值，再除以该区域的标准差。原文报告正常跟踪时这一比值在20到60之间，降到7左右通常意味着遮挡或跟踪失败，此时可以暂停更新。整个跟踪器在原文实验中达到每秒669帧。

MOSSE为什么有效？Henriques等人给出的解释是：它等价于在所有循环移位样本上做岭回归，而循环结构使岭回归可以在频域闭式求解；由此还可以把它推广到核方法与多通道特征，这就是\term{核相关滤波}（kernelized correlation filter, KCF）\scite{vision-kcf2015}。为简洁，下面用一维实信号说明，二维与多通道情形逐元素推广。设基样本$\bm{x}\in\mathbb{R}^n$，$P$是循环移位算子$(P\bm{x})[j]=x[j-1]$（下标按模$n$计），训练样本是$P^u\bm{x}$，$u=0,\ldots,n-1$，对应的回归目标$y_u$是零位移处为1、随位移增大而平滑衰减的高斯。把样本按行排成数据矩阵$\bm{X}$，其第$u$行为$(P^u\bm{x})^\top$，这样的矩阵称为\term{循环矩阵}（circulant matrix）：每一行都是上一行循环右移一位，见图\ref{fig:vision-circulant-shifts}。岭回归
\[
  \min_{\bm{w}}\ \lVert\bm{X}\bm{w}-\bm{y}\rVert_2^2+\lambda\lVert\bm{w}\rVert_2^2
\]
的解满足正规方程$(\bm{X}^\top\bm{X}+\lambda\bm{I})\bm{w}=\bm{X}^\top\bm{y}$，岭回归的一般性质可参考\BookChapterRef{chap:regression}{卷四的回归章}。逐项看这些乘积：$(\bm{X}\bm{w})[u]=\sum_jx[j-u]w[j]$是$\bm{w}$与$\bm{x}$的循环互相关，由上面的相关定理，其DFT为$\widehat{w}\odot\widehat{x}^*$；$(\bm{X}^\top\bm{v})[j]=\sum_ux[j-u]v[u]$是$\bm{x}$与$\bm{v}$的循环卷积，其DFT为$\widehat{x}\odot\widehat{v}$。于是$\bm{X}^\top\bm{X}\bm{w}$的DFT为$\widehat{x}\odot\widehat{x}^*\odot\widehat{w}$，正规方程在每个频率上都变成一个标量方程，解为
\begin{equation}
  \widehat{w}=\frac{\widehat{x}\odot\widehat{y}}{\widehat{x}\odot\widehat{x}^*+\lambda}\text{。}
  \label{eq:vision-video-dcf}
\end{equation}
这一步的本质是：DFT的基向量（复指数）都是循环移位$P$的特征向量，任何循环矩阵都是$P$的多项式，因而被DFT同时对角化，求逆变成逐元素相除。原文用DFT矩阵的对角化记号推导，结果写作$\widehat{x}^*\odot\widehat{y}$的形式；共轭落在哪一侧取决于移位方向和变换矩阵的记号约定，实现时应按所用FFT的约定重新核对。式\eqref{eq:vision-video-dcf}与单个训练样本、$\epsilon=\lambda$时的式\eqref{eq:vision-video-mosse}形式相同，KCF原文也指出线性核、单通道时的对偶相关滤波（dual correlation filter, DCF）在极限情形下等价于MOSSE。

核方法把线性岭回归换成在特征空间中求解，解写成对偶系数$\bm{\alpha}=(\bm{K}+\lambda\bm{I})^{-1}\bm{y}$，其中$\bm{K}$的元素是样本两两之间的核函数值。这正是\term{核技巧}（kernel trick）：只通过核函数计算样本间的内积，而不显式构造特征，核岭回归在上面引用的同一章中有完整推导。一般的$\bm{K}$没有可利用的结构，求逆代价为$O(n^3)$。KCF证明：只要核函数满足$\kappa(\bm{x},\bm{x}')=\kappa(M\bm{x},M\bm{x}')$对任意置换矩阵$M$成立（高斯核、多项式核等都满足），循环移位样本的核矩阵仍是循环矩阵。于是
\[
  \widehat{\alpha}=\frac{\widehat{y}}{\widehat{k}^{\bm{xx}}+\lambda},\qquad
  \widehat{f}(\bm{z})=\widehat{k}^{\bm{xz}}\odot\widehat{\alpha}\text{，}
\]
其中$\bm{k}^{\bm{xx}'}$称为核相关，第$u$个元素为$\kappa(\bm{x}',P^u\bm{x})$；$\bm{f}(\bm{z})$是新图像块$\bm{z}$所有循环位移上的响应，取最大值的位移就是目标的移动量。高斯核的核相关也能用FFT一次算出：$\bm{k}^{\bm{xx}'}=\exp\bigl(-\sigma^{-2}(\lVert\bm{x}\rVert_2^2+\lVert\bm{x}'\rVert_2^2-2\mathcal{F}^{-1}(\sum_c\widehat{x}_c^*\odot\widehat{x}'_c))\bigr)$，$c$遍历特征通道，原文实现中再把括号内的量除以元素个数。训练与检测的代价都是$O(n\log n)$，与线性滤波相同。

原文的跟踪流程有意保持简单，见算法\ref{alg:vision-video-kcf}。跟踪区域取目标大小的2.5倍，以提供背景作为负样本；特征用单元大小为4像素、每单元31个通道的HOG（见\ref{sec:vision-features-hog}节），乘余弦窗后进入计算。全部视频使用同一组参数：HOG特征下高斯核带宽$\sigma=0.5$、更新率0.02，回归目标的空间带宽为$\sqrt{mn}/10$（$m,n$是以HOG单元计的目标高和宽），$\lambda=10^{-4}$。高斯回归目标的峰放在左上角，即零位移处，并循环延伸到其他三个角，这与循环相关的零位移位置一致。原文的流程不估计尺度变化，也没有失败检测。在50个序列的基准上，HOG特征的KCF在20像素门限处的平均精度为73.2\%、速度为每秒172帧，线性核的DCF为72.8\%、每秒292帧；改用原始像素时两者分别降到56.0\%和45.1\%。核函数带来的提升不大，特征从像素换成HOG的提升要大得多。

\begin{algorithm}[基于核相关滤波的跟踪]
  \label{alg:vision-video-kcf}
  \begin{algorithmic}[1]
    \Require 首帧图像与目标框，核函数$\kappa$，正则系数$\lambda$，更新率$\eta$
    \Ensure 每帧的目标位置
    \State 在首帧以目标为中心裁取放大的图像块，提取特征并乘余弦窗，得$\bm{x}$
    \State 构造峰值位于零位移的高斯回归目标$\bm{y}$
    \State $\widehat{\alpha}\gets\widehat{y}/(\widehat{k}^{\bm{xx}}+\lambda)$
    \For{每个新帧}
      \State 在上一位置裁取同样大小的块，提取特征并加窗，得$\bm{z}$
      \State $\bm{r}\gets\mathcal{F}^{-1}(\widehat{k}^{\bm{xz}}\odot\widehat{\alpha})$，按最大响应的循环位移修正目标位置
      \State 在新位置裁取训练块$\bm{x}'$，按第3行计算$\widehat{\alpha}'$
      \State $\widehat{\alpha}\gets(1-\eta)\widehat{\alpha}+\eta\widehat{\alpha}'$，$\bm{x}\gets(1-\eta)\bm{x}+\eta\bm{x}'$
    \EndFor
  \end{algorithmic}
\end{algorithm}

用线性核与$n=4$的一维信号手算一遍。基样本$\bm{x}=(2,1,0,1)$，回归目标$\bm{y}=(1,0.25,0,0.25)$，$\lambda=0.1$。取$\omega_4=e^{-2\pi\mathrm{i}/4}$，DFT为$\widehat{x}=(4,2,0,2)$、$\widehat{y}=(1.5,1,0.5,1)$，二者都是实数。按式\eqref{eq:vision-video-dcf}，$\widehat{w}=(4\times1.5/16.1,\ 2/4.1,\ 0,\ 2/4.1)\approx(0.373,0.488,0,0.488)$，逆变换得$\bm{w}\approx(0.337,0.093,-0.151,0.093)$。检查训练样本：四个循环移位与$\bm{w}$的内积为$(0.860,0.373,-0.115,0.373)$，形状与$\bm{y}$一致、幅度被正则收缩。第三个频率上$\widehat{x}=0$，该频率不含任何信息，若没有$\lambda$就会出现零除零；正则项让这一频率的滤波器取0。下一帧目标右移一格，图像块变为$\bm{z}=(1,2,1,0)$。四个循环位移$P^u\bm{z}$与$\bm{w}$的内积为$(0.373,-0.115,0.373,0.860)$，最大值出现在$u=3$，即$P^3\bm{z}=P^{-1}\bm{z}=\bm{x}$：把$\bm{z}$左移一格才与训练样本对齐，所以目标右移了一格。用FFT计算时，四个响应由一次逐元素乘积和一次逆变换同时得到。

\begin{figure}[htbp]
  \centering
  % 相关滤波的循环移位样本与回归目标；宽约80mm。基样本x=(2,1,0,1)，第u行为P^u x，(Px)[j]=x[j-1]；回归目标y=(1,0.25,0,0.25)。
\begin{tikzpicture}[x=0.85cm,y=0.75cm,font=\figurefont\small]
  \node[anchor=south] at (1.5,0.75) {训练样本$P^u\bm{x}$};
  \node[anchor=south] at (5.6,0.75) {目标$y_u$};
  \foreach \u/\a/\b/\c/\d/\y in {0/2/1/0/1/1,1/1/2/1/0/0.25,2/0/1/2/1/0,3/1/0/1/2/0.25}{
    \node[anchor=east] at (-0.6,-\u) {$u=\u$};
    \foreach \j/\v in {0/\a,1/\b,2/\c,3/\d}{
      \pgfmathsetmacro{\shade}{\v*35}
      \fill[BookBlue!\shade!white] (\j-0.5,-\u-0.5) rectangle (\j+0.5,-\u+0.5);
      \draw[BookRule,line width=0.5pt] (\j-0.5,-\u-0.5) rectangle (\j+0.5,-\u+0.5);
      \node at (\j,-\u) {$\v$};
    }
    \pgfmathsetmacro{\yshade}{\y*45}
    \fill[BookTeal!\yshade!white] (5.1,-\u-0.5) rectangle (6.1,-\u+0.5);
    \draw[BookRule,line width=0.5pt] (5.1,-\u-0.5) rectangle (6.1,-\u+0.5);
    \node at (5.6,-\u) {$\y$};
  }
  \draw[-{Stealth[length=2mm]},BookMuted,line width=0.7pt] (3.7,-1.5) -- (4.9,-1.5);
\end{tikzpicture}

  \caption{相关滤波的训练样本与回归目标。基样本$\bm{x}=(2,1,0,1)$的四个循环移位按行排成循环矩阵，颜色深浅表示数值；右列是每个移位对应的回归目标，零位移处为1，向两侧循环衰减。数值为正文手算算例的教学构造。}
  \label{fig:vision-circulant-shifts}
\end{figure}

相关滤波的局限来自它的两个假设。循环移位只是平移的近似，靠近边界的移位样本含有绕回的错误内容，余弦窗和扩大的搜索区只能减轻这一点；目标移出搜索区或发生大幅尺度变化时，响应图上找不到正确的峰。按固定比例的在线更新使模型能追随缓慢的外观变化，也使遮挡期间的背景被逐渐写进模型，MOSSE的峰值旁瓣比正是为此设置的停止条件。相关滤波留下的核心运算是“模板特征与搜索区域特征的相关”，下一小节的孪生网络保留了这一运算，把特征换成离线学到的深层特征，并放弃了在线更新。

\subsubsection{基于孪生网络的目标模型}
\label{sec:vision-video-siamrpn}

相关滤波的特征是人工设计的灰度或HOG，滤波器也只从当前视频学习。能否从大量视频中离线学到“模板与候选是否相像”的比较方式，跟踪时只做前向计算？\term{孪生网络}（Siamese network）用同一组参数分别编码两个输入，再比较二者的特征，正适合这种判断。SiamFC和SiamRPN等早期孪生跟踪器已经采用这一结构，但骨干都是AlexNet这类浅层网络；直接换成ResNet等深层网络，精度并没有相应提升。SiamRPN++的作者把原因归结为平移不变性被破坏：深层网络的填充使特征带有位置偏好，而训练时目标总在搜索区域中心，网络于是学会了“中心更可能是目标”的捷径，而不是真正比较外观\scite{vision-siamrpn2019}。

SiamRPN++的主要想法有三点。第一，训练时让目标在搜索区域内随机偏移，称为空间感知采样，打破“目标总在中心”的统计偏好，使深层骨干可以使用；原文的消融显示，偏移范围取$\pm64$像素较为合适。第二，把模板与搜索特征的比较改为逐通道相关，减少参数并让两支特征的尺度更平衡。第三，在骨干的多层特征上分别做比较与预测，再加权融合，使浅层的定位细节和深层的语义区分共同决定结果。

模型架构见表\ref{tab:vision-video-siamrpn-arch}。首帧以目标框为中心裁出$127\times127\times3$的模板图像，后续帧以上一帧位置为中心裁出$255\times255\times3$的搜索图像$\bm{S}_t$，两者经过共享的ResNet-50骨干（标准ResNet-50的结构见\ref{sec:vision-ivm-resnet}节）。原ResNet最后两个阶段的累计步长为16和32像素，对密集预测过粗；原方法把第四、五阶段的空间步长改为1，并用空洞卷积补回感受野，使这两个阶段的有效步长都保持为8。第三、四、五阶段的输出各接一个$1\times1$卷积，把通道降到256。由于保留了所有填充，模板特征的空间尺寸为$15\times15$，原方法只裁取中心$7\times7$作为模板特征，每个特征单元仍能覆盖整个目标。设某层模板特征为$\bm{Z}\in\mathbb{R}^{h_z\times w_z\times d}$、搜索特征为$\bm{S}\in\mathbb{R}^{h_s\times w_s\times d}$，逐通道相关对每个通道$k$计算
\begin{equation}
  r_{u,v,k}=\sum_{i,j}z_{i,j,k}\,s_{u+i,v+j,k}\text{，}
  \label{eq:vision-video-depthwise-correlation}
\end{equation}
其中小写元素取出相应位置及通道，$(u,v)$枚举模板在搜索网格中的放置位置，输出$\bm{R}\in\mathbb{R}^{(h_s-h_z+1)\times(w_s-w_z+1)\times d}$。与把所有通道一起求内积、只得到单通道响应的相关不同，逐通道相关保留了$d$个通道的比较结果，后续卷积再组合这些通道。它和相关滤波中的运算相同，区别在于参与相关的是学到的特征，而且模板不随视频更新。

每个SiamRPN模块内部的计算顺序是：模板支与搜索支各经一层不共享的卷积（卷积加批归一化）调整特征，因为锚框分类与框回归对两支的要求并不对称；两支逐通道相关；再接一个卷积—批归一化—ReLU块融合通道；最后由一层卷积分别输出分类与回归结果。分类头为每个位置的$k$个锚框输出前景与背景两个得分，共$2k$个通道；回归头为每个锚框输出四个偏移，共$4k$个通道。原文取$k=5$，输出的空间尺寸为$25\times25$。三层的输出空间尺寸相同，直接加权相加：$\bm{S}_{\mathrm{all}}=\sum_{l=3}^{5}\alpha_l\bm{S}_l$，$\bm{B}_{\mathrm{all}}=\sum_{l=3}^{5}\beta_l\bm{B}_l$，权重$\alpha_l,\beta_l$随网络一起学习。

\begin{table*}[tbp]
  \centering
  \small
  \begin{tabularx}{\linewidth}{@{}lXXlll@{}}
    \toprule
    阶段或层组 & 层类型 & 卷积核与步长 & 输出通道 & 重复次数 & 输出尺寸\\
    \midrule
    conv1 & 卷积 & $7\times7$，步长2 & 64 & 1 & 输入的1/2\\
    pool & 最大池化 & $3\times3$，步长2 & 64 & 1 & 输入的1/4\\
    conv2\_x & 瓶颈残差块 & $1\times1$、$3\times3$、$1\times1$ & 256 & 3 & 输入的1/4\\
    conv3\_x & 瓶颈残差块 & 同上，首块步长2 & 512 & 4 & 输入的1/8\\
    conv4\_x & 瓶颈残差块 & 同上，步长改为1，用空洞卷积 & 1024 & 6 & 输入的1/8\\
    conv5\_x & 瓶颈残差块 & 同上，步长改为1，用空洞卷积 & 2048 & 3 & 输入的1/8\\
    调整层 & 卷积 & $1\times1$，接在conv3至conv5后 & 256 & 各1 & 模板15，裁为7\\
    SiamRPN模块 & 逐通道相关头 & 各卷积核大小论文未给出 & $2k$与$4k$ & 3 & $25\times25$\\
    加权融合 & 学习权重的加权和 & — & $2k$与$4k$ & 1 & $25\times25$\\
    \bottomrule
  \end{tabularx}
  \caption{SiamRPN++的逐层结构，输入为$127\times127\times3$的模板与$255\times255\times3$的搜索图像，$k=5$。前六行沿用标准ResNet-50，只改动第四、五阶段的步长；空洞率与搜索支中间特征的空间尺寸论文未给出，输出尺寸以相对输入的步长表示。}
  \label{tab:vision-video-siamrpn-arch}
\end{table*}

训练使用同一目标在同一视频不同帧中的模板与搜索图像对，原工作的图像对来自ImageNet VID、YouTube-BB、COCO和ImageNet DET等带框数据，空间感知采样在这一步加入随机偏移。锚框与真实框的交并比决定它是前景、背景还是忽略。设搜索裁剪中的锚框中心和尺寸为$(x_a,y_a,w_a,h_a)$，真实框为$(x^*,y^*,w^*,h^*)$，回归目标为
\[
  \delta_x^*=\frac{x^*-x_a}{w_a},\quad
  \delta_y^*=\frac{y^*-y_a}{h_a},\quad
  \delta_w^*=\log\frac{w^*}{w_a},\quad
  \delta_h^*=\log\frac{h^*}{h_a}\text{。}
\]
中心差除以锚框尺寸，使不同尺度能使用同一输出规则；尺寸比取对数，使放大和缩小在相反方向上对称表达。总目标是分类交叉熵与正锚框上的平滑$L_1$回归误差之和：小残差使用二次项以保持梯度平稳，大残差使用线性项以减轻离群样本的影响。原工作从ImageNet分类预训练的骨干出发，三个$1\times1$调整层随机初始化；用8块GPU同步随机梯度下降，每批共128个图像对。前5轮只以0.001的学习率预热训练预测分支，后15轮端到端训练整个网络，学习率从0.005指数衰减到0.0005，骨干的学习率取预测分支的十分之一，权重衰减0.0005、动量0.9。训练完成后参数固定，跟踪时不再为每段视频重新训练；原文报告在Titan Xp上的速度为每秒35帧。

推理时，首帧模板特征只计算一次并保持不变，这就是SiamRPN++的初始状态；此后每帧只更新位置。预测偏移按$x=x_a+w_a\delta_x$、$y=y_a+h_a\delta_y$、$w=w_a\exp\delta_w$、$h=h_a\exp\delta_h$反解为裁剪框，再乘以搜索区域缩放的逆比例并加回裁剪原点，得到原帧位置。例如教学锚框$(32,32,20,10)$预测$(0.2,0,\log1.2,0)$，得到裁剪框$(36,32,24,10)$；若裁剪原点为$(20,10)$且未缩放，原帧中心就是$(56,42)$。实际选框之前，还要用余弦窗口惩罚离上一位置过远的候选，用尺度惩罚抑制尺寸和长宽比的突变，再按既定平滑规则更新尺寸，并用更新后的位置裁取下一帧搜索区。余弦窗口在这里的作用与均值漂移的局部搜索相同，表达的是“目标在相邻帧间不会跳得太远”。

这些约束表达了“目标连续运动”的偏好，也带来相应的失败情形。假设汽车A在首帧中心$x=40$像素，搜索范围覆盖下一帧的$x=20$至80像素。候选A在$x=48$处得分0.80，另一辆相似汽车B在$x=70$处得分0.85。余弦窗口会压低离上一位置较远的B，因此通常能保住A；但当A突然加速、完全遮挡或离开搜索范围时，同样的约束会把结果拉向错误位置。固定首帧模板避免了把错误目标写进模板，却难以适应外观的大幅变化；在线改写模板属于额外设计，跨镜头的身份恢复也超出了局部搜索的输入条件。

三种目标模型在同一单目标跟踪任务上的差别见表\ref{tab:vision-video-sot-compare}。从核加权直方图到相关滤波，改变的是目标模型的判别性：直方图只描述目标自身的颜色，相关滤波把所有背景位移当作负样本学习，因而保留了空间排列，代价是循环边界效应和在线更新带来的漂移。从相关滤波到SiamRPN++，改变的是特征与比较方式的来源：从人工特征加当前视频的在线训练，换成大规模数据上离线学到的深层特征和锚框预测，换来了对外观变化的更强区分和直接的尺度估计，代价是需要大量带框视频训练、需要GPU，并且模板固定后无法适应目标外观的长期变化。表中速度来自各自论文、硬件不同，只能说明量级。

\begin{table*}[tbp]
  \centering
  \small
  \begin{tabularx}{\linewidth}{@{}lXXX@{}}
    \toprule
    比较项 & 核加权直方图（均值漂移） & 相关滤波（MOSSE、KCF） & 孪生网络（SiamRPN++）\\
    \midrule
    目标模型 & 首帧区域的核加权颜色直方图 & 在所有循环移位上闭式求得的线性或核滤波器 & 离线学到的深层特征模板\\
    定位方式 & 均值漂移沿相似度上升，每帧数次迭代 & 一次FFT得到所有循环位移的响应 & 逐通道相关后预测锚框得分与偏移\\
    尺度处理 & 三种带宽中取最优再平滑 & 原文流程不估计 & 锚框回归给出宽高\\
    在线更新 & 假定可周期性更新目标模型 & 每帧按固定比例插值 & 不更新，模板固定为首帧\\
    训练数据 & 无 & 无，当前视频即训练数据 & 大规模带框视频与图像\\
    论文报告的速度 & 每帧平均4.19次迭代 & MOSSE每秒669帧，KCF每秒172帧 & 每秒35帧（GPU）\\
    典型失败 & 同色干扰或同色背景，目标移出窗口 & 遮挡时吸收背景，循环边界，尺度变化 & 相似干扰者，外观大变，离开搜索区\\
    \bottomrule
  \end{tabularx}
  \caption{单目标跟踪三种目标模型的对照。比较固定为同一段视频、首帧给定框的单目标跟踪；速度取自各论文，硬件与输入尺寸不同，不能直接比较。}
  \label{tab:vision-video-sot-compare}
\end{table*}

\subsection{跨帧对应的建立}
\label{sec:vision-video-correspondence}

有了初始状态，每到一帧新画面，就要决定当前观测的哪一部分延续历史中的哪一部分。这一步的核心选择是对应的粒度。粒度最细的是像素：只问第一帧的每个像素在第二帧移到了哪里，不关心它属于哪个对象。粒度最粗的是对象：只问当前的哪个检测框延续哪条轨迹，不关心框内每个像素怎样移动。第三种是表面点：既跟踪物体表面的单个点，又要求它在共同的三维坐标中随时间移动，从而把相机运动与物体运动分开。

三种粒度回答的是不同的问题，也不能相互替代。某个像素的位移估计很精确，并不自动把遮挡前后的两辆相似车区分开；两条车辆轨迹的身份正确，也说不清车门是否在转动；二维位移和身份都正确，仍不能告诉我们手臂向右是因为手臂在动还是相机在转。下面依次以RAFT讲像素级的对应，以卡尔曼预测加一对一分配及其改进ByteTrack讲对象级的对应，以D4RT讲表面点级的对应。

\subsubsection{像素级的对应}
\label{sec:vision-video-raft}

像素级的对应就是估计光流，即两幅图像间像素的二维位移场，定义见\ref{sec:vision-videofound-motion}节。对第一帧的位置$\bm{p}=(x,y)^\top$，位移$\bm{u}(\bm{p})\in\mathbb{R}^{2}$给出第二帧的对应位置$\bm{p}+\bm{u}(\bm{p})$。横轴向右、纵轴向下，位移以当前网格的像素为单位。光流描述图像中的运动，相机移动、物体运动及透视共同影响它。

光流约束方程、孔径问题与经典解法已在\ref{sec:vision-videofound-motion}节讲过，这里只回顾RAFT要改动的几个结论。亮度守恒假设同一表面点在相邻帧中亮度不变，即$I(\bm{p}+\bm{u},t+1)=I(\bm{p},t)$。位移很小时，对左边作一阶泰勒展开，得到光流约束方程$I_xu+I_yv+I_t=0$，其中$I_x,I_y,I_t$是图像对横坐标、纵坐标和时间的偏导，$(u,v)$是位移的两个分量。每个像素只有一个方程、两个未知数，只能确定沿梯度方向的位移分量，这是孔径问题的代数形式。\ref{sec:vision-videofound-horn-schunck}节的Horn--Schunck方法在整幅图上最小化约束方程残差的平方与位移场梯度的平方之和，用平滑项把信息从纹理丰富处传到均匀区域。\ref{sec:vision-videofound-lucas-kanade}节的Lucas--Kanade方法假设小窗口内位移相同，把窗口内各像素的方程联立成最小二乘问题，系数矩阵是结构张量，即梯度外积之和$\sum\bigl(\begin{smallmatrix}I_x^2&I_xI_y\\I_xI_y&I_y^2\end{smallmatrix}\bigr)$；只有它的两个特征值都不太小时解才稳定，这与\ref{sec:vision-features-harris}节Harris角点的判据相同。位移超出线性化的范围时，两种方法都借助图像金字塔由粗到细地估计。

这些结论指出了经典光流的三个薄弱点。亮度守恒在光照变化和反光处不成立；线性化只在小位移内有效，金字塔的粗层又会把快速移动的小物体平滑掉；平滑项或窗口内位移相同的假设由人设定，在运动边界处会把两侧的位移混在一起。RAFT逐一替换了这三处：用学到的特征之间的内积代替亮度差，用全部像素对的相关代替线性化，用学到的循环更新代替手工平滑项。它保留的是经典方法的迭代结构：Horn--Schunck的迭代求解和Lucas--Kanade的迭代精化，都在反复“按当前估计读取对应、再修正估计”，RAFT的更新模块做的是同一件事。

在RAFT之前，学习式光流网络大多沿用\ref{sec:vision-videofound-lucas-kanade}节金字塔Lucas--Kanade的由粗到细思路：先在低分辨率上估计大位移，再逐级上采样并修正。这种做法有一个结构性的弱点：快速移动的小物体在最粗一级可能已经消失，粗级的错误传到细级后很难纠正；而级数固定，也就限定了修正次数。RAFT的主要想法是只在一个分辨率上维护光流，先一次性算出所有像素对之间的相关，再用一个共享参数的循环更新模块反复读取局部相关并修正位移\scite{vision-raft2020}。所有像素对的相关保留了大位移的候选，多次迭代则让每次修正都能利用上一轮的结果。

模型架构分三部分。第一部分是特征提取：两帧$H\times W\times3$的图像经共享的特征编码器得到约$1/8$分辨率的特征，记为$\bm{f},\bm{g}\in\mathbb{R}^{h\times w\times 256}$，$h=H/8$、$w=W/8$；另一个结构相同的上下文编码器只处理第一帧，输出的上下文特征在每一轮都注入更新模块；两者的区别只在归一化方式，特征编码器用实例归一化，上下文编码器用批归一化。第二部分是相关体：把两帧各有的$m=hw$个位置两两做内积，
\begin{equation}
  c_{i,j}=\sum_{k=1}^{d}f_{i,k}g_{j,k}\text{，}
  \label{eq:vision-video-all-pairs}
\end{equation}
得到可以视为$\mathbb{R}^{h\times w\times h\times w}$的四维张量，前两维是第一帧的查询位置，后两维是第二帧的候选位置。再沿后两维分别以核大小1、2、4、8做平均池化，得到四层相关金字塔，前两维保持不变。第三部分是迭代更新：设第$r$轮光流为$\bm{u}^{(r)}$，对第一帧每个位置$\bm{p}$，以第二帧中的$\bm{p}+\bm{u}^{(r)}(\bm{p})$为中心，在每一层取半径$\rho_{\mathrm{l}}=4$的局部网格读取相关值；第$\ell$层的读取中心要除以$2^\ell$，中心不是整数时用双线性插值。论文的公式以$L_1$距离写出这一邻域，公开实现取$(2\rho_{\mathrm{l}}+1)\times(2\rho_{\mathrm{l}}+1)$的方形网格，四层拼接为每个位置$4\times81=324$维的相关特征。半径在各层相同，所以越粗的层覆盖的范围越大：最粗一层上半径4对应原图256像素的范围。相关特征再与当前光流、上下文特征组合，作为更新模块的输入$\bm{x}^{(r)}$。

更新模块是卷积形式的门控循环单元。设旧隐状态为$\bm{h}^{(r)}$，二者位于同一粗网格。更新门$\bm{z}=\sigma(\operatorname{Conv}([\bm{h}^{(r)},\bm{x}^{(r)}]))$决定写入多少新信息，重置门$\bm{g}$由另一卷积及sigmoid计算，决定旧状态参与候选的比例；方括号表示沿通道拼接。候选状态为$\widetilde{\bm{h}}=\tanh(\operatorname{Conv}([\bm{g}\odot\bm{h}^{(r)},\bm{x}^{(r)}]))$，新状态为$\bm{h}^{(r+1)}=(1-\bm{z})\odot\bm{h}^{(r)}+\bm{z}\odot\widetilde{\bm{h}}$，$\odot$表示逐元素乘法。门控循环的一般性质可参考\BookChapterRef{chap:recurrent-networks}{卷五的循环网络章}。从新隐状态预测增量$\Delta\bm{u}^{(r)}$，更新为
\[
  \bm{u}^{(r+1)}=\bm{u}^{(r)}+\Delta\bm{u}^{(r)}\text{。}
\]
相关体只算一次，每轮变化的是读取中心和隐状态；各轮共享同一组门控参数。这与每次重新运行互不联系的匹配器不同，也与由粗到细方法中各级使用不同网络不同。

更新模块内部的计算顺序如下（表\ref{tab:vision-video-raft-arch}）。相关特征经$1\times1$卷积到256通道、再经$3\times3$卷积到128通道；当前光流经$7\times7$卷积到128通道、再经$3\times3$卷积到64通道；二者拼接后经一层$3\times3$卷积得到128通道的运动特征，再与上下文特征和光流本身拼接成$\bm{x}^{(r)}$。完整模型把上式中的卷积门控单元拆成两个串联的单元，第一个用$1\times5$卷积、第二个用$5\times1$卷积，隐状态均为128通道，这样在参数量不明显增加的情况下扩大了水平和竖直方向的感受范围。新隐状态经$3\times3$卷积（256通道）和$1\times1$卷积（2通道）输出增量$\Delta\bm{u}^{(r)}$。论文正文没有写明隐状态的初始化方式。整个网络约480万个参数，加上下面的上采样模块约530万；小型版本RAFT-S把各层通道减小、残差单元换成瓶颈结构，并只用一个$3\times3$的卷积门控单元，约100万个参数。

\begin{table*}[tbp]
  \centering
  \small
  \begin{tabularx}{\linewidth}{@{}llXlll@{}}
    \toprule
    阶段或层组 & 层类型 & 卷积核与步长 & 输出通道 & 重复次数 & 输出尺寸\\
    \midrule
    \multicolumn{6}{@{}l}{\emph{特征编码器与上下文编码器（结构相同，前者两帧共享）}}\\
    起始层 & 卷积 & $7\times7$，步长2 & 64 & 1 & $H/2\times W/2$\\
    第一组 & 残差单元 & 论文未给出 & 64 & 2 & $H/2\times W/2$\\
    第二组 & 残差单元 & 论文未给出，组内分辨率减半 & 128 & 2 & $H/4\times W/4$\\
    第三组 & 残差单元 & 论文未给出，组内分辨率减半 & 192 & 2 & $H/8\times W/8$\\
    输出层 & 卷积 & $3\times3$ & 256 & 1 & $H/8\times W/8$\\
    \midrule
    \multicolumn{6}{@{}l}{\emph{相关金字塔（无参数）}}\\
    全对相关 & 特征内积 & — & — & 1 & $h\times w\times h\times w$\\
    池化 & 沿后两维平均池化 & 核与步长$1,2,4,8$ & — & 4层 & $h\times w\times\frac{h}{2^\ell}\times\frac{w}{2^\ell}$\\
    \midrule
    \multicolumn{6}{@{}l}{\emph{更新模块（各轮共享参数）}}\\
    相关支 & 卷积 & $1\times1$，$3\times3$ & 256，128 & 各1 & $h\times w$\\
    光流支 & 卷积 & $7\times7$，$3\times3$ & 128，64 & 各1 & $h\times w$\\
    运动特征 & 卷积 & $3\times3$ & 128 & 1 & $h\times w$\\
    门控循环 & 卷积门控循环单元 & $1\times5$，$5\times1$ & 隐状态128 & 各1 & $h\times w$\\
    光流头 & 卷积 & $3\times3$，$1\times1$ & 256，2 & 各1 & $h\times w\times2$\\
    上采样掩码 & 两层卷积 & 论文未给出 & $8\times8\times9$ & 1 & $h\times w\times576$\\
    \bottomrule
  \end{tabularx}
  \caption{RAFT（完整版本）的逐层结构，输入为两帧$H\times W\times3$的图像，$h=H/8$、$w=W/8$。编码器的分组与通道取自论文附录的结构图；残差单元内部的卷积核与下采样位置论文未写明。更新模块在训练时展开12轮，测试轮数另行设定。}
  \label{tab:vision-video-raft-arch}
\end{table*}

粗网格上的光流还要恢复到输入分辨率。RAFT为每个高分辨率位置学习邻近九个粗网格光流的非负、和为1的权重$\omega_{\bm{p},j}$，取$\widehat{\bm{u}}(\bm{p})=8\sum_{j=1}^{9}\omega_{\bm{p},j}\bm{u}_{j}$。因粗网格一步约对应原图八个像素，前面的8不能漏掉；权重经softmax归一化，输出是局部粗位移的凸组合，因此能在运动边界处选择性地取用一侧的光流，而不是像双线性插值那样跨边界混合。

训练一条样本包含两帧、真实二维位移以及有效监督位置集合$V$。若一共展开$R$轮更新，对每一轮的上采样结果都计算误差，较晚的轮次权重更大：
\begin{equation}
  L_{\mathrm{flow}}
  =\sum_{r=1}^{R}\gamma^{R-r}
    \frac{1}{|V|}\sum_{\bm{p}\in V}
    \lVert\widehat{\bm{u}}^{(r)}(\bm{p})-\bm{u}^{*}(\bm{p})\rVert_1\text{，}
  \label{eq:vision-video-raft-loss}
\end{equation}
其中$\widehat{\bm{u}}^{(r)}$已经恢复到原图像素单位，$\bm{u}^{*}$为真实位移，$\gamma$在原实验中取0.8。对每轮都施加监督，使中间结果也朝正确方向移动，而不是只要求最后一轮正确。训练时展开$R=12$轮。真实光流很难在自然视频中大规模标注，所以原工作从随机初始化开始，先在带光流真值的合成数据FlyingChairs上训练10万步、再在FlyingThings上训练10万步，然后在Sintel、KITTI-2015、HD1K与FlyingThings的混合数据上微调10万步，最后在KITTI-2015上再微调5万步；各阶段用AdamW更新全部参数，梯度裁剪到$[-1,1]$，学习率在$10^{-4}$到$4\times10^{-4}$之间，按阶段设定。原方法还对递推中旧光流的梯度作截断：每轮只沿增量$\Delta\bm{u}$一支反向传播，训练的计算图不能理解为无限轮迭代的展开。

推理时执行规定的更新轮数，测试轮数可以与训练时不同，需另行报告；原文在Sintel上取32轮、在KITTI上取24轮，计时则按10轮报告。处理视频时，原工作还用上一对帧的光流前向投影作为下一对帧的初始化。例如汽车边缘实际向右移动5个像素，教学更新可以从0开始，依次增加3、1.5和0.5，得到$0,3,4.5,5$。这串数值只演示增量与读取位置的关系，不是RAFT收敛的保证。相关体的存储随位置数按$m^2$增长，输入分辨率必须计入资源预算；遮挡位置在第二帧中没有可见的对应，网络给出的稠密值仍只是估计。

\subsubsection{对象级的对应}
\label{sec:vision-video-association}
\label{sec:vision-video-motion-assignment}
\label{sec:vision-video-bytetrack}

多目标跟踪通常采用“先检测、后关联”的方式：检测器在每一帧独立给出若干候选框，关联步骤把它们与历史轨迹配对。像素位移在这里不够用，因为关联要回答的是身份问题。关联有两个基本要求：同一检测不能同时延续两条轨迹，同一轨迹在当前帧也不应接收两个检测；而且由于对象在运动，应该拿检测与轨迹在当前时刻的预测位置比较，而不是与上一帧的位置比较。

SORT把这两个要求分别交给卡尔曼滤波和匈牙利算法：用线性运动模型预测每条轨迹的当前框，用框重叠度形成代价，再求一对一的最优分配\scite{vision-sort2016}。这套流程本身没有需要学习的参数，跟踪质量主要取决于检测器；原文使用以VGG16为骨干的Faster R-CNN检测结果，并在验证序列上比较了不同检测器，指出检测质量对跟踪表现的影响很大。不计检测，跟踪部分在单个CPU核上每秒可处理约260帧。

运动预测部分回顾\term{卡尔曼滤波}（Kalman filter）的结论，它在线性高斯模型下按预测与观测的不确定性加权合并二者；采样间隔变化时，还需重新计算状态转移与过程噪声协方差，相关推导可参考\BookChapterRef{chap:dynamical-systems-and-state-estimation}{卷二的动力系统章}。SORT的状态为$\bm{x}=(x,y,s,a,\dot{x},\dot{y},\dot{s})^\top$，其中$(x,y)$是框中心，$s$是框面积，$a$是宽高比并假定恒定，带点的量是相应的速度；观测$\bm{z}=(x,y,s,a)^\top$由检测框给出。线性运动近似下，
\begin{equation}
  \bm{x}^{-}=\bm{F}\bm{x}^{+}_{\mathrm{old}},\qquad
  \bm{P}^{-}=\bm{F}\bm{P}^{+}_{\mathrm{old}}\bm{F}^{\top}+\bm{Q},\qquad
  \bm{z}=\bm{H}\bm{x}+\bm{\varepsilon}\text{。}
  \label{eq:vision-video-kalman-predict}
\end{equation}
$\bm{F}$把速度乘时间差加到位置与面积上，$\bm{H}$取出状态中可观测的四项，$\bm{P}$是状态协方差，$\bm{Q}$是过程噪声协方差，$\bm{\varepsilon}$是协方差为$\bm{R}$的观测噪声。匹配确定后，用创新$\bm{e}=\bm{z}-\bm{H}\bm{x}^{-}$更新
\[
  \bm{K}=\bm{P}^{-}\bm{H}^{\top}
  (\bm{H}\bm{P}^{-}\bm{H}^{\top}+\bm{R})^{-1},
  \qquad
  \bm{x}^{+}=\bm{x}^{-}+\bm{K}\bm{e}\text{，}
\]
协方差则更新为$\bm{P}^{+}=(\bm{I}-\bm{K}\bm{H})\bm{P}^{-}$。增益$\bm{K}$在预测越不确定、观测越可靠时越大，修正越依赖检测。教学一维中心预测为48、检测为50，预测与检测方差分别为4和1时，增益为$4/(4+1)=0.8$，修正中心为49.6，更新方差为$(1-0.8)\times4=0.8$。没有匹配的轨迹只做预测，并标明最近一次实测的时刻。身份由之前的分配决定，不由这个加权平均重新生成。

分配部分需要先写出要优化的问题。设预测框为$\widehat{\bm{b}}_i$、检测框为$\bm{d}_j$，代价$c_{i,j}=1-\operatorname{IoU}(\widehat{\bm{b}}_i,\bm{d}_j)$。用$x_{i,j}\in\{0,1\}$表示轨迹$i$是否接收检测$j$，一对一分配要求
\begin{equation}
  \min_{x}\sum_{i,j}c_{i,j}x_{i,j}
  \quad\text{s.t.}\quad
  \sum_j x_{i,j}\leq1,\ \ \sum_i x_{i,j}\leq1\text{。}
  \label{eq:vision-video-assignment}
\end{equation}
\term{匈牙利算法}（Hungarian algorithm）在多项式时间内求出这个问题的最优解，即共同决定哪些轨迹接收哪些检测。为了允许“不匹配”，可以添加虚拟行列并给它们设定未匹配代价，或在求解前把重叠太小的组合排除；SORT在分配后拒绝交并比低于设定下限的配对，该下限在验证数据上调定。与所有预测框的交并比都低于下限的检测被视为新目标：用检测框初始化状态，速度置零，速度分量的协方差取较大值以表达“速度尚未观测”，并经过一段试用期，积累足够匹配后才确认。未匹配的轨迹进入丢失计数，SORT的原始设置在一帧未匹配后即删除轨迹，以控制错误身份的蔓延；作者给出的理由是匀速模型对真实运动的预测很差，且重新识别不在该工作的范围内。

表\ref{tab:vision-association-cost}给出教学代价，说明为什么要共同分配。若按行先处理轨迹A，它会占用代价最小的检测1，B只能接检测2，总成本为$0.10+0.80=0.90$。共同分配则选择A接检测2、B接检测1，总成本为$0.20+0.11=0.31$。检测3在空间门限之外，留给新目标规则检查。全局分配改善的是这个既定的代价目标；代价若不能区分真实身份，最优配对也可能错误。

\begin{table}[htbp]
  \centering
  \begin{tabular}{lccc}
    \toprule
    历史轨迹 & 检测1 & 检测2 & 检测3\\
    \midrule
    A & 0.10 & 0.20 & 不允许\\
    B & 0.11 & 0.80 & 不允许\\
    \bottomrule
  \end{tabular}
  \caption{两条轨迹与三个检测的教学关联代价，越小越合适。“不允许”由配对前的门限排除。例中两个有效匹配均优于约定的未匹配选择，检测3不参加这次分配。}
  \label{tab:vision-association-cost}
\end{table}

这套流程有明确的失效情形。强相机运动会使所有框一起移动，急转弯会破坏线性外推，两辆相似车交换位置时框重叠无法区分它们。在代价中加入外观距离可以缓解后一种情形，但需要固定外观表示的训练数据、归一化、权重和门限。还有一个更常见的问题出在检测这一侧：通常只有得分高于门限的检测才进入关联，而被遮挡的目标恰恰得分偏低。

ByteTrack针对的正是这个问题\scite{vision-bytetrack2022}。直接丢弃低分检测，被部分遮挡的目标会在关联阶段缺席，轨迹因此中断或在遮挡后换成新编号；简单降低门限又会把背景误检一起放进来，产生虚假轨迹。ByteTrack的主要想法是：低分框本身不足以区分目标与背景，但与已有轨迹的预测位置吻合的低分框很可能是被遮挡的目标，因此可以用轨迹历史来筛选低分框，而不是一律接受或一律丢弃。

ByteTrack的架构由检测器和两次关联组成，流程见算法\ref{alg:vision-video-byte}。检测器对每帧输出框和得分，按高分门限$\tau_{\mathrm{high}}$分成高分组与低分组；原文取$\tau_{\mathrm{high}}=0.6$，低分组的下限$\tau_{\mathrm{low}}$在正文中出现，但论文没有给出取值。所有轨迹先用卡尔曼滤波预测当前位置。第一次关联把高分组与全部轨迹（包括处于丢失状态的轨迹）按式\eqref{eq:vision-video-assignment}分配，代价可以只用交并比，也可以加入外观距离。第二次关联只把低分组与第一次未匹配的轨迹分配，代价只用交并比，因为低分框往往对应严重遮挡或运动模糊，外观特征不可靠。两次分配中交并比小于0.2的配对都被拒绝。第二次仍未匹配的低分框视为背景直接舍弃，不为它们新建身份；第一次未匹配的高分框建立新轨迹；第二次仍未匹配的轨迹转为丢失状态，保留30帧以待重新出现，超过则删除。丢失轨迹不出现在当前帧的输出中。

\begin{algorithm}[ByteTrack的逐帧关联]
  \label{alg:vision-video-byte}
  \begin{algorithmic}[1]
    \Require 当前帧检测集合$D$，轨迹集合$T$（含丢失轨迹），高分门限$\tau_{\mathrm{high}}$
    \Ensure 更新后的轨迹集合与当前帧的输出
    \State $D_{\mathrm{high}}\gets\{d\in D:\ \text{得分}>\tau_{\mathrm{high}}\}$，$D_{\mathrm{low}}\gets D\setminus D_{\mathrm{high}}$
    \State 用卡尔曼滤波预测$T$中每条轨迹的当前框
    \State 第一次关联：$T$与$D_{\mathrm{high}}$按交并比（可加外观距离）求一对一分配，匹配的轨迹用检测更新
    \State 记未匹配的高分检测为$D_{\mathrm{remain}}$，未匹配的轨迹为$T_{\mathrm{remain}}$
    \State 第二次关联：$T_{\mathrm{remain}}$与$D_{\mathrm{low}}$只按交并比分配，匹配的轨迹用检测更新
    \State 仍未匹配的低分检测舍弃；仍未匹配的轨迹标为丢失，丢失超过30帧的删除
    \State 为$D_{\mathrm{remain}}$中的每个检测建立新轨迹
    \State 输出未处于丢失状态的轨迹的框与编号
  \end{algorithmic}
\end{algorithm}

ByteTrack的学习只发生在检测器上。检测器是以YOLOX-X为骨干的YOLOX，从COCO预训练权重初始化；在MOT17测试设置下，原工作在MOT17、CrowdHuman、Cityperson和ETHZ的合并数据上训练80轮，输入尺寸为$1440\times800$，多尺度训练时短边在576到1024之间变化，并使用Mosaic与Mixup增广；优化器为随机梯度下降，权重衰减$5\times10^{-4}$、动量0.9，初始学习率$10^{-3}$，1轮预热后按余弦退火，8块GPU上批大小为48。损失就是检测器自身的分类与框回归目标。两次关联没有可学习参数，高低分门限、拒绝配对的交并比下限和最大丢失帧数都是事先选定的超参数；原文的消融显示，与SORT相比，ByteTrack的结果对$\tau_{\mathrm{high}}$在0.2到0.8之间的变化更不敏感，因为第二次关联几乎考虑了所有检测框。因此，检测器的网络训练与使用时的两次分配是两项过程，关联改变的是框状态和身份记录，没有在线重新训练检测网络。

设高分门限为0.6，汽车A、B的当前检测得分分别为0.9和0.3。第一次只匹配A，B暂时未匹配。设B的预测框与低分B框的交并比为0.7，与另一个得分0.2的背景框为0.05，第二次允许交并比至少0.5，则前者可匹配，后者在求分配前被排除。B用检测更新位置、重置丢失计时并保留原编号；背景框没有合适的历史，被舍弃。图\ref{fig:vision-bytetrack-association}展示这种条件关系。低分并不等于背景，也不等于一定存在目标，轨迹历史决定它的用途。若下一帧连低分B也消失，B转为丢失状态；完全没有检测候选时，低分阶段也没有新证据可用，持续的运动预测可以等待目标重新出现，却不能作为已经观测到目标的证明。

\begin{figure}[htbp]
  \centering
  % 构造示例：横向位置仅示意关联，不表示实测轨迹。
\begin{tikzpicture}[x=1cm,y=1cm,font=\figurefont\small,
  arr/.style={-{Stealth[length=2mm]},BookInk,line width=0.8pt}]
  \node[anchor=west] at (0,3.3) {上一帧};
  \node[anchor=west] at (3.5,3.3) {当前预测与检测};
  \node[anchor=west] at (7.3,3.3) {关联结果};
  \foreach \y/\name/\col in {2.3/A/BookBlue,0.9/B/BookTeal}{
    \draw[\col,line width=1pt] (0.2,\y-0.35) rectangle (1.6,\y+0.35);
    \node at (0.9,\y) {\name};
    \draw[arr] (1.8,\y) -- (3.4,\y);
    \draw[\col,dashed,line width=0.9pt] (3.6,\y-0.45) rectangle (5.1,\y+0.45);
    \draw[\col,line width=1.1pt] (3.8,\y-0.35) rectangle (5.2,\y+0.35);
    \node at (4.5,\y) {\name};
    \draw[arr] (5.4,\y) -- (7.1,\y);
  }
  \node[above] at (6.2,2.3) {高分0.9};
  \node[above] at (6.2,0.9) {低分0.3};
  \node[anchor=west] at (7.3,2.3) {第一次延续A};
  \node[anchor=west] at (7.3,0.9) {第二次延续B};
  \draw[BookMuted,line width=0.9pt] (3.8,-0.8) rectangle (5.2,-0.1);
  \node at (4.5,-0.45) {背景};
  \draw[arr] (5.4,-0.45) -- (7.1,-0.45);
  \node[above] at (6.2,-0.45) {低分0.2};
  \node[anchor=west] at (7.3,-0.45) {无匹配，舍弃};
\end{tikzpicture}

  \caption{两次关联的教学序列。第一次用高分框延续A，第二次只用位置相符的低分框延续B。虚线框表示预测位置，实线框表示当前检测；右侧孤立低分框没有可信历史而被舍弃，不据此新建身份。}
  \label{fig:vision-bytetrack-association}
\end{figure}

表\ref{tab:vision-video-mot-pipeline}按步骤对照两种关联流程。二者都不含需要学习的关联网络，所以表中列出的是流程中的状态维度、门限和保留规则，而不是网络层。

\begin{table*}[tbp]
  \centering
  \small
  \begin{tabularx}{\linewidth}{@{}lXX@{}}
    \toprule
    步骤 & SORT & ByteTrack\\
    \midrule
    检测 & Faster R-CNN（VGG16骨干）逐帧检测 & YOLOX（YOLOX-X骨干）逐帧检测，输入$1440\times800$\\
    检测分组 & 不分组 & 按$\tau_{\mathrm{high}}=0.6$分为高分组与低分组，下限$\tau_{\mathrm{low}}$论文未给出\\
    运动预测 & 七维状态$(x,y,s,a,\dot{x},\dot{y},\dot{s})$，匀速卡尔曼滤波 & 卡尔曼滤波预测，状态设计沿用SORT一类方法，论文未单列\\
    第一次关联 & 代价$1-\operatorname{IoU}$，匈牙利算法；交并比低于$\operatorname{IoU}_{\min}$的配对拒绝，下限在验证集上调定，论文未给出取值 & 高分检测与全部轨迹（含丢失轨迹），代价为交并比或外观距离；交并比小于0.2的配对拒绝\\
    第二次关联 & 无 & 低分检测与第一次未匹配的轨迹，代价只用交并比\\
    新建轨迹 & 与所有预测框交并比都低于下限的检测，速度置零，经试用期确认 & 第一次关联后仍未匹配的高分检测；未匹配的低分检测舍弃\\
    删除轨迹 & 连续1帧未匹配即删除 & 未匹配轨迹转为丢失，保留30帧后删除\\
    可学习部分 & 仅检测器 & 仅检测器\\
    \bottomrule
  \end{tabularx}
  \caption{SORT与ByteTrack逐帧关联流程的对照。两者的关联部分都没有网络层，差别在检测怎样进入分配、轨迹怎样新建与删除；门限取各自论文在MOT基准上的设置。}
  \label{tab:vision-video-mot-pipeline}
\end{table*}

比较SORT式的关联与ByteTrack，二者使用相同的卡尔曼预测和一对一分配，ByteTrack改变的只是哪些检测进入分配以及分几次进入，以及丢失轨迹保留多久。这一改变换来了遮挡期间的轨迹连续性，代价是多了一组需要调节的门限，并且在低分框与轨迹恰好重叠的背景位置上可能把误检接到真实轨迹上。两者共同的局限是：框重叠只在相邻帧位移较小时可靠，长时间遮挡后的重新识别需要外观表示或\ref{sec:vision-video-recovery}节的恢复规则。

\subsubsection{表面点级的对应}
\label{sec:vision-video-d4rt}

移动相机拍摄行走者时，手臂在图像中向右移动，可能是手臂在动，也可能是相机在转。RAFT给出的是两幅图像之间的位置差，对象编号给出的是实例的连续性，两者都没有给出共同三维坐标中的运动。要分开相机运动和物体运动，以往的做法通常分别运行深度估计、相机位姿估计和点跟踪模型，或对每个场景单独做优化。前者的各项结果彼此不一定一致，后者对每段新视频都要花费大量迭代。

D4RT的主要想法是把完整视频只编码一次，再用一个轻量解码器回答“某个表面点在某时刻、以某帧相机为参照时位于何处”的查询\scite{vision-d4rt-query2025}。深度、点轨迹和相机参数都可以表示为这类查询的特例，因此它们来自同一个表示，也就能相互一致；每条查询独立解码，又使查询数量可以按需增减。

模型架构由编码器和查询解码器组成，配置见表\ref{tab:vision-video-d4rt-config}。编码器是视觉Transformer，帧内局部自注意力层与跨帧全局自注意力层交替排列。为支持任意宽高比，输入视频先缩放成固定的正方形分辨率再切成$2\times16\times16$的时空小块；原始宽高比另编码为一个词元，与视频词元一起送入Transformer。主训练的48帧、$256\times256$输入因此得到$24\times16\times16=6144$个视频词元。视频$\bm{V}$经编码器得到共享的视频特征$\bm{F}=E_\theta(\bm{V})$。一条查询为$q=(u,v,t_{\mathrm{src}},t_{\mathrm{tgt}},t_{\mathrm{cam}})$：$(u,v)$是来源帧上的归一化像素位置，$t_{\mathrm{src}}$指定在哪帧选点，$t_{\mathrm{tgt}}$指定询问这个点在哪个时刻的位置，$t_{\mathrm{cam}}$指定用哪帧的相机坐标表示结果。查询词元由二维坐标的傅里叶特征、三个离散时间嵌入及来源点周围$9\times9$的RGB小块编码相加而成：小块提供所选表面点的局部外观，时间嵌入说明要查询的状态与参考系。解码器是一个小型的交叉注意力Transformer，让每个查询词元交叉读取$\bm{F}$，最后输出13维：三维位置$\bm{p}(q)=D_\phi(q,\bm{F})\in\mathbb{R}^{3}$、二维位置（2维）、可见性（1维）、点运动（3维）、表面法向（3维）和置信度（1维）。查询之间没有自注意力，所以增减另一批查询不会改变当前查询的结果。选点时刻、点状态时刻和坐标系时刻是三项独立条件，不能只用一个“当前帧”省略它们。

不同的查询组合给出不同的几何量。固定来源点和$t_{\mathrm{cam}}$、改变$t_{\mathrm{tgt}}$，得到该点在共同坐标系中的运动轨迹；三个时间都等于同一帧时，$\bm{p}=(p_x,p_y,p_z)^\top$的第三项就是该像素的深度。在主点为$(0.5,0.5)$的归一化针孔模型下，投影满足$u=0.5+f_xp_x/p_z$、$v=0.5+f_yp_y/p_z$，所以可从多个非退化点算出$f_x=p_z(u-0.5)/p_x$和$f_y=p_z(v-0.5)/p_y$，再取中位数；靠近分母为零的点不能参与。相机外参则由同一组点在两个参考系中的坐标恢复：固定来源与目标时刻，只把$t_{\mathrm{cam}}$分别设为$i$和$j$，得到两组对应点$\bm{a}_k,\bm{b}_k\in\mathbb{R}^3$，求最贴合的刚性变换$\bm{b}_k\approx\bm{R}\bm{a}_k+\bm{t}$。先求均值$\bar{\bm{a}},\bar{\bm{b}}$，形成$\bm{M}=\sum_k(\bm{a}_k-\bar{\bm{a}})(\bm{b}_k-\bar{\bm{b}})^\top$；若$\bm{M}=\bm{U}\bm{\Sigma}\bm{V}^{\top}$，取$\bm{R}=\bm{V}\operatorname{diag}(1,1,\det(\bm{V}\bm{U}^{\top}))\bm{U}^{\top}$、$\bm{t}=\bar{\bm{b}}-\bm{R}\bar{\bm{a}}$。去均值先消去平移，奇异值分解求出使总平方误差最小的旋转，行列式修正保证结果不是镜像反射。单目结果的尺度依赖监督与评价标定，不能未经核实就把距离记为米。

训练一条样本由视频和一批有几何监督的查询组成，目标包括三维点、二维位置、法向、可见性或点运动，某项标注缺失时不计算该项误差。主要的三维误差先让预测点集、真实点集分别除以各自的平均深度，再逐坐标施加$g(x)=\operatorname{sign}(x)\log(1+|x|)$后比较$L_1$差。按平均深度归一化减轻了单目尺度不确定的影响，带符号对数减轻了远处点的数值支配；二者都改变了误差口径，不能把训练目标称作原始米制距离误差。辅助项分别比较二维坐标、法向余弦、可见性的二元交叉熵和运动向量，并对三维误差加入学习的置信权重及防止置信全降为零的惩罚。原工作用VideoMAE V2的预训练权重初始化编码器（VideoMAE的遮蔽重建思路见\ref{sec:vision-ivm-videomae}节），附录的消融显示这比随机初始化有明显提升。训练数据混合了BlendedMVS、Co3Dv2、Dynamic Replica、Kubric、MVS-Synth、PointOdyssey、ScanNet、ScanNet++、TartanAir、Virtual KITTI、Waymo Open等公开数据及内部数据。主训练用48帧、$256\times256$的视频，每样本解码2048条随机及区域偏置抽样的查询，在64个TPU芯片上以每芯片批大小1、用AdamW训练50万步，约两天完成，编码器与解码器端到端更新。

\begin{table}[htbp]
  \centering
  \small
  \begin{tabularx}{\linewidth}{@{}lXX@{}}
    \toprule
    配置项 & 编码器 & 查询解码器\\
    \midrule
    结构 & 视觉Transformer（ViT-g变体），帧内局部与全局自注意力交替 & 交叉注意力Transformer，查询之间无自注意力\\
    图像块大小 & $2\times16\times16$（时间$\times$高$\times$宽） & 每条查询一个词元，含$9\times9$像素小块的编码\\
    词元数 & 6144个视频词元加1个宽高比词元（48帧、$256\times256$） & 训练时每样本2048条查询\\
    层数 & 40 & 8\\
    宽度、头数、MLP倍率 & 论文未给出 & 论文未给出\\
    参数量 & 约10亿 & 约1.44亿\\
    \bottomrule
  \end{tabularx}
  \caption{D4RT主模型的配置。消融实验另用ViT-L编码器配6层解码器；编码器规模从ViT-B（约9000万参数）增大到ViT-g时，深度估计与相机旋转误差等指标明显改善。}
  \label{tab:vision-video-d4rt-config}
\end{table}

使用时参数固定，同一个编码结果可以回答新的点与时间查询。要得到密集轨迹，可以把已被可见轨迹覆盖的位置标为已访问，只从未访问位置发起新查询，以减少重复解码；被遮挡的位置不能仅因模型输出了一个点就标成已看到。我们用同一个手腕点作教学检查。固定参考系为第0帧，时刻0、1、2的点分别为$(0,0,2)$、$(0.1,0,2)$、$(0.2,0,2)$，每步横向位移0.1个共同尺度单位。若时刻1相机向右平移0.05而方向不变，该点在时刻1相机坐标中的横向值为$0.1-0.05=0.05$，对应归一化投影为$0.5+1\times0.05/2=0.525$。只把各帧相机坐标相减，会漏掉相机运动的贡献。相机估计有误时，三维轨迹可能平滑却不准确；表面点编号也不能替代整个人的永久身份。VGGT-$\Omega$把前馈多视图重建扩展到静态与动态场景，可作为与\ref{sec:vision-recognition-vggt}节VGGT的衔接入口\scite{vision-vggt-omega2026}，其结果同样需要说明坐标、尺度和可见性。

三种粒度的对应在表\ref{tab:vision-video-correspondence-compare}中对照。从RAFT到对象级关联，对应单位由像素变为对象框，换来了身份和多目标管理，丢掉了框内的细节运动；SORT与ByteTrack的差别只在检测进入分配的方式。从RAFT到D4RT，对应单位由图像位置变为三维表面点，换来了相机与物体运动的分离以及跨越多帧的长程对应，代价是需要整段视频输入、大规模几何监督和更高的编码成本。

\begin{table}[htbp]
  \centering
  \small
  \begin{tabularx}{\linewidth}{@{}lXXXX@{}}
    \toprule
    比较项 & RAFT & SORT式关联 & ByteTrack & D4RT\\
    \midrule
    对应单位 & 像素 & 检测框与轨迹 & 检测框与轨迹 & 三维表面点\\
    输入 & 两帧图像 & 每帧检测 & 每帧高低分检测 & 整段视频与查询\\
    输出 & 稠密二维位移 & 身份编号与框 & 身份编号与框 & 指定时刻与参考系下的三维点\\
    可学习部分 & 全部网络 & 仅检测器 & 仅检测器 & 编码器与解码器\\
    遮挡时 & 给出估计值 & 只做运动预测 & 低分框可延续轨迹 & 输出可见性，位置仍为估计\\
    \bottomrule
  \end{tabularx}
  \caption{三种粒度跨帧对应的对照。比较固定为同一段车辆或行人视频；“可学习部分”指需要用标注数据训练的参数，门限等超参数另由验证选择。}
  \label{tab:vision-video-correspondence-compare}
\end{table}

\subsection{跟踪状态的更新}
\label{sec:vision-video-state-update}

对应建立之后，还要决定把什么写入状态、以后怎样读取。对框轨迹，状态就是卡尔曼滤波的均值与协方差；但抠取汽车、替换背景或计算精细轮廓需要逐帧掩码，这时状态还要包含目标的像素外观。\term{视频目标分割}（video object segmentation, VOS）给定目标在某些帧的掩码或提示，要求输出它在其余各帧的掩码。这里的难点是记忆：只保留上一帧，遮挡之后就无从恢复；保留所有帧，存储随视频长度线性增长；而写进记忆的预测掩码一旦有错，错误还会传给后续帧。

两种代表性的回答都把历史保存为可以按内容读取的键值记忆，用当前帧的特征去查询历史，差别在于记忆怎样组织、初始条件从哪里来。XMem只接受首帧的完整掩码，重点在于按时间尺度分层组织记忆，让长视频的存储有界；SAM 2允许从点或框开始，并在任意帧追加纠正提示，重点在于把交互式提示与流式记忆结合起来。

\subsubsection{XMem}
\label{sec:vision-video-xmem}

在XMem之前，基于记忆网络的视频目标分割方法每隔若干帧就把一帧的特征存入记忆，显存随视频长度线性增长，长视频只能丢弃早期记忆；只保留上一帧状态的循环方法占用固定存储，却容易随时间漂移。XMem借用认知心理学中Atkinson--Shiffrin的多级记忆模型，主要想法是用三种不同时间尺度的记忆分工：快速更新的感知记忆保存最近的变化，工作记忆保存近期较完整的特征，长期记忆用少量原型保存反复有用的历史\scite{vision-xmem2022}。工作记忆满了以后，把其中常被读取的内容压缩进长期记忆，从而在有界存储下保留长期信息。

模型架构中有两个编码器，各组件的配置见表\ref{tab:vision-video-xmem-config}。两者都以ResNet为基础，去掉分类头和最后一个卷积阶段，输出步长为16。查询编码器（ResNet-50）只读当前图像，其输出再经三个独立的$3\times3$卷积，分别得到键、缩放项和选择项；键在每个位置是一个$d_k=64$维向量。值编码器（ResNet-18）同时读取图像与目标掩码，产生$d_v=512$维的对象值，即“这个位置的外观加上它是否属于目标”。当前帧的键作为查询$\bm{q}_i\in\mathbb{R}^{d_k}$，与记忆中的历史键$\bm{k}_j\in\mathbb{R}^{d_k}$比较，然后把对应的对象值$\bm{v}_j\in\mathbb{R}^{d_v}$加权读出：
\begin{equation}
  a_{i,j}=\frac{\exp s_{i,j}}{\sum_\ell\exp s_{i,\ell}},
  \qquad
  \bm{r}_i=\sum_j a_{i,j}\bm{v}_j\text{。}
  \label{eq:vision-video-memory-read}
\end{equation}
相似度不是简单的内积。XMem用查询的通道选择权重$e_{i,c}\in[0,1]$和记忆项的缩放系数$\alpha_j\geq1$调节键距离：
\[
  s_{i,j}=-\alpha_j\sum_c e_{i,c}(k_{j,c}-q_{i,c})^2\text{。}
\]
通道选择决定这次比较关注哪些外观维度，缩放系数决定该记忆项在距离增大时衰减多快，两者都由学习的模块产生；实际读取还可以只保留得分最高的若干项。选择项经sigmoid落在$[0,1]$，缩放项经“平方加1”保证不小于1。读出值$\bm{r}_i$与感知记忆的隐状态拼接后送入解码器；解码器每次上采样2倍，直到步长4，每一级融合查询编码器对应层的跳跃连接，最后由一个$3\times3$卷积输出单通道的对数几率，再双线性上采样到输入分辨率，得到当前帧的掩码。若两个候选的读取权重为0.8和0.2，值中某个分量分别为0.9和0.1，则该分量读出0.74，再由解码器与当前特征共同形成像素判断，不能直接把0.74当作最终掩码概率。

三种记忆的写入规则不同。感知记忆是64通道的隐状态，每帧由门控循环单元读取解码器的多尺度特征更新；每当写入新的工作记忆帧时，再用值编码器的特征经另一个门控循环单元做一次深度更新。工作记忆每隔$r$帧写入一次当前帧的键与值，原文取$r=5$。长期记忆只由巩固过程写入。巩固并不是把所有旧帧平均成一幅图。设工作记忆最多保存$T_{\max}$个记忆帧，达到容量时仍保留首帧标注和最近$T_{\min}-1$个帧，把其余帧中的键值位置作为候选。候选$j$的使用量是它历次读取权重的累计和除以在库时长，按使用量选出$P$个原型键；除以时长，使较早入库本身不会成为被选中的唯一原因。每个原型再查询全部候选，汇聚出原型值
\begin{equation}
  \bm{v}^{\mathrm{p}}_k
  =\sum_j
   \frac{\exp s(\bm{k}^{\mathrm{c}}_j,\bm{k}^{\mathrm{p}}_k)}
        {\sum_\ell\exp s(\bm{k}^{\mathrm{c}}_\ell,\bm{k}^{\mathrm{p}}_k)}
   \bm{v}^{\mathrm{c}}_j\text{，}
  \label{eq:vision-video-xmem-consolidation}
\end{equation}
其中上标$\mathrm{c}$表示候选，$\mathrm{p}$表示原型，相似度沿用前述学习距离。选键控制保留多少位置，汇聚值则让未被选为原型的候选仍能贡献信息。原型加入长期记忆后才删除相应的工作候选；长期记忆达到自己的容量后，再按使用量淘汰。原文默认$T_{\min}=5$、$T_{\max}=10$、$P=128$，长期记忆至多保存1万个元素，读取时只保留相似度最高的30项；在480p视频上，整个模型占用的显存不超过1.4 GB。

\begin{table}[htbp]
  \centering
  \small
  \begin{tabularx}{\linewidth}{@{}lXX@{}}
    \toprule
    组件 & 结构与规模 & 输入与输出\\
    \midrule
    查询编码器 & ResNet-50去掉分类头与第五阶段，三个$3\times3$卷积投影 & 当前图像$\to$步长16的键（64维）、缩放项与选择项\\
    值编码器 & ResNet-18去掉分类头与第五阶段 & 图像与掩码$\to$步长16的对象值（512维）\\
    感知记忆 & 门控循环单元，隐状态64通道，每帧更新 & 解码器多尺度特征$\to$隐状态\\
    工作记忆 & 每5帧写入一帧的键与值，容量$T_{\max}=10$帧 & 首帧标注及最近$T_{\min}-1=4$帧常驻\\
    长期记忆 & 每次巩固选$P=128$个原型，上限1万个元素 & 原型键与汇聚的原型值\\
    解码器 & 逐级2倍上采样至步长4，融合跳跃连接，$3\times3$卷积输出 & 读出值与隐状态$\to$单通道对数几率\\
    \bottomrule
  \end{tabularx}
  \caption{XMem各组件的配置。编码器沿用标准ResNet（见\ref{sec:vision-ivm-resnet}节），只截去最后一个卷积阶段；读取时保留相似度最高的30项。}
  \label{tab:vision-video-xmem-config}
\end{table}

教学中三个车身候选的使用量设为0.6、0.3、0.1，只保留第一个为原型。若该原型查询候选所得权重为0.6、0.3、0.1，对应某值分量为0.9、0.7、0.1，长期值为$0.54+0.21+0.01=0.76$，并非直接复制首项的0.9。后来遮挡解除，当前查询同时读取近期值和长期原型；只有长期键仍与真实车身相符，这个0.76才可能提供有用信息。被压掉的原始车牌像素不能由这个分量复原。

学习这套读取需要逐帧的分割标注。一条训练样本从标注视频抽取若干有时间间隔的帧，以首帧掩码建库，依次读取、解码，把预测作为后续帧的条件，最后比较各帧掩码。损失是等权组合的自举交叉熵（bootstrapped cross entropy）与Dice损失：前者只对损失最大的一部分像素计算交叉熵，使训练集中于难分的边界；后者以$1-(2\sum_i p_im_i+\epsilon)/(\sum_i p_i+\sum_i m_i+\epsilon)$约束预测与真实区域的整体重叠，$p_i$是像素概率，$m_i\in\{0,1\}$是标签，$\epsilon>0$处理空区域。原工作先用变形静态图像生成的三帧合成序列预训练15万步（批大小16），再在YouTubeVOS与DAVIS上主训练11万步（批大小8），从标注视频抽取八帧并逐渐扩大采样间隔；八帧而不是三帧，是为了让感知记忆不过分依赖初始状态。每次训练至多随机保留三个历史帧作为工作记忆，深度更新以0.2的概率执行，对应$r=5$。两阶段都用AdamW，学习率$10^{-5}$、权重衰减0.05，主训练在8万步后学习率降为十分之一，全部网络一起更新。巩固与淘汰是推理时的容量管理规则，不需要为长期原型求梯度。使用阶段网络参数固定，写入的是特征、预测掩码和隐状态；错误掩码写入会污染后续读取，原型压缩也可能丢掉车牌等细节，有限容量不提供长期身份正确的保证。

\subsubsection{SAM 2}
\label{sec:vision-video-sam2}

XMem要求首帧给出完整而准确的掩码，这在实际标注中代价不小；一旦后续某帧分错，也没有让用户纠正的入口。另一方面，\ref{sec:vision-recognition-sam}节的SAM可以从点或框提示分割单幅图像，却没有跨帧记忆。SAM 2把二者结合：让用户在任意帧用点、框或掩码给出提示，模型把结果写入流式记忆并传播到其余帧；某帧分错时，用户在该帧追加提示，修正就会通过记忆影响后续帧\scite{vision-sam2}。本节依据论文2024年10月的修订版，其中更新后的模型对应公开实现中的SAM 2.1。

模型架构按帧依次处理，配置见表\ref{tab:vision-video-sam2-config}。图像编码器是经MAE预训练的分层视觉Transformer Hiera，对每帧只运行一次，输出步长4、8、16、32的四级特征\scite{vision-hiera2023}。特征金字塔网络融合步长16与32两级，得到每帧的图像嵌入，对$1024\times1024$的输入其形状为$256\times64\times64$；步长4与8两级不进入记忆注意力，而是直接加到掩码解码器的上采样层中，用来恢复边界细节。记忆注意力模块由$L=4$个Transformer块组成，每块依次作当前帧内的自注意力、对记忆库的交叉注意力和MLP，空间位置用二维旋转位置编码表示。记忆库包括两类内容：空间记忆是若干帧的特征网格，对象指针是从这些帧的掩码输出词元得到的短向量，概括“目标是什么”。二者空间分辨率不同，前者提供位置细节，后者提供对象层面的条件。读过记忆的特征与提示编码器的输出一起进入掩码解码器，解码器结构沿用SAM，输出目标掩码、预测IoU，以及判断目标在当前帧是否出现的存在得分。最后，记忆编码器把预测掩码下采样，与未读取记忆的图像特征相加并融合，得到这一帧的空间记忆写入记忆库。记忆库以先进先出队列保留至多$N$个近期非提示帧，原文默认$N=6$；带用户提示的帧放在另一个至多$M$帧的队列中，论文没有给出$M$的取值。空间记忆被投影到64维以减小存储，256维的对象指针拆成4个64维词元参与交叉注意力。近期帧的记忆加入时间位置编码，以表达短时运动；提示帧的记忆不加，因为它们在推理时可能来自与训练时很不相同的时间范围。被预测为遮挡的帧，其记忆上还加一个学习的遮挡嵌入。记忆总量有界，所以长视频的存储不随长度增长。

\begin{table*}[tbp]
  \centering
  \small
  \begin{tabularx}{\linewidth}{@{}lXX@{}}
    \toprule
    组件 & 配置 & 输出\\
    \midrule
    图像编码器（Hiera-B+，默认） & 四阶段通道112、224、448、896，块数2、3、16、3，头数2、4、8、16；第12、16、20块用全局注意力，其余为窗口注意力；分类配置约7000万参数 & 步长4、8、16、32的四级特征\\
    图像编码器（Hiera-L） & 四阶段通道144、288、576、1152，块数2、6、36、4，头数2、4、8、16；第23、33、43块用全局注意力；分类配置约2.14亿参数 & 同上\\
    特征金字塔 & 融合步长16与32两级 & $256\times64\times64$的图像嵌入\\
    记忆注意力 & 4个块，每块自注意力、交叉注意力、MLP，二维旋转位置编码 & 以记忆为条件的$256\times64\times64$特征\\
    记忆编码器 & 卷积下采样预测掩码，与图像嵌入相加，再经轻量卷积融合 & 64维空间记忆\\
    记忆库 & 至多$N=6$个近期帧与$M$个提示帧（$M$论文未给出），对象指针拆为4个64维词元 & 供交叉注意力读取的键值\\
    掩码解码器 & 双向Transformer块，沿用SAM；另加遮挡预测词元与MLP头 & 候选掩码、预测IoU、出现得分\\
    \bottomrule
  \end{tabularx}
  \caption{SAM 2的组件配置，输入分辨率$1024\times1024$。Hiera的通道、块数、头数与参数量取自Hiera论文的配置表（参数量按图像分类配置计），全局注意力的位置取自SAM 2论文；Hiera的MLP倍率在两篇论文的配置表中都未列出。}
  \label{tab:vision-video-sam2-config}
\end{table*}

两点设计需要特别注意。其一，记忆编码器直接使用预测掩码，所以错误掩码会改变保存的内容，而不仅是在显示上留一个错误轮廓。其二，单个点击可能有歧义（点在车门上，指的是车门还是整辆车），解码器为此输出多个候选掩码；若没有新提示消除歧义，就按预测IoU选一个传播，不能同时把多个不相容的解释称作已确认的目标。目标不存在或被遮挡时，存在得分让模型输出空掩码，但状态仍然保留。

训练分为两个阶段。先在SA-1B静态图像上预训练约9万步，批大小256，图像缩放为$1024\times1024$。再联合使用图像与视频分割数据训练约15万步，视频数据包括论文发布的SA-V，每次迭代按数据规模的比例抽取一整批图像或一整批视频。为了让网络学会利用交互，训练用8帧序列模拟交互：至多选两个帧给提示，初始提示以50\%的概率取真实掩码、25\%取真实区域内的前景点、25\%取包围框；之后根据当前预测与真实掩码的差异采样纠正点击，网络在收到点击后重新计算，从而学到“这处遗漏如何被纠正”。掩码的焦点损失与Dice损失分别约束难分像素和整体重叠，预测IoU用绝对误差约束，存在得分用交叉熵约束，四项的权重之比为$20:1:1:1$；多候选输出时只监督分割损失最小的那张掩码。真实目标不在帧中时，只监督存在得分，不要求凭空生成目标掩码。图像编码器从MAE预训练的Hiera权重初始化，其余模块随联合训练学习，用AdamW按既定日程更新。原文最后还在编辑最多的困难视频上用16帧序列微调5万步，学习率减半，并冻结图像编码器以放进显存。使用时写记忆库不是重新训练这些参数；论文报告Hiera-B+与Hiera-L版本的速度分别为每秒43.8帧和30.2帧。

在四帧教学序列中，第1帧前景点选中汽车，第2帧读取首帧与近期记忆，第3帧车尾漏分。在第3帧追加车尾前景点后，解码器重新输出掩码，这张新掩码经记忆编码成为后续条件；重跑第4帧时才能利用它。若第3帧同时误包含另一辆车，还需负点或更明确的框来排除，追加一个前景点不必解决全部歧义。纠正影响多远，取决于实际重新传播了哪些帧。在线前向传播只能读已到达的帧，离线交互则可从不同提示时刻向指定方向重新处理给定视频，二者的观测范围和等待需要分别记录。用户点选一辆车后的持续分割，与自动发现并命名所有车辆是不同任务。

XMem与SAM 2都用键值记忆把历史带到当前帧，比较见表\ref{tab:vision-video-memory-compare}。SAM 2沿用了XMem这类方法的“当前特征查询历史记忆”的读法，改变的是初始条件和记忆的组织：用提示代替必须完整的首帧掩码，用对象指针补充空间记忆，用存在得分显式处理目标消失，并允许在任意帧纠正。换来的是交互能力和图像视频统一的训练，代价是更大的图像编码器与训练数据，以及记忆只保留有限近期帧，长期信息依赖提示帧和对象指针；XMem的分层巩固则在不需要用户参与的长视频中更直接地控制存储。

\begin{table}[htbp]
  \centering
  \small
  \begin{tabularx}{\linewidth}{@{}lXX@{}}
    \toprule
    比较项 & XMem & SAM 2\\
    \midrule
    初始条件 & 首帧完整掩码 & 任意帧的点、框或掩码提示\\
    记忆内容 & 键值特征，分感知、工作、长期三级 & 近期帧与提示帧的空间记忆，加对象指针\\
    容量控制 & 工作记忆满后按使用量巩固为原型 & 只保留有限近期帧，提示帧另行保留\\
    纠错方式 & 无交互入口 & 在出错帧追加提示并重新传播\\
    目标消失 & 由解码器输出决定 & 存在得分显式判断\\
    训练数据 & 合成序列与标注视频 & 图像与视频分割数据联合训练\\
    \bottomrule
  \end{tabularx}
  \caption{两种记忆式视频目标分割模型的对照。比较固定为同一段视频中单个目标的逐帧掩码输出。}
  \label{tab:vision-video-memory-compare}
\end{table}

\subsection{遮挡与重现的处理}
\label{sec:vision-video-recovery}

前面各步都假定当前帧有可靠的观测；遮挡打破了这一前提。状态因此要区分两种来源：可靠匹配后，状态有当前观测支持；暂时失去匹配后，位置来自外推，身份来自历史。输出应能区分这两种来源。沿“可靠匹配—暂时丢失—重新匹配—终止轨迹”的顺序，等待时间内若出现候选，需要重新核验空间范围、外观和历史的一致性；超过协议规定的期限再终止轨迹，必要时将再次出现的对象作为新身份，而不强行连接。

前文的模型分别给出了这一阶段的具体做法：MOSSE用峰值旁瓣比判断遮挡并暂停更新，卡尔曼滤波在丢失期间只做预测，ByteTrack用低分框和丢失轨迹的保留期接回目标，XMem依靠长期原型在遮挡后重新读到目标，SAM 2用存在得分在遮挡期间输出空掩码。它们面对同一个两难：固定模板或冻结记忆避免把错误目标写进历史，却难以适应外观的大幅变化；更新模板或持续写入能追随姿态与光照，却可能在遮挡时吸收背景或另一个对象。实际管理可以对低可信观测停止写入、保留早期可靠模板，或允许用户纠正，但这些措施都应按具体实现验证。跨镜头的相似外观不足以证明是同一实例，历史中已经淘汰的细节也不能靠延长推理重新获得。

\section{活动与事件的判断}
\label{sec:vision-video-activities}

身份把对象在时间上连起来，活动判断还要解释对象之间发生了什么。“手、杯子和桌子”在拿起与放下时都可能出现，区别在于运动方向、接触关系和状态顺序。本节的输入是视频帧或预先提取的片段特征，必要时加上前一节的对象轨迹；输出从预设的活动词表中选取，但形式各不相同。难点在于活动的证据分布在时间上：长动作需要较广的覆盖，短暂运动需要足够密集的观测，而事件的起止边界本身也常有标注歧义。

本节按输出类型展开。给定一段已裁剪的片段，只需判断它属于哪种活动，这是片段活动的分类；面对未裁剪的长视频，还要找出每次事件的起止时间，这是事件区间的定位；在线场景中事件尚未结束就要作出判断，这是未完成事件的预判；得到多个事件后，还要解释它们之间的顺序与重复。后一种输出往往以前一种为基础：区间定位通常复用片段分类的特征，预判要在定位的框架下限制可见范围，顺序解释则读取定位出的事件记录。

\subsection{片段活动的分类}
\label{sec:vision-video-activity-evidence}
\label{sec:vision-video-activity-output}
\label{sec:vision-video-clip-label}

片段活动分类的输入是一段给定时间范围的视频，输出是这段范围内的活动类别。单标签分类假设一个片段只有一个主要类别，通常用softmax及类别交叉熵；若同一人同时站立、拿杯并说话，就需要多标签输出，可为各类别使用独立的sigmoid及二元交叉熵，再由阈值决定保留哪些标签。阈值的选择受类别不均衡与标注遗漏影响。

整段标签还没有说明是谁在做。教学画面中，左边的人拿杯、右边的人看手机，两人都站立；整帧输出“拿杯、看手机、站立”没有完成动作到人物的绑定。AVA数据集为此在关键时刻标注每个人的框及其多个原子动作\scite{vision-ava2018}。原论文的基线方法从关键帧的ResNet-50特征生成人物候选框，再在RGB与光流两路I3D特征的每个时间切片上，按同一个空间框作区域池化，两路融合后输出人物框修正和各动作的sigmoid得分。对一个人物候选，设真实多标签为$\bm{y}\in\{0,1\}^{C}$、预测为$\bm{p}\in(0,1)^C$，动作误差为$-\sum_c[y_c\log p_c+(1-y_c)\log(1-p_c)]$，各类独立计分，站立与拿杯可以同时为真；人物位置则由框回归另行监督。按固定空间框跨时刻取特征，并不等于已经跟上了人物的真实运动，大幅移动的人可能部分离开这个区域；关键时刻的多标签也不直接给出整次拿杯从开始到结束的范围。

无论单标签、多标签还是人物级输出，核心都是怎样从片段中形成活动证据，这一步的难点是时间覆盖与运动细节难以兼得。连续读取少量相邻帧能看清局部运动，却可能只看到“伸手”而看不到杯子之后是否离开桌面；在整段上稀疏取帧能覆盖全过程，却丢掉了相邻帧之间的运动。TSN选择稀疏覆盖，用分段采样和得分共识让二维网络看到完整过程；I3D选择连续片段，把二维卷积扩展到时间维，直接从相邻帧学习运动。二者都可以再加一路光流输入补充运动信息。

在深层网络之前，片段活动分类的代表做法是改进的密集轨迹（improved dense trajectories, iDT）。它的基础是密集轨迹：在8个空间尺度上（相邻尺度相差$1/\sqrt{2}$倍）每隔5个像素采样特征点，用中值滤波后的稠密光流逐帧推进，得到点的轨迹；每条轨迹至多延续15帧以避免漂移，静止的轨迹和突然大幅跳动的轨迹被删去\scite{vision-dense-trajectories2011}。沿每条轨迹取一个$32\times32$像素、15帧长的时空管道，划分为$2\times2\times3$个单元，计算四种描述子：轨迹形状（归一化位移串，30维），HOG（梯度方向直方图，96维，描述外观），HOF（光流方向直方图，108维，描述运动），以及MBH（光流水平与竖直分量各自的梯度方向直方图，192维，描述运动边界）。MBH对光流再求一次导数，局部恒定的相机平移在求导后被消去，这是它比HOF更能抵抗相机运动的原因。

改进版本进一步显式估计相机运动\scite{vision-idt2013}。在相邻两帧之间，用SURF特征匹配和从光流中采样的点得到候选对应，用RANSAC估计一个单应；人体检测框内的匹配点事先去掉，以免人的运动被当作相机运动。按单应把第二帧变形后重新计算光流，HOF与MBH都在这一变形光流上计算；变形后最大位移小于1像素的轨迹被视为相机运动造成的背景轨迹而删除。编码时，每种描述子先用主成分分析降到一半维数，再用256个高斯分量的混合模型编码为Fisher向量，长度为$2DK$（$D$为降维后的维数，$K=256$），经幂归一化与$L_2$归一化后拼接，交给线性支持向量机分类。这条路线已经有了后来深层网络的分工：外观与运动分开描述，运动依赖预先计算的光流，整段视频汇总成一个固定长度的向量。它的局限也正是下面两个模型要回答的：光流与轨迹的计算昂贵，描述子由人设计，Fisher向量编码又不保留局部特征的先后。TSN把HOG与HOF的分工换成RGB与光流两路学习的卷积网络，并借用iDT的相机运动补偿得到“变形光流”作为附加输入；I3D则试图从连续的RGB帧直接学习运动，同时保留光流一路作为补充。

\subsubsection{TSN}
\label{sec:vision-video-tsn}

TSN提出时，主流的双流卷积网络从视频中只取单帧RGB或一小段堆叠光流作为输入，只能看到很短的时间范围；在整段视频上密集取帧又计算量过大且相邻帧高度冗余。当时的动作数据集规模也较小，深层网络容易过拟合。\term{时间分段网络}（temporal segment network, TSN）的主要想法是：把视频均分为若干段，每段只取一个短片段，用同一个网络处理各片段，再把各片段的得分汇总成视频预测\scite{vision-tsn2016}。稀疏取样使计算量与视频长度无关，分段又保证取样覆盖整个时间范围；由于损失定义在汇总结果上，网络从视频标签学习，而不要求每个片段单独都表现出该活动。

模型架构如下。给定视频被均分为$K$段，原工作训练时取$K=3$；第$k$段中随机取一个片段$\bm{V}_k$，RGB分支取单帧，光流分支取若干连续帧间的水平和垂直光流，按通道堆叠。每个片段裁剪为$224\times224$，经共享参数的BN-Inception网络得到类别得分$\bm{s}_k=f_\theta(\bm{V}_k)\in\mathbb{R}^C$；光流分支堆叠的帧数论文未给出。BN-Inception是在GoogLeNet上加入批归一化并作若干改动的网络：$5\times5$卷积换成两个串联的$3\times3$卷积，$28\times28$尺度上的Inception模块由两个增加到三个，模块之间不再设单独的池化层，改由模块内步长为2的卷积和池化下采样（逐层结构见表\ref{tab:vision-video-bninception}），共约1360万个参数\scite{ioffe2015batchnorm}。它的基本单元是Inception模块：输入同时进入四条分支，即$1\times1$卷积；$1\times1$降维后接$3\times3$卷积；$1\times1$降维后接两个串联的$3\times3$卷积（代替GoogLeNet中的$5\times5$卷积）；以及平均或最大池化后接$1\times1$投影。四条分支的输出沿通道拼接，每个卷积之后都接批归一化与ReLU。不同分支看到不同大小的邻域，拼接使下一层同时获得多种尺度的特征。采用均值共识时，先汇总得分，再做softmax及视频标签损失：
\begin{equation}
  \bm{g}=\frac{1}{K}\sum_{k=1}^{K}\bm{s}_k,\qquad
  p_c=\frac{\exp g_c}{\sum_j\exp g_j},\qquad
  L=-\log p_y\text{，}
  \label{eq:vision-video-tsn}
\end{equation}
其中$y$是给定范围的真实活动类别。RGB与光流两个分支分别训练，测试时把两路的类别得分按规定权重加权融合。

均值共识为什么能用同一个视频标签训练多段，可以从梯度看出。softmax交叉熵对汇总得分的导数为$p_c-\bm{1}[c=y]$，而$g_c$对每段$s_{k,c}$的导数为$1/K$，所以
\[
  \frac{\partial L}{\partial\theta}
  =\sum_{k=1}^{K}\sum_c
    \frac{p_c-\bm{1}[c=y]}{K}
    \frac{\partial s_{k,c}}{\partial\theta}\text{，}
\]
这里$\bm{1}[c=y]$是类别一致时为1的指示量。各片段共用$\theta$，一次更新把各段的贡献相加，而不是分别训练$K$个分类器。每段收到的误差信号相同，都来自视频层面的预测；目标是降低视频层面的错误，并不要求每段单独分类都等于$y$。把视频标签复制为逐帧标签，会改变监督的含义。

学习方法上，TSN针对小数据集提出了几项做法。两个分支都从ImageNet预训练的网络出发；光流先按线性变换离散到0至255，使取值范围与RGB图像相同，再把第一层卷积核在RGB三个通道上取平均，复制到光流的各个通道，称为跨模态初始化。微调时只重新估计第一个批归一化层的统计量，冻结其余批归一化层的均值和方差，称为部分批归一化，以减轻小数据上的统计偏差；全局池化之后另加dropout，比例在RGB分支取0.8、光流分支取0.7。数据增强使用角点裁剪与尺度抖动：输入固定为$256\times340$，裁剪区域的宽和高从$\{256,224,192,168\}$中随机选取，再缩放到$224\times224$。每批在各段随机取片段，批大小256、动量0.9，以随机梯度下降最小化式\eqref{eq:vision-video-tsn}；RGB分支学习率从0.001开始、每2000步降为十分之一、4500步停止，光流分支从0.005开始、在12000和18000步各降为十分之一、20000步停止。光流用TV-L1算法提取。推理时在整段视频上取25帧RGB或25组光流，每帧取四角、中心及其水平翻转共10个裁剪，在softmax之前融合各帧与两路得分，RGB与光流的权重为$1:1.5$。这些重复读取都应计入开销；固定数量的稀疏取样仍没有覆盖短暂动作的保证。

\begin{table*}[tbp]
  \centering
  \small
  \begin{tabularx}{\linewidth}{@{}llXlll@{}}
    \toprule
    阶段或层组 & 层类型 & 卷积核与步长 & 输出通道 & 重复次数 & 输出尺寸\\
    \midrule
    conv1 & 卷积 & $7\times7$，步长2 & 64 & 1 & $112\times112$\\
    pool1 & 最大池化 & $3\times3$，步长2 & 64 & 1 & $56\times56$\\
    conv2 & $1\times1$降维（64）后接卷积 & $3\times3$，步长1 & 192 & 1 & $56\times56$\\
    pool2 & 最大池化 & $3\times3$，步长2 & 192 & 1 & $28\times28$\\
    inception 3a、3b & Inception模块，平均池化分支 & 四分支，见正文 & 256、320 & 2 & $28\times28$\\
    inception 3c & Inception模块，步长2 & 池化分支直通 & 576 & 1 & $14\times14$\\
    inception 4a至4d & Inception模块，平均池化分支 & 四分支 & 576 & 4 & $14\times14$\\
    inception 4e & Inception模块，步长2 & 池化分支直通 & 1024 & 1 & $7\times7$\\
    inception 5a、5b & Inception模块，平均与最大池化分支 & 四分支 & 1024 & 2 & $7\times7$\\
    pool & 平均池化 & $7\times7$，步长1 & 1024 & 1 & $1\times1$\\
    分类头 & TSN所加dropout与线性层 & — & $C$ & 1 & $C$\\
    \bottomrule
  \end{tabularx}
  \caption{TSN所用BN-Inception的逐层结构，输入为$224\times224$的单帧RGB；光流分支只改变第一层的输入通道数。各Inception模块内部的分支通道数见BN论文附录。原表对两个步长为2的模块记录的是模块前的空间尺寸，这里按步长写出模块输出的尺寸。}
  \label{tab:vision-video-bninception}
\end{table*}

在九秒教学视频中，伸手、拿杯和收手分别主要发生在三个时间段。连续抽取开头三帧可能都只见伸手，跨三段取样却能覆盖完整过程。图\ref{fig:vision-temporal-sampling}对照了两种覆盖范围。若三段对“拿杯、放杯”的得分依次为$(1,0)$、$(3,0)$、$(2,0)$，均值为$(2,0)$，拿杯概率约为0.881；交换这三个得分的顺序，均值完全不变。均值共识因此不编码段间次序，拿起与放下若只在顺序上不同，单靠这种汇总无法区分。

\begin{figure}[htbp]
  \centering
  % 九秒教学序列，样本线不表示实际方法的固定帧数。
\begin{tikzpicture}[x=0.82cm,y=1cm,font=\figurefont\small]
  \foreach \a/\b/\name in {0/3/伸手,3/6/拿杯,6/9/收手}{
    \fill[BookPaper] (\a,2.45) rectangle (\b,3.05);
    \node at ({(\a+\b)/2},2.75) {\name};
    \draw[BookRule,line width=0.5pt] (\a,0.2) -- (\a,3.1);
  }
  \draw[BookRule,line width=0.5pt] (9,0.2) -- (9,3.1);
  \draw[-{Stealth[length=2mm]},BookInk,line width=0.8pt]
    (0,0) -- (9.3,0);
  \node[anchor=east] at (9.3,-0.65) {时间（秒）};
  \foreach \x in {0,3,6,9}{
    \draw[BookInk,line width=0.7pt] (\x,0.05) -- (\x,-0.05);
    \node[below] at (\x,-0.08) {\x};
  }
  \node[anchor=east] at (-0.2,1.95) {跨段取样};
  \draw[BookMuted,line width=0.6pt] (0,1.95) -- (9,1.95);
  \foreach \x in {1.5,4.5,7.5}{
    \draw[BookBlue,line width=1.5pt] (\x,1.7) -- (\x,2.2);
  }
  \node[anchor=east] at (-0.2,0.95) {连续片段};
  \draw[BookMuted,line width=0.6pt] (0,0.95) -- (9,0.95);
  \fill[BookTeal!12] (3.6,0.65) rectangle (5.4,1.25);
  \foreach \x in {3.7,3.95,4.2,4.45,4.7,4.95,5.2}{
    \draw[BookTeal,line width=1pt] (\x,0.7) -- (\x,1.2);
  }
\end{tikzpicture}

  \caption{稀疏覆盖与连续片段提供的时间证据不同。图中三段和动作位置为教学设定，竖线表示采样帧。跨段采样覆盖伸手、拿杯和收手，但稀疏点间仍可能漏掉短动作；连续片段保留局部运动，却未必覆盖全程。}
  \label{fig:vision-temporal-sampling}
\end{figure}

\subsubsection{I3D}
\label{sec:vision-video-i3d}

TSN的每个RGB片段只是一帧，局部运动只能靠另一路光流提供，而光流需要单独计算。三维卷积网络可以直接从连续帧学习时空特征，但此前的三维网络（如C3D）参数多、只能从头训练，在当时的小数据集上难以充分训练。I3D的主要想法是“膨胀”：把ImageNet上训练好的二维卷积网络的每个卷积核和池化核沿时间维扩展，使网络直接继承图像网络的参数，然后在大规模视频数据Kinetics上继续训练\scite{vision-i3d2017}。这样，三维网络从一个已经能识别外观的状态出发，只需再学习时间上的变化。

膨胀需要保证初始化时的响应与原网络一致。若二维卷积核为$\bm{W}$，把它沿时间复制$q$份，并将每份设为$\bm{W}/q$。对没有运动、每帧相同的输入，远离时间填充边界的卷积响应满足
\begin{equation}
  \sum_{r=1}^{q}\frac{\bm{W}}{q}*\bm{X}
  =\bm{W}*\bm{X}\text{，}
  \label{eq:vision-video-inflation}
\end{equation}
即三维网络在静态重复的“无聊视频”上给出与原图像网络相同的激活。对一般的连续输入，响应是$\sum_{r=1}^{q}\bm{W}_r*\bm{X}_{t+r}$，只有在训练使各时间切片的权重$\bm{W}_r$分化之后，网络才可能表达先后和运动；复制本身并没有学会方向。教学中二维滤波对静态手部给出6，复制三份后每份贡献2，总和仍是6。若训练后某个标量时间核变成$(-1,0,1)$，输入局部亮度$(2,3,5)$得3，倒序$(5,3,2)$得$-3$，说明时间权重可以表达方向；这组手设权重不是I3D实际学得的参数。

模型架构以ImageNet上预训练的Inception-V1（带批归一化）为原型，逐层结构见表\ref{tab:vision-video-i3d-arch}。原工作以每秒25帧处理视频，训练时输入64帧、$224\times224$的连续RGB片段，形状为$64\times224\times224\times3$，覆盖2.56秒。所有二维卷积核与池化核都扩展为三维：第一层是步长为2的$7\times7\times7$卷积；前两个最大池化层用$1\times3\times3$的核、时间步长为1，不在时间维下采样，以免过早丢掉时间分辨率；其余两个最大池化层在时间和空间上采用对称的核与步长；最后的平均池化用$2\times7\times7$的核，再接$1\times1\times1$卷积输出类别得分。作者的考虑是：图像网络在两个空间方向上对称地扩大感受野是自然的，时间方向却应当依帧率和图像尺寸决定扩大的快慢，太快会把不同物体的边缘混在一起，太慢又抓不住场景动态。论文图中给出了几处理论感受野（时间以帧计，空间以像素计）：第一个池化后为$7\times11\times11$，第二个池化后为$11\times27\times27$，第三个池化后为$23\times75\times75$，第四个池化后为$59\times219\times219$，平均池化后为$99\times539\times539$。膨胀的Inception模块有四条分支：$1\times1\times1$卷积；$1\times1\times1$卷积后接$3\times3\times3$卷积的两条分支；$3\times3\times3$最大池化后接$1\times1\times1$卷积；输出沿通道拼接，每个卷积后接批归一化与ReLU。网络输出$\mathbb{R}^{t'\times h'\times w'\times d'}$的时空特征，经平均池化与分类层得到类别得分。I3D也采用双流：另一路以连续光流片段为输入、结构相同，两路分别训练，测试时把两路的预测取平均；双流合计约2500万个参数。

\begin{table*}[tbp]
  \centering
  \small
  \begin{tabularx}{\linewidth}{@{}llXlll@{}}
    \toprule
    阶段或层组 & 层类型 & 卷积核与步长 & 输出通道 & 重复次数 & 输出尺寸（$t\times h\times w$）\\
    \midrule
    卷积1 & 三维卷积 & $7\times7\times7$，步长2 & 论文未给出 & 1 & $32\times112\times112$\\
    池化1 & 三维最大池化 & $1\times3\times3$，步长$(1,2,2)$ & 同上 & 1 & $32\times56\times56$\\
    卷积2、3 & 三维卷积 & $1\times1\times1$；$3\times3\times3$ & 论文未给出 & 各1 & $32\times56\times56$\\
    池化2 & 三维最大池化 & $1\times3\times3$，步长$(1,2,2)$ & 同上 & 1 & $32\times28\times28$\\
    第一组模块 & 膨胀Inception模块 & 四分支，见正文 & 论文未给出 & 2 & $32\times28\times28$\\
    池化3 & 三维最大池化 & $3\times3\times3$，步长2 & 同上 & 1 & $16\times14\times14$\\
    第二组模块 & 膨胀Inception模块 & 四分支 & 论文未给出 & 5 & $16\times14\times14$\\
    池化4 & 三维最大池化 & $2\times2\times2$，步长2 & 同上 & 1 & $8\times7\times7$\\
    第三组模块 & 膨胀Inception模块 & 四分支 & 论文未给出 & 2 & $8\times7\times7$\\
    池化5 & 三维平均池化 & $2\times7\times7$，步长1 & 同上 & 1 & $7\times1\times1$\\
    分类层 & 三维卷积 & $1\times1\times1$ & $C$ & 1 & $7\times C$，在时间上取平均\\
    \bottomrule
  \end{tabularx}
  \caption{膨胀Inception-V1（I3D）的逐层结构，输入为$64\times224\times224\times3$的RGB片段。层类型、核与步长取自原文的结构图，模块数2、5、2与GoogLeNet的三组Inception模块一致；论文没有列出各层通道数，输出尺寸是按步长并假设边界补齐推算的。}
  \label{tab:vision-video-i3d-arch}
\end{table*}

训练目标是片段上的类别交叉熵，与式\eqref{eq:vision-video-tsn}中$K=1$的情形相同。学习分两步：先在Kinetics上从膨胀的ImageNet权重开始训练11万步，验证损失饱和时学习率降为十分之一；再迁移到UCF-101、HMDB-51等较小数据集，至多训练5000步，可以只训练新的分类层，也可以微调全部参数，应记录采用哪一种。优化器为动量0.9的随机梯度下降，每个卷积层后接批归一化与ReLU（输出类别得分的最后一层除外）。训练时把短边缩放到256像素后随机裁剪$224\times224$，在时间上随机选择起始帧，过短的视频循环补足，并作左右翻转；光流用TV-L1算法提取。推理时网络在时间维上以全卷积方式处理整段视频（约250帧、10秒），取中心$224\times224$裁剪，再对时间上的预测取平均。同样数量的输入帧仍可能对应不同的观察范围：若输入16帧、原帧率30、每隔2帧取一帧，首末观测间隔是$15\times2/30=1$秒；改变原帧率或步长，网络就会看到不同速度和范围的运动。

TSN与I3D面对同一个难点作出了相反的选择，比较见表\ref{tab:vision-video-activity-compare}。I3D把TSN的单帧片段换成连续片段、二维网络换成膨胀的三维网络，换来了从RGB直接学习局部运动和顺序的能力，代价是每个片段的计算量和显存大幅增加，单个片段的时间覆盖也有限；TSN则以极低的计算覆盖整段视频，却必须依赖光流分支提供运动，均值共识也不保留顺序。二者并不互斥，后续工作常把连续片段网络放进分段框架中使用。\ref{sec:vision-ivm-timesformer}节的TimeSformer与\ref{sec:vision-ivm-videomae}节的VideoMAE可以替换I3D作为片段编码器，用法相同：按预训练时的帧数与步长组织输入，在主干之后接分类层，并用视频标签微调；它们改变的是主干内部建模时空关系的方式，片段采样、标签形式与评价协议仍按本小节规定。无论用哪种主干，背景与物体的静态线索都可能支配类别判断，检验运动判断时可以比较正放、倒放及保持背景而替换动作的受控片段。

\begin{table}[htbp]
  \centering
  \small
  \begin{tabularx}{\linewidth}{@{}lXX@{}}
    \toprule
    比较项 & TSN & I3D\\
    \midrule
    时间覆盖 & 分段稀疏取样，覆盖整段 & 一个连续片段，覆盖有限\\
    运动来源 & 光流分支 & 三维卷积直接学习，光流分支可补充\\
    顺序敏感性 & 均值共识不保留段间顺序 & 片段内的时间权重可表达方向\\
    初始化 & ImageNet二维网络，跨模态初始化 & ImageNet二维网络膨胀为三维\\
    单次计算 & 低，与视频长度无关 & 高，随片段帧数增长\\
    \bottomrule
  \end{tabularx}
  \caption{两种片段活动分类模型的对照。比较固定为同一组已裁剪片段的单标签分类，评价为视频级准确率。}
  \label{tab:vision-video-activity-compare}
\end{table}

\subsection{事件区间的定位}
\label{sec:vision-video-temporal-boundaries}

未裁剪的长视频中可能有多次事件，也可能大部分时间没有任何预设动作。事件区间的定位要回答“哪类事件在何时开始、何时结束”。输入是按时间排列的片段特征序列，通常由上一小节的I3D等主干在固定步长上预先提取；输出是带类别的时间区间集合，或每个时刻的步骤标签。难点有三：同一类事件的持续时间可能相差几十倍，边界附近的画面本身含糊，逐帧判断又容易在一个长步骤中间来回跳变。

这里有两种输出形式，对应两类任务。\term{时间动作定位}（temporal action localization）在未裁剪视频中输出预设动作类别及发生区间，例如两次拿杯分别位于第9至17秒和第25至30秒；区间之间允许有空白、重叠和重复。\term{动作分割}（action segmentation）则为每个时间位置分配一个步骤类别，再把相邻的同类位置合并成片段，例如拿杯、倒水和放杯，通常要求每个位置都在规定步骤词表中作出选择。ActionFormer直接预测区间：在各时间位置判断是否有事件，并回归到两侧边界的距离；MS-TCN先逐帧预测，再由多阶段网络修正碎片化，区间由连续同类标签得到。

\subsubsection{ActionFormer}
\label{sec:vision-video-actionformer}

早期的时间动作定位方法常常分两阶段，先生成候选区间，再对候选分类；或者在每个位置放置多个预设长度的锚窗口。前者需要分别设计和训练两个阶段，后者需要事先枚举持续时间，而事件长度的跨度又很大。ActionFormer的主要想法是去掉锚窗口，做单阶段预测：用Transformer编码时间序列，构造多尺度的特征金字塔，让每个尺度的每个时间位置直接判断类别并回归到事件左右边界的距离\scite{vision-actionformer2022}。短事件由细尺度负责，长事件由粗尺度负责，于是无须枚举长度。

模型架构如下，逐层结构见表\ref{tab:vision-video-actionformer-arch}。输入是片段特征序列$\bm{Z}\in\mathbb{R}^{n\times d}$，$n$是时间位置数，每个位置对应固定的时间步长。在THUMOS14上，原工作用Kinetics预训练的双流I3D提取特征：每次输入16帧、以步长4滑动，两路各取最后全连接层之前的1024维特征，拼接成$d=2048$。投影部分是两层核宽为3的一维卷积，输出512通道并接ReLU；作者发现加入位置嵌入反而降低效果，因而默认不用，这也使推理时可以直接输入任意长度的完整序列。编码器由7个Transformer单元组成，每个单元的计算为
\[
  \begin{gathered}
  \bar{\bm{Z}}^{\ell}=\alpha^{\ell}\operatorname{MSA}\bigl(\operatorname{LN}(\bm{Z}^{\ell-1})\bigr)+\bm{Z}^{\ell-1},\\
  \widehat{\bm{Z}}^{\ell}=\bar{\alpha}^{\ell}\operatorname{MLP}\bigl(\operatorname{LN}(\bar{\bm{Z}}^{\ell})\bigr)+\bar{\bm{Z}}^{\ell},\qquad
  \bm{Z}^{\ell}=\operatorname{\downarrow}(\widehat{\bm{Z}}^{\ell})\text{，}
  \end{gathered}
\]
其中$\operatorname{LN}$是层归一化，$\operatorname{MSA}$是只在宽度为19的局部时间窗口内计算的多头自注意力，即每个位置只与邻近位置交换信息，使计算量随序列长度线性增长；MLP用GELU激活；$\alpha^{\ell},\bar{\alpha}^{\ell}$是逐通道可学习的缩放系数；$\downarrow$是步长为2的深度一维卷积，前两个单元不下采样，后五个单元各下采样2倍。第1至第6个单元的输出构成6层特征金字塔，长度从$n$依次减半到$n/32$，第$\ell$层的一个位置覆盖的时间是输入步长的$2^{\ell-1}$倍。同样宽度19的窗口在下采样16倍的层上覆盖304个输入步长，这就是局部注意力与金字塔配合的用意。多头注意力的头数论文未给出。分类头与回归头在各尺度共享参数，各由三层核宽为3的一维卷积组成，前两层为512通道并接层归一化与ReLU：分类头最后一层输出$C$个通道并接sigmoid，回归头最后一层输出到左、右边界的两个距离并接ReLU，保证非负。设某位置换算到物理时间后为$\tau$，该尺度一步对应$\Delta$秒，网络输出的距离$a,b$以尺度步为单位，则候选区间为
\begin{equation}
  [\widehat{\tau}_{\mathrm{s}},\widehat{\tau}_{\mathrm{e}}]
  =[\tau-\Delta a,\ \tau+\Delta b]\text{。}
  \label{eq:vision-video-boundaries}
\end{equation}
若输入特征曾按视频时长重新采样，还需先恢复原始的时间映射；不能把不同尺度的原始输出不加换算地相减。

例如$\tau=12$秒、左右距离已经换算为3秒和5秒，输出为$[9,17]$秒。若模型实际在每步0.5秒的网格上输出$a=6,b=10$，仍得到同一区间。图\ref{fig:vision-temporal-boundaries}用这个例子说明时间点与左右边界的关系。一个位置预测为正类，不表示只有该瞬间发生动作，它还要通过两侧距离描述持续范围。

\begin{figure}[htbp]
  \centering
  % 上下两部分是独立的教学时间范围；纵向位置不承载数值。
\begin{tikzpicture}[x=1cm,y=1cm,font=\figurefont\small]
  \node[anchor=west] at (0,3.7) {时间点预测左右边界};
  \draw[-{Stealth[length=2mm]},BookInk,line width=0.8pt]
    (1,2.45) -- (10,2.45) node[below left] {秒};
  \fill[BookTeal!14] (2,2.85) rectangle (8.4,3.2);
  \draw[BookTeal,line width=1pt] (2,2.7) -- (2,3.35);
  \draw[BookTeal,line width=1pt] (8.4,2.7) -- (8.4,3.35);
  \fill[BookInk] (4.4,3.02) circle (2pt);
  \draw[{Stealth[length=1.7mm]}-{Stealth[length=1.7mm]},BookInk,line width=0.8pt]
    (2,3.02) -- (4.4,3.02);
  \draw[{Stealth[length=1.7mm]}-{Stealth[length=1.7mm]},BookInk,line width=0.8pt]
    (4.4,3.02) -- (8.4,3.02);
  \node[above] at (3.2,3.1) {3秒};
  \node[above] at (6.4,3.1) {5秒};
  \foreach \x/\lab in {2/9,4.4/12,8.4/17}{
    \draw[BookInk,line width=0.7pt] (\x,2.4) -- (\x,2.5);
    \node[below] at (\x,2.4) {\lab};
  }
  \node[anchor=west] at (0,1.55) {逐帧标签与片段碎裂};
  \node[anchor=east] at (1.2,0.9) {真实};
  \node[anchor=east] at (1.2,0.2) {预测};
  \foreach \x/\gt/\pr in {1.5/拿/拿,2.8/倒/倒,4.1/倒/拿,5.4/倒/倒,6.7/放/放,8/放/放}{
    \draw[BookRule,fill=BookPaper,line width=0.5pt] (\x,0.62) rectangle ++(1.2,0.55);
    \node at (\x+0.6,0.9) {\gt};
    \draw[BookRule,fill=BookPaper,line width=0.5pt] (\x,-0.08) rectangle ++(1.2,0.55);
    \node at (\x+0.6,0.2) {\pr};
  }
  \draw[BookAmber,line width=1.1pt] (4.1,-0.08) rectangle (5.3,0.47);
\end{tikzpicture}

  \caption{时间点回归区间与逐帧步骤分段的教学对照。上半图中12秒位置预测左右距离，形成9至17秒区间；下半图为另一段六个等长时间格的标签，孤立的“拿”把真实连续的“倒”分成两段。两种输出需要分别检查边界与碎片化。}
  \label{fig:vision-temporal-boundaries}
\end{figure}

训练首先要决定哪些位置是正样本。每层规定可回归的距离范围，相邻层的范围大致按2倍扩大，一个真实区间只分配给其长度落在该范围内的层；在这一层上，只有落在事件中心邻域内的位置才作为正样本，称为中心采样。对中心为$c$的事件，第$\ell$层上落在以$c$为中心、半径为1.5倍该层步长的区间内的位置算作正样本。一个位置若落入多个真实区间，按训练实现的规则选择其一。中心采样使正监督集中在信息较充分的事件中心附近，它只用于训练时分配标签，推理时并不需要真实边界。

损失由两项组成。大量背景位置会让普通分类误差由负类主导，所以分类用\term{焦点损失}（focal loss）：在正确标签概率$p_{\mathrm{true}}$的负对数误差之外乘$(1-p_{\mathrm{true}})^\gamma$，减轻已经容易判对的位置的贡献，其设计与\ref{chap:vision-image-recognition}章的密集检测相同。区间回归不仅希望左右距离接近，还希望解码后的范围重叠。令$\widehat{I}$为预测区间、$I^*$为真实区间，时间交并比是交集时长除以并集时长，距离交并比误差为
\[
  1-\operatorname{tIoU}(\widehat{I},I^*)+\frac{(c(\widehat{I})-c(I^*))^2}{\ell_{\mathrm{cover}}^2}\text{，}
\]
其中$c$是区间中心，$\ell_{\mathrm{cover}}>0$是同时覆盖两个区间的最短区间长度。第一项鼓励重叠，第二项在两区间不重叠时仍能给出指向正确位置的梯度。分类误差在全部位置上累积，回归误差只在正位置上累积，再除以正样本数量；实现要为没有正样本的批次设置非零的归一化保护。两项损失的权重之比取$1:1$。每批对序列裁剪或填充并屏蔽无效位置，用Adam加预热日程更新时间编码器和两个预测头；在THUMOS14上，输入长度上限为2304个位置（约5分钟），训练50轮、前5轮线性预热，初始学习率$10^{-4}$并按余弦衰减，批大小2，权重衰减$10^{-4}$，另用参数的指数滑动平均和梯度裁剪稳定训练。输入特征是预先提取的，提供特征的视频主干不随之更新。

\begin{table*}[tbp]
  \centering
  \small
  \begin{tabularx}{\linewidth}{@{}lXXllX@{}}
    \toprule
    阶段或层组 & 层类型 & 卷积核与步长 & 输出通道 & 重复次数 & 输出尺寸\\
    \midrule
    投影 & 一维卷积，ReLU & 核3，步长1 & 512 & 2 & $n\times512$\\
    单元0 & 局部自注意力Transformer单元 & 窗口19，不下采样 & 512 & 1 & $n\times512$\\
    单元1 & 同上，金字塔第1层 & 窗口19，不下采样 & 512 & 1 & $n\times512$，范围$[0,4)$\\
    单元2至6 & 同上，金字塔第2至6层 & 窗口19，深度卷积下采样2倍 & 512 & 5 & $n/2$至$n/32$长，范围$[4,8)$至$[64,\infty)$\\
    分类头 & 一维卷积，层归一化，ReLU & 核3，步长1 & 512、512、$C$ & 3层 & 各层长度$\times C$\\
    回归头 & 一维卷积，末层ReLU & 核3，步长1 & 512、512、2 & 3层 & 各层长度$\times2$\\
    \bottomrule
  \end{tabularx}
  \caption{ActionFormer的逐层结构（论文附录表D），输入为$n\times d$的片段特征序列。回归范围在THUMOS14上以秒计，相邻层按2倍扩大；Transformer单元的头数、MLP倍率与总参数量论文未给出。}
  \label{tab:vision-video-actionformer-arch}
\end{table*}

推理时输入完整的特征序列，解码各尺度所有位置的候选区间，按得分筛选后用Soft-NMS处理重复：它逐次选出得分最高的候选，按与已选区间的重叠程度降低其余候选的得分，而不是直接删除，以减少同一次事件被重复报告。它仍可能把真实的紧邻重复动作误压低，需要按次数与边界分别检查。输出类别来自训练词表，不能响应任意文字查询；原方法利用给定序列的前后上下文，若事件正在发生，右端点尚未观测到，在线使用必须区分预测的结束位置与已经确认的结束位置。

\subsubsection{MS-TCN}
\label{sec:vision-video-mstcn}

动作分割要给长视频的每个时刻都打上步骤标签。逐帧分类器只看局部特征，在“拿杯—倒水—放杯”这样的长过程中，常会在倒水中间偶然跳回拿杯，造成过度分段；用一个时间卷积网络扩大感受范围能缓解，但一次预测仍难以消除所有孤立跳变。MS-TCN的主要想法是多阶段修正：第一阶段从视频特征给出初步的逐帧概率，后续每个阶段只读取前一阶段的概率序列，利用更大的时间上下文判断哪些跳变不合理并加以修正\scite{vision-mstcn2019}。

模型架构中，每个阶段都是单阶段时间卷积网络。第一阶段输入$\mathbb{R}^{n\times d}$的特征序列，先用$1\times1$卷积投影通道，再叠加核宽3、\term{空洞卷积}（dilated convolution）率依次为$1,2,4,\ldots$的残差时间卷积层，最后输出每个时刻$C$个类别的概率。空洞卷积在核元素之间跳过若干位置，可以在不增加参数的前提下扩大覆盖范围。第$\ell$层的空洞率$d_\ell=2^\ell$，步长为1且不下采样时，每层新增$2d_\ell$个网格的感受范围，因此$L$层覆盖的长度为
\[
  1+2\sum_{\ell=0}^{L-1}2^\ell=2^{L+1}-1\text{。}
\]
每层内部的计算顺序是：空洞卷积，ReLU，$1\times1$卷积，加残差，即$\widehat{\bm{H}}^{(l)}=\operatorname{ReLU}(\bm{W}_1*\bm{H}^{(l-1)}+\bm{b}_1)$、$\bm{H}^{(l)}=\bm{H}^{(l-1)}+\bm{W}_2*\widehat{\bm{H}}^{(l)}+\bm{b}_2$，其中$*$表示一维卷积，$\bm{W}_1$是核宽3的空洞卷积，$\bm{W}_2$是$1\times1$卷积，每层之后还使用dropout。原工作每阶段用10层，空洞率从1加倍到512，覆盖2047个网格；这是网格长度，还要乘特征的时间步长才换算为秒，在每秒15帧的特征上约为136秒。所有层的通道数都是64。后续阶段结构相同，但输入换成$\mathbb{R}^{n\times C}$的概率序列，原工作共用4个阶段，结构见表\ref{tab:vision-video-mstcn-arch}。

\begin{table}[htbp]
  \centering
  \small
  \begin{tabularx}{\linewidth}{@{}lXXlll@{}}
    \toprule
    层组 & 层类型 & 核与空洞率 & 输出通道 & 重复次数 & 输出尺寸\\
    \midrule
    输入投影 & $1\times1$卷积 & 核1 & 64 & 1 & $n\times64$\\
    残差层 & 空洞卷积、ReLU、$1\times1$卷积、残差、dropout & 核3，空洞率$1,2,\ldots,512$ & 64 & 10 & $n\times64$\\
    输出层 & $1\times1$卷积、softmax & 核1 & $C$ & 1 & $n\times C$\\
    阶段重复 & 后三个阶段以前一阶段的概率为输入 & 同上 & 同上 & 4个阶段 & $n\times C$\\
    \bottomrule
  \end{tabularx}
  \caption{MS-TCN的结构。第一阶段的输入为$n\times d$的I3D特征，后续阶段的输入为$n\times C$的概率序列；dropout比例论文未给出。时间维步长始终为1，不下采样。}
  \label{tab:vision-video-mstcn-arch}
\end{table}后续阶段不读原始特征，所以多阶段并不增加新的视频证据，只是学习修正前一阶段已有的判断。卷积是非因果的，每个位置都会读取未来时刻的特征，因此不能直接用于不允许等待的在线输入。

训练在每个阶段都施加两项损失。分类项是逐帧交叉熵
\[
 L_{\mathrm{cls}}=-\frac{1}{n}\sum_t\log p_{t,y_t}\text{；}
\]
平滑项惩罚相邻时刻对数概率的跳变，并对过大的跳变截断：
\begin{equation}
  L_{\mathrm{smooth}}
  =\frac{1}{nC}\sum_{t=2}^{n}\sum_{c=1}^{C}
  \min\!\left\{
  \bigl(\log p_{t,c}-\log p_{t-1,c}\bigr)^2,\ \eta^2
  \right\}\text{，}
  \label{eq:vision-video-smoothing}
\end{equation}
这里$n$为序列长度，$C$为类别数，$\eta>0$为截断上限。截断的作用是：真实边界处概率本来就应该大幅变化，截断后这些大跳变的惩罚不再继续增长，平滑项主要作用于小幅抖动。计算时把前一时刻的概率当作固定目标，只对$p_{t,c}$一侧求梯度。每阶段目标为$L_{\mathrm{cls}}+\lambda L_{\mathrm{smooth}}$，原工作取$\lambda=0.15$、$\eta=4$，全部阶段求和后在视频批次上平均。输入的是预先提取的I3D特征，时间分辨率为每秒15帧；时间卷积从随机初始化开始，用学习率0.0005的Adam更新。

在一次教学跳变中，“倒水”概率从0.8降到0.1，对数差绝对值为$\log 8\approx2.079$，未截断的平方约为4.324。若示例设$\eta=2$，这一对只贡献4；从0.8到0.7的变化则仅贡献约0.0178。推理时按最后一阶段的概率逐时刻取类别，再合并连续同类标签。图\ref{fig:vision-temporal-boundaries}下半的真实标签为“拿、倒、倒、倒、放、放”，预测为“拿、倒、拿、倒、放、放”，六格中只错一格，却多出一个虚假的拿杯段并把倒水切成两段。只报告帧准确率$5/6$会掩盖这种步骤结构错误；反过来，过强的平滑也可能删掉真实的短步骤。

ActionFormer与MS-TCN都从预提取的片段特征出发，比较见表\ref{tab:vision-video-event-compare}。两者的差别在于把“事件”表示成什么：ActionFormer表示成带类别的区间集合，每个区间由一个位置和两段距离给出，允许空白、重叠和重复，适合稀疏分布的动作；MS-TCN表示成逐帧标签序列，区间是合并的结果，适合首尾相接的步骤序列，但不能表达同一时刻的两个事件。ActionFormer依靠多尺度和回归处理时长差异，代价是需要中心采样等标签分配规则以及Soft-NMS后处理；MS-TCN依靠多阶段修正和平滑损失处理碎片化，代价是全序列非因果计算和可能抹掉短步骤。

\begin{table}[htbp]
  \centering
  \small
  \begin{tabularx}{\linewidth}{@{}lXX@{}}
    \toprule
    比较项 & ActionFormer & MS-TCN\\
    \midrule
    输出形式 & 带类别的区间集合 & 逐帧步骤标签，合并为片段\\
    处理时长差异 & 多尺度金字塔与边界回归 & 空洞卷积扩大感受范围\\
    处理碎片化 & Soft-NMS抑制重复区间 & 多阶段修正与截断平滑损失\\
    允许空白与重叠 & 允许 & 每个时刻一个类别，不允许重叠\\
    主要评价 & 多个tIoU门限下的平均精度 & 帧准确率、片段F1与编辑距离\\
    \bottomrule
  \end{tabularx}
  \caption{两种事件区间定位模型的对照。比较固定为同一组预提取的片段特征；评价指标的定义见\ref{sec:vision-video-output-protocol}节。}
  \label{tab:vision-video-event-compare}
\end{table}

\subsection{未完成事件的预判}
\label{sec:vision-video-proactive-response}

前两小节都假定事件已经结束，模型可以读到它的全部画面。人正靠近热水壶时，拿起壶的画面尚未到达。完整录像的分类可以在动作结束后回看，在线助手却可能需要在手伸向危险位置时就提醒。输入因此变成截至当前时刻的帧，输出是在某个时刻发出（或暂不发出）的判断。应先约定触发事件、允许的响应区间和误报代价，再区分已经发生但证据不足、尚未完成和当前无需回答这几种状态。预测之后可能发生什么，与确认已经看见了什么，是两种不同的输出；\ref{sec:vision-video-dense-state}节V-JEPA 2.1的动作预期头属于前者，它只读取动作开始前的画面。

SPOT-Bench把提问和后续视频保持为连续的输入流，评价主动查询、无请求的干预以及多个并存查询的响应，而不只在人工暂停的几个时刻问答\scite{vision-spot-timeliness2026}。它的Timeliness-F1同时检查语义、时机和整条时间线上的输出频率。设事件的可接受响应区间为$[t_s,t_e]$，输出发生于$\tau$。为了在区间内给足额信用、同时惩罚过早与过晚，论文定义
\begin{equation}
 \mathcal{T}(\tau)=
 \begin{cases}
  \exp\!\bigl(-(\tau-t_s)^2/(2\sigma_{\mathrm{early}}^2)\bigr),
    &\tau<t_s,\\
  1,&t_s\leq\tau\leq t_e,\\
  \exp\!\bigl(-(\tau-t_e)^2/(2\sigma_{\mathrm{late}}^2)\bigr),
    &\tau>t_e\text{，}
 \end{cases}
 \label{eq:vision-video-timeliness}
\end{equation}
其中两个$\sigma$以秒为单位，原协议令$\sigma_{\mathrm{early}}>\sigma_{\mathrm{late}}$，即同样的偏离长度下更严厉地惩罚迟到。区间由任务和人工时机标注确定，不能看到模型何时回答再把区间移过去；内容错误的及时提醒也不算正确匹配。

先考虑每个事件只允许一个有效响应的情形。按时间处理所有输出，与尚未占用且语义相符的事件槽匹配，重复或无效的输出计误报，未填的槽计漏报。设匹配次数为$N_{\mathrm{TP}}$，时间加权信用总和为$x_{\mathrm{TP}}=\sum_{\mathrm{matched}}\mathcal{T}(\tau)$，则精确率与召回率分别为$x_{\mathrm{TP}}/(N_{\mathrm{TP}}+N_{\mathrm{FP}})$和$x_{\mathrm{TP}}/(N_{\mathrm{TP}}+N_{\mathrm{FN}})$，再取调和均值得到Timeliness-F1。分母是整数的响应数与事件数，分子才是时间加权的信用。允许每槽多个候选时，需遵守论文附录的占用和去重规则，比较时固定这一设置。教学事件区间为$[5,6]$秒，设提前尺度1秒、迟到尺度0.5秒。6.5秒的正确输出得到$\exp(-0.5)\approx0.607$；若7秒又重复提醒，已经填满的槽不能再得信用，计一次误报。此时精确率为$0.607/2$，召回率为0.607，Timeliness-F1约为0.404。完全静默会漏掉该槽，一直重复则降低精确率。

论文同时给出AsynKV作为基线，说明持续响应的记忆怎样组织。它使用已有的Qwen2.5-VL权重，不添加新的训练目标。每个查询保存编号、提出时刻、是否已回应以及最近带时刻的答案，后续上下文包含这份答案记录，使模型知道刚才已经响应过。近期帧的键值缓存留在GPU，较早帧移到CPU形成有界历史；收到查询后按相似度取回相关帧的缓存，放到近期帧之前并重新安排时间位置。这里的缓存是语言模型已经算过的内部键值，不是原始视频，也不能从中恢复未保留的像素。没有待回应查询的时段用于摘要、缓存搬运和压缩，原方法寻找相邻最相似的两帧并丢弃较旧的一项来控制长期容量。原配置每秒读一帧、GPU近期16帧、CPU长期至多256帧、每查询取回4帧；这是该实验的协议，不是视频助手的一般最优设置。更高帧率或繁忙查询下，后台任务仍可能与响应争用资源，应实际记录硬件、排队、漏帧和响应耗时。对“靠近—伸手—拿起”的各时刻还应保存当时可见的帧，6.5秒的修订不能回填为5秒时已经正确预警。

\subsection{事件顺序与重复的解释}
\label{sec:vision-video-event-relations}

得到事件区间之后，可以把每个结果保存为一条记录：对象身份、类别、开始时间、结束时间和来源位置。两个事件是否相邻、重叠或重复，应由这些记录计算，而不是从视频的总标签推测。“拿起”与“放下”还要结合杯子相对桌面的状态变化，即\ref{sec:vision-video-continuous-attributes}节的状态序列；车辆转弯则要结合\ref{sec:vision-video-tracking}节的轨迹方向和道路参照。

记录的质量直接决定解释的结果。两次拿杯若因区间过宽而被合并，次数会少算；一次倒水若被分成三个小段，次数会多算。例如真实拿杯为$[9,17]$和$[19,24]$秒，若预测的单个区间为$[9,24]$秒，它与第一次的时间交并比只有$8/15\approx0.53$，计数也从2变为1。前一段结束恰好早于后一段开始，只提供时间先后，并不证明前者导致后者。更复杂的跨事件语言解释将在\ref{chap:vision-multimodal-understanding}章中连接到答案证据的核验。

\section{视频匹配与片段检索}
\label{sec:vision-video-retrieval}

前面各节从视频出发，发现其中的对象与事件；检索则从查询出发，决定哪些内容值得返回。输入是一个查询（一张车辆照片、一段示例视频或一句文字）和一个视频库，输出是按相关性排序的视频、对象轨迹或时间片段。难点有三：查询与视频之间的相关性需要先定义清楚；视频库可能很大，必须先用低成本的表示快速召回候选；长视频的关键证据可能只占几秒，粗粒度的表示会把它稀释掉。

检索按流程分为三步。先规定匹配要求，即什么样的结果算作正确、返回单位是什么；再召回候选片段，这一步决定用什么表示与查询比较；最后核验召回结果，检查对象、动作和时间范围是否真正满足查询。前面得到的帧、轨迹、片段和事件表示都可以作为召回的索引，但原始画面与时间戳必须保留，供核验回查。

\subsection{匹配要求的规定}
\label{sec:vision-video-query}

视频查视频可能寻找相同事件或相似运动，图像查视频可能寻找相似外观，给定车辆照片则可能要求同一实例。即使三种查询都用向量相似度排序，正例的定义也不同。任务协议还要规定返回单位是整段视频、身份轨迹还是时间区间，以及边界可以偏差多少。

文字查询“有人把杯子放到桌上”增加了语言与视频的对应。图文表示的对齐见\ref{sec:vision-models-contrastive}节；这里先把文字编码器视为能够输出可比较向量的接口，重点追踪它与哪些帧、哪些范围比较。精确的语言指代及查询区间定位将在\ref{chap:vision-multimodal-understanding}章展开。

\subsection{候选片段的召回}
\label{sec:vision-video-candidates}

召回的目标是在可接受的成本下，把真正相关的片段尽量留在候选集中。核心的设计选择是用什么表示与查询比较，以及在什么粒度上建立索引。每段视频一个向量便于快速搜索，但位置和时间被压缩掉后，很难说清证据来自哪里；局部索引保存更多表示和时间信息，成本更高，却能把命中回溯到实际观测。

三种路线对这一选择给出了不同回答。基于视觉表示的召回直接用前文得到的帧、轨迹或片段特征与图像或视频查询比较，不需要新的训练，但无法直接处理文字查询。基于帧聚合文字匹配的召回以CLIP4Clip为代表，把图文对齐模型逐帧编码的向量聚合成视频表示，再用视频文本配对训练，使文字查询能与视频直接比较。基于分层索引的长历史召回以TRACE为代表，面对单个长视频中的复杂问题，不再依赖一个预先算好的向量，而是先看粗略概览，再按问题回查原始短片段，逐轮积累证据。三者依次增加了对查询语义和时间细节的处理能力，也依次增加了每次查询的计算量。

\subsubsection{基于视觉表示的召回}
\label{sec:vision-video-visual-matching}

图像查询可以使用\ref{sec:vision-recognition-retrieval}节的描述向量，与逐帧全图或对象裁剪比较；活动查询可以与\ref{sec:vision-video-activity-evidence}节得到的连续片段表示比较。索引条目至少保留视频编号、时间范围及必要的区域坐标，表示来源还应注明训练与归一化方式。建立索引只是编码和存储，没有重新学习表示。

给定查询向量$\bm{q}$和候选帧向量$\bm{v}_t$，帧级最大相似度$\max_t\langle\bm{q},\bm{v}_t\rangle$强调曾经出现过的强匹配；沿轨迹取均值则积累同一对象的外观证据。前者可能被一次错误匹配抬高，后者依赖身份关联。对运动事件，独立的帧得分无法区分同样画面元素的不同顺序，需要使用保存局部时间关系的片段表示。

相似度计算的候选集合来自实际建立的索引，不能搜索尚未编码或已经删除的位置。若采用余弦比较，先把非零向量除以各自的$L_2$范数；零向量应排除或按协议标为未知。教学车辆查询$\bm{q}=(1,0)$，两条轨迹的区域表示分别为$(0.8,0.6)$与$(0.6,0.8)$，二者已是单位向量，得分为0.8与0.6，前者先返回。若整帧背景表示反而让第二条得分更高，区域回查就能查出两种排序依据的差异。以汽车图片查录像时，整帧可能因为同一片道路和天空而命中，裁剪区域却显示车型不同。若索引窗口为10秒，排序只说明这个10秒窗口值得检查，不能认定动作恰好持续10秒。

\subsubsection{基于帧聚合文字匹配的召回}
\label{sec:vision-video-clip4clip}

车辆照片可以直接匹配车辆区域，文字查询却需要把“把杯子放到桌上”这样的描述与视频画面连接起来。CLIP4Clip之前的文本视频检索方法多数使用若干专家模型离线提取的视频特征，特征提取器不随检索任务更新；专门在视频文本数据上预训练又代价很高。CLIP在大规模图文对上学到的对齐已经能把文字与单幅图像比较，问题是怎样把它迁移到由多帧组成的视频上。CLIP4Clip的主要想法是：直接用CLIP的图像编码器逐帧编码，把帧向量聚合成视频向量，再在视频文本配对上端到端微调，并比较不同的聚合方式\scite{vision-clip4clip2021}。

模型架构由三部分组成。视频端按每秒1帧的速率均匀取帧，原工作在MSR-VTT上至多取12帧，每帧缩放裁剪到$224\times224$，经CLIP的ViT-B/32图像编码器得到帧向量$\bm{f}_t\in\mathbb{R}^{512}$；文字端截断为至多32个词元，经CLIP的文字编码器得到$\bm{w}\in\mathbb{R}^{512}$。中间的相似度计算器有三类。无参数型对帧向量求均值：
\begin{equation}
  \bm{v}=\operatorname{norm}\Bigl(\frac{1}{n}\sum_{t=1}^{n}\bm{f}_t\Bigr)\text{，}
  \label{eq:vision-video-clip4clip-mean}
\end{equation}
其中$n$为取帧数，$\operatorname{norm}$把非零向量除以其$L_2$范数，再与归一化的文字向量求余弦相似度；时间索引只用于取帧，不参与均值运算，所以交换帧序结果不变。序列型用LSTM（长短期记忆网络）处理帧序列，或先给帧向量加上时间位置嵌入、再用4层Transformer编码器处理，然后对输出序列取均值，使时间关系可以写入视频表示；Transformer在时间维上的信息交换可参考\BookChapterRef{chap:transformer}{卷五的Transformer章}。紧耦合型则把文字词元与帧向量拼成一个序列送入Transformer，让两种模态在编码中直接交互，最后由线性层输出一个匹配分数。前两类是双编码器，视频表示可以预先算好；紧耦合型必须对每个查询与每个视频重新联合编码，无法预先建立视频索引。

各部分的配置见表\ref{tab:vision-video-clip4clip-config}，其中两个编码器沿用CLIP的结构\scite{radford2021}。图像编码器把$224\times224$的帧切成$32\times32$的图像块，得到$7\times7=49$个块词元，再加一个类别词元；每个Transformer块的计算为$\bm{z}'=\bm{z}+\operatorname{MSA}(\operatorname{LN}(\bm{z}))$、$\bm{z}''=\bm{z}'+\operatorname{MLP}(\operatorname{LN}(\bm{z}'))$，12层之后取类别词元的输出，经线性投影得到512维帧向量。文字编码器取最高层在句末标记处的激活作为句子表示。原文还试验了“三维线性”的图像块嵌入：用核为$3\times32\times32$、时间维步长与填充均为1的三维卷积代替二维图像块投影，使相邻帧在嵌入时就交换信息，初始化时把CLIP的二维投影放在时间核的中间一片、其余两片置零。序列型与紧耦合型中的位置嵌入由CLIP文字编码器的位置嵌入重复得到，Transformer编码器用CLIP图像编码器对应层的权重初始化，LSTM和线性层随机初始化。

\begin{table*}[tbp]
  \centering
  \small
  \begin{tabularx}{\linewidth}{@{}lXXlllX@{}}
    \toprule
    模块 & 图像块大小 & 词元数 & 层数 & 宽度 & 头数 & 说明\\
    \midrule
    图像编码器（ViT-B/32） & $32\times32$ & 每帧$49+1$ & 12 & 768 & 12 & 输出投影到512维；MLP倍率与参数量论文未给出\\
    文字编码器 & — & 至多32 & 12 & 512 & 8 & 句末标记处的输出作为句子表示\\
    序列型：LSTM & — & 帧数（12） & 1 & 论文未给出 & — & 输出序列再取均值\\
    序列型：Transformer & — & 帧数（12） & 4 & 论文未给出 & 论文未给出 & 加时间位置嵌入\\
    紧耦合型：Transformer & — & 1个文字向量加帧数 & 4 & 论文未给出 & 论文未给出 & 首个输出经两层线性层与ReLU得到匹配分数\\
    \bottomrule
  \end{tabularx}
  \caption{CLIP4Clip各模块的配置。图像与文字编码器的层数、宽度和头数取自CLIP论文的配置表；MSR-VTT实验取12帧、32个文字词元，ActivityNet与DiDeMo作段落检索时两者都取64。}
  \label{tab:vision-video-clip4clip-config}
\end{table*}

训练数据是视频及其文字描述，不需要动作边界。设一个批次含$B$对匹配的视频和文字，$\bm{v}_i,\bm{w}_j$为归一化后的表示，$z_{i,j}=\langle\bm{v}_i,\bm{w}_j\rangle/\tau$为相似度，$\tau>0$是温度。两个方向分别在同一矩阵的行和列中识别正确配对：
\begin{equation}
  \begin{aligned}
    L_{\mathrm{v2t}}&=-\frac{1}{B}\sum_i
      \log\frac{\exp z_{i,i}}{\sum_j\exp z_{i,j}},\\
    L_{\mathrm{t2v}}&=-\frac{1}{B}\sum_i
      \log\frac{\exp z_{i,i}}{\sum_j\exp z_{j,i}}\text{，}
  \end{aligned}
  \label{eq:vision-video-retrieval-loss}
\end{equation}
总损失为二者之和。$L_{\mathrm{v2t}}$让每个视频在批内所有文字中认出自己的描述，$L_{\mathrm{t2v}}$让每句文字在批内所有视频中认出自己的视频；批内其他配对充当负候选。如果两段视频实际上都符合同一描述，就会产生假负例，标注的一一对应并不意味着语义上的唯一匹配。以$B=2$的教学分数矩阵$\bm{Z}=\bigl(\begin{smallmatrix}2&0\\1&3\end{smallmatrix}\bigr)$为例，两行正确配对的概率都为$\exp2/(\exp2+1)\approx0.881$，所以$L_{\mathrm{v2t}}\approx0.127$；第一列正确概率为$\exp2/(\exp2+\exp1)\approx0.731$，第二列为$\exp3/(1+\exp3)\approx0.953$，故$L_{\mathrm{t2v}}\approx(0.313+0.049)/2=0.181$。两个方向的差别来自各自要区分的候选，不是把同一误差重复计算两次。

学习方法上，两个编码器从CLIP权重初始化，序列型和紧耦合型新增的模块从随机初始化开始。由于CLIP参数已经过大规模训练，而视频文本数据较少，原工作对CLIP参数使用很小的学习率（$10^{-7}$），对新增模块使用较大的学习率（$10^{-4}$），用Adam端到端更新，学习率按余弦日程衰减，批大小128，训练5轮。原文还尝试在HowTo100M的约38万段视频上用多示例对比损失对CLIP作视频文本的二次预训练，作为初步探索。评价报告文字查视频的R@1、R@5、R@10与中位排名。原论文的观察是：在其实验的数据规模上，无参数的均值聚合已经与序列型相当，紧耦合型在较小数据集上明显较差，作者把后者归因于新增的交互模块缺少足够数据训练。均值配置没有新增参数，但仍然微调编码器，“无参数聚合”不等于没有训练。

推理时预先编码视频库，收到查询后只需编码文字、计算相似度并排序。教学查询“有人把杯子放到桌上”对三个视频的分数设为0.82、0.70和0.21，依次返回第一、第二、第三个候选。第二个视频若是拿起杯子，较高的相似度可能来自相同的对象；若第一段长一分钟而放杯只占两秒，全视频均值又可能稀释这次事件。前者是顺序表达不足，后者是时间粒度不足。图库编码与查询必须使用相同的预处理和权重；更换编码器权重后，旧向量不能不加核验地继续与新查询比较。这一流程排序的是事先给定的视频或片段，没有回归动作边界；若先切成窗口再检索，窗口长度与重叠就是外部条件，不能把窗口端点解释成模型推断的动作起止。

\subsubsection{基于分层索引的长历史召回}
\label{sec:vision-video-history-retrieval}

查询“两次敲门都发生在开门之前吗”需要三段事件证据，找到第一次敲门还不能决定答案。在长视频中均匀抽帧可能漏掉短暂的敲门，全视频均值也会稀释它；如果历史压缩后只剩“有人在门口”，后来就不能判断敲了几次。CLIP4Clip式的单个视频向量回答不了这类问题，需要先用粗略概览确定候选范围，再在局部增加帧密度和分辨率，且每次都回到原始画面。

TRACE按这个思路组织长视频问答中的证据收集\scite{vision-trace-anchors2026}。它的主要想法是以原始短片段作为视觉锚点，逐轮添加锚点并检查答案是否随证据稳定下来。锚点包含时间范围、抽帧方式和空间分辨率，其内容是能直接查看的原始帧，而不是上一轮生成的文字摘要。方法不训练新的参数，而是用已有的视觉语言模型执行候选提议、证据判断与回查控制；语言接口沿用\ref{chap:vision-multimodal-models}章，这里解释它实际取得了哪些视频观测。TRACE自身没有网络层，它的“结构”是表\ref{tab:vision-video-trace-components}中的一组调用环节；主实验以Gemini-2.5-Pro为骨干，该模型的规模报告未公开，论文还在共享骨干的条件下，把它与其他智能体方法和直接均匀抽帧的做法放在相同的帧预算下比较。

\begin{table*}[tbp]
  \centering
  \small
  \begin{tabularx}{\linewidth}{@{}lXX@{}}
    \toprule
    环节 & 输入 & 输出\\
    \midrule
    概览（粗层索引） & 完整视频，均匀取32帧 & 带原始时刻的32帧概览\\
    时间提议 & 概览、问题、选项、已有锚点、未探索区间、未满足的证据需要 & 候选时间范围及观察方式（帧率、分辨率）\\
    锚点提取（细层索引） & 候选范围内的原始视频 & 与问题相关的短片段，即视觉锚点；无关则不保留\\
    锚点排序 & 锚点、问题、选项 & 按区分选项或填补缺口的作用排定的锚点顺序\\
    证据判断 & 某一前缀的原始锚点、问题、选项，不继承上一前缀的输出 & 答案候选、可答或冲突状态、已观察的区分事实、剩余缺口\\
    轨迹协调 & 全部前缀的判断记录 & 停止、删除、重读或扩展四种控制动作\\
    最终复查 & 最早稳定前缀的原始片段、问题、选项 & 只含答案的单次回答，与稳定前缀一致才接受停止\\
    \bottomrule
  \end{tabularx}
  \caption{TRACE的调用环节。各环节都由同一个已有的视觉语言模型执行，方法不引入可训练参数；概览帧和每次判断重新输入的帧都计入实际的视觉输入成本。}
  \label{tab:vision-video-trace-components}
\end{table*}

流程如下。初始从完整视频均匀取32帧形成概览，记录其原始时刻，这是分层索引的粗层。时间提议模块结合概览、问题、选项、已经查看的锚点、未探索的区间和尚未满足的证据需要，提出候选时间范围及观察方式；局部提取按该方式读取原视频，保留与问题有关的短锚点，这是细层。优先级模块再按锚点区分答案或填补缺口的作用排定顺序。这个顺序是证据加入的顺序，不是事件发生的顺序，事件本身仍按原始时间戳排列。设排好序的锚点为$a_{(1)},\ldots,a_{(n)}$，为了查明新增画面是否改变回答，形成前缀证据集合
\begin{equation}
  \mathcal{E}_j=\{a_{(1)},\ldots,a_{(j)}\},
  \qquad
  r_j=\operatorname{EA}(\mathcal{E}_j,q,O)\text{，}
  \label{eq:vision-video-trace-prefix}
\end{equation}
其中$q$是问题，$O$是答案选项，$\operatorname{EA}$是一次证据判断调用，输出答案候选、可回答或冲突状态、已观察到的区分事实及剩余缺口。每次调用只接收当前的原始证据、同一问题和选项，不继承上一前缀的答案或文字推理；这样比较$r_j$和$r_{j+1}$时，输入的差别只是新增的锚点，而不是让模型顺着自己的旧答案继续解释。

独立前缀的代价是反复输入相同的画面。若三个锚点分别含4、3、3帧，三次判断依次输入4、7、10帧，共21帧，而不是只支付10帧。初始概览、局部筛选中给模型看的画面以及最终复查也都计入实际的视觉输入成本；用于覆盖检查的不同时间戳数与用于成本检查的重复输入帧数是两种量。轨迹协调模块读取这串判断，寻找最早的稳定前缀：后来的锚点不改变答案、不降低可回答性、不与已记录事实冲突，也没有未满足的证据需要。若画面冗余或错误，就删除锚点并重建判断；若候选中看不清动作，就以更高密度或分辨率重读；若仍缺事件，则扩展未探索范围。

在教学视频中，第一个锚点显示10秒敲门，第二个显示40秒开门，此时只能确认一次敲门在开门之前，第二次敲门仍是缺口。扩大搜索后，第三个锚点显示25秒另一次敲门，三者的时间顺序才支持“两次都在开门前”。即使答案编号没有变化，新锚点补上第二次敲门也说明早先的证据并不充分，不能仅凭连续两次回答相同就停止。停止前再独立调用一次，只给最早稳定前缀的原始片段、问题及选项，答案一致才报告找到稳定前缀；若预算先耗尽，只能给出带“未稳定”状态的强制答案，不能报告覆盖充分。稳定是模型在已取得证据上的停止判断，不是录像中绝无第三次敲门的证明。此流程假设离线完整原视频可以回查；在线时只能在已到达且实际保留的历史范围内执行。若要从成本与错误代价共同推出停止规则，需要比较立即作答的条件风险与继续查询后的期望风险及成本，并保留剩余预算；这一决策比较可参考\BookChapterRef{chap:decision-theory-and-dynamic-risk}{卷二“策略的价值与优化”章}。TRACE的稳定前缀规则本身没有给出这种最优性保证。

三种召回路线的对照见表\ref{tab:vision-video-retrieval-compare}。CLIP4Clip在视觉表示召回的基础上，用图文对齐的帧编码和视频文本配对训练，使文字查询可以与视频比较，代价是训练和对细节、顺序的压缩；TRACE不再预先压缩成一个向量，而是对每个查询按需回查原始片段，换来了对长视频中短暂事件和多事件顺序的处理能力，代价是每次查询都要多次调用视觉语言模型、重复输入画面，并且只适用于原视频仍可读取的情形。

\begin{table}[htbp]
  \centering
  \small
  \begin{tabularx}{\linewidth}{@{}lXXX@{}}
    \toprule
    比较项 & 基于视觉表示 & 基于帧聚合文字匹配 & 基于分层索引的长历史\\
    \midrule
    查询类型 & 图像或视频 & 文字 & 关于单个长视频的问题与选项\\
    索引单位 & 帧、轨迹或片段向量 & 每段视频一个向量 & 粗概览加按需提取的原始锚点\\
    训练需求 & 无新训练 & 视频文本配对微调 & 无新训练，调用已有模型\\
    时间细节 & 取决于索引粒度 & 被均值压缩 & 局部加密回查\\
    每次查询成本 & 向量比较 & 文字编码加向量比较 & 多次视觉语言模型调用\\
    \bottomrule
  \end{tabularx}
  \caption{三种候选片段召回路线的对照。前两种在视频库中排序候选，第三种在单个长视频中收集证据；比较时须固定索引内容与实际送入模型的帧。}
  \label{tab:vision-video-retrieval-compare}
\end{table}

\subsection{检索结果的核验}
\label{sec:vision-video-retrieval-check}

召回只说明候选值得检查。视频级检索要检查相关性和排序，实例查询还要核验对象是否同一，区间查询则要核验起止与重复次数。比如正确的视频中有两次拿杯，只返回覆盖两次的长片段，可能满足“找一段含拿杯的视频”，却不满足“找到第二次拿杯的时间”。查询语义改变了成功的条件。

精细核验可以沿候选范围重新运行检测、关联和状态判断，即把\ref{sec:vision-video-tracking}节和\ref{sec:vision-video-activities}节的方法用在少量候选上，再结合查询要求选择结果。核验依赖的知识也要区分：预设活动检测依赖固定类别的标注，任意文字定位还依赖语言对齐与指代监督；前者不能仅因输出了区间，就被视为已经解决了后者。

\section{时序判断与历史依据的检验}
\label{sec:vision-video-evaluation}

一次交通分析可以给对车型却交换了身份，检索命中正确的视频却圈错了转弯区间。视频系统的输出由多个环节串成，单一的总分无法说明哪个环节出了问题。本节从三个方面检验时序判断：输出本身是否在共同的观察协议下被正确度量；错误沿着检测、关联、事件和检索的链条从哪里进入；使用记忆的系统，其收益是否真的来自正确的历史内容。

\subsection{输出与观察协议的核验}
\label{sec:vision-video-output-protocol}

比较之前要先固定观察协议。视频划分应尽量以源视频、对象或场景为单位，避免把相邻帧分到训练和测试中形成近重复泄漏。比较还要固定抽帧、分辨率、初始化、未来帧的可见范围和人工纠正的预算。离线算法能读完整的事件，在线算法可能尚未看到结尾，两者的差异应先从输入解释，再比较输出。

度量的定义已在\ref{sec:vision-video-recognition-challenges}节给出，这里关心的是怎样用它们得出可靠的结论。第一，同一输出要同时报告互补的指标。十帧拆成两个编号的算例中，MOTA为0.9、IDF1为0.6、HOTA约为0.72：只报MOTA会掩盖身份断裂，只报IDF1又看不出检测本身是否准确；实际还应报告身份切换次数、轨迹中断以及不同遮挡长度下的表现。同理，视频目标分割要同时看$\mathcal{J}$与$\mathcal{F}$，事件区间要报告多个tIoU门限下的mAP而不是单一门限，逐帧步骤要同时报告帧准确率与编辑分数或F1，避免长步骤的正确帧掩盖短步骤被删或反复切碎。第二，指标的计算细节要随结果一起给出：区间检测的平均精度要求类别一致、按得分排序并处理同一真实事件的重复命中；检索的Recall@$K$和中位排名以约定的相关候选为依据，同一描述对应多段视频时，单一正例的约定会低估实际效果。第三，在线输出还要检查时机，未完成事件的响应用\ref{sec:vision-video-proactive-response}节的Timeliness-F1同时检查内容与时机，并保存每次输出当时可见的帧。

\subsection{错误传播链的核验}
\label{sec:vision-video-error-trace}

交通应用会把本章的各个环节串起来：检测提供汽车范围，关联维持车辆编号，轨迹与道路方向共同支持转弯判断，检索返回需要复核的片段。错误也沿着同一路径传播。天气导致检测置信降低，可能造成漏配；低帧率使运动预测失准，可能引发身份交换；过宽的事件区间又可能把两次转弯算成一次。用一个总分类分数难以识别该修哪一环。

对一条错误记录，可以沿链条逐环回看：目标是否确实被采到，候选框是否正确，关联代价和分配结果如何，掩码写入记忆的内容是什么，事件的左右边界落在哪里。若车辆框始终正确而编号中途交换，增加分类器训练并不能直接修复关联；若身份不变但转弯区间被拉长，应检查时间特征和边界恢复。在线记录还要保留首次输出、每次修订及其当时可见的帧，避免把后来借助结尾画面修正的答案当作早期成绩。日志应区分当前观测、历史读出和外推结果，以便确定错误信息从哪里进入。

\subsection{历史依赖的核验}
\label{sec:vision-video-memory-specificity}

汽车遮挡后，使用记忆的分割得分高于不用记忆，仍不能证明收益来自这辆车的正确历史。改变记忆还可能同时改变特征尺度、缓存长度或位置编码；即使放入不相关但形状合适的表示，网络也可能获得类似的改善。因此需要固定读取时的其余计算，只替换实际被消费的值，再看哪些收益随内容的丢失而消失。

\term{记忆特异性}（memory specificity）描述记忆收益对所取内容的依赖程度，而不只是记忆通路是否存在。相关的读取替换审计在已有冻结视频模型、WorldMem和SAM 2上实施了不同的对照\scite{vision-memory-specificity2026}。它检查特定读取位置的直接影响，无须重新训练，但要求能访问模型实际消费的记忆接口；移除整个缓存不是同一项干预。这里固定的是计算模型，替换的是内部读取值，得到的结论针对该接口对输出的影响；若要解释现实事件的因果效应，还需建立现实变量、生成机制及其识别条件，可参考\BookChapterRef{chap:causal-inference}{卷二“因果响应”章}。

设$S(\bm{v})$为消费记忆值$\bm{v}$时的任务分数，统一令分数越大越好。正确值为$\bm{v}_{\mathrm{corr}}$，去除内容后的参考分数为$S_0$；模型支持无记忆模式时用该模式，否则保留原接口而把值置零。为了把不同基线上的变化分开，定义
\begin{equation}
  B=S(\bm{v}_{\mathrm{corr}})-S_0,\qquad
  R(\bm{v})=\frac{S(\bm{v})-S_0}{B}\quad(B>0)\text{。}
  \label{eq:vision-video-memory-recovery}
\end{equation}
$B$是正确记忆带来的收益，$R$是替换后仍保留的收益比例。只在$B>0$时解释这个比例；$R$也不裁剪到$[0,1]$，因为替换后的分数可以低于参考，也可能略高于正确记忆。

在目标读取步骤，查询、已选择的位置、条目数量、时间及位置编码、模型参数和此前的运行历史全部固定，只把消费的值换成有效但错误的记忆。同一轨迹中错误时刻的记忆保留了较多场景上下文，其他轨迹或其他对象的记忆改变了更多内容，不含特定身份的原型进一步去除实例信息但保留部分表示统计。这些对照移除的信息不同，结果要逐项报告，不能由一种随机噪声对照就断言系统完全不使用历史。教学掩码分数设正确记忆为0.8、参考为0.4，则$B=0.4$。同场景的错误历史得0.7，恢复比例为$(0.7-0.4)/0.4=0.75$；其他对象的历史得0.45，比例为0.125；若某对照得0.35，比例为$-0.125$，不应强制记为零。前者说明相当一部分收益在这个对照下保留下来，后者说明换了对象后多数收益消失。

原论文在冻结视频模型的指定设置中发现，某些不含实例内容的对照仍保留了较多收益；WorldMem对场景上下文和具体位置表现出不同程度的依赖；SAM 2的空间记忆被替换后，目标分割明显受损，在目标重新出现时尤其如此。模型、任务与消费接口各不相同，这组结果说明既不能概括为所有记忆都忽略历史，也不能把SAM 2的内容依赖推广为所有视频系统的保证。身份和事件输出需要各自重做审计，不能用掩码分数的变化代替。

\section{本章小结}
\label{sec:vision-video-summary}

回到开篇的三个问题。十帧中的同一辆车要计作一辆，靠的是跨帧对应与状态维护。单目标跟踪的目标模型从均值漂移的颜色直方图、相关滤波的判别式滤波器发展到SiamRPN++离线学到的深层特征，三者都保留“在上一位置附近局部搜索”的结构；像素对应从Lucas--Kanade与Horn--Schunck的亮度守恒和人工平滑，换成RAFT的全对相关与学习的迭代更新；卡尔曼预测加一对一分配及ByteTrack维持对象身份，D4RT在共同三维坐标中分开相机与物体运动，XMem和SAM 2用键值记忆把目标的像素外观带到后续帧。拿起与放下的区分，靠的是保留时间顺序的证据：密集轨迹路线已经把外观与运动分开描述，却在编码时丢掉了顺序；TSN以稀疏分段覆盖全过程，同样不编码段间顺序；I3D以连续片段学习局部运动，ActionFormer和MS-TCN进一步给出事件的起止区间或步骤序列，再由事件记录和对象状态解释先后与重复。找到“杯子放到哪里”的片段，靠的是可回查的索引：视觉表示和CLIP4Clip在视频库中召回候选，TRACE在长视频中按问题回查原始片段，最后仍要核验对象、动作与时间范围。

本章第一节提出的五类挑战，在各节中得到的回答程度并不相同。运动对应与身份维护有了较完整的方法链，但每个环节都依赖明确的条件：经典方法依赖亮度守恒、小位移和人工设定的目标模型，学习方法依赖与训练分布相近的输入，并且大多给不出“何时不可信”的信号。时间尺度的跨度靠多尺度表示和稀疏覆盖部分解决，在线因果约束与长视频的记忆负担则主要靠限制读取范围和压缩历史来折中。这些方法的共同条件是保留证据的来源。每个结果都应说明实际看过哪些帧、哪些历史仍可访问、哪些位置由预测或外推补出；汇总依赖身份，身份依赖对应，事件依赖状态，错误也沿同一条链传播，因此需要按输出分别度量，并用内容替换检查记忆收益是否来自正确的历史。

仍未解决的问题同样清楚。长时间遮挡和跨镜头后的身份恢复、细粒度活动的顺序判断、有限预算下的长视频观察、未完成事件的及时响应，都还没有可靠的通用方案。增加记忆容量只有在保留了任务可用信息时才有用，缩短等待则可能减少可见的运动阶段；更值得继续检验的是怎样把局部状态、原始证据覆盖和及时输出放在同一预算协议下比较。\ref{chap:vision-video-restoration}章沿着可靠的跨帧对应，讨论如何恢复和传输连续画面本身。

\paragraph{扩展阅读}
对象跟踪方向，DeepSORT在SORT的卡尔曼预测与匈牙利分配之外加入重识别网络提取的外观特征，用外观距离和运动距离共同决定关联，以减少遮挡后的身份切换\scite{vision-deepsort2017}OC-SORT仍只用运动信息关联，但在目标重新被检测到时，用实际观测回溯修正遮挡期间累积的卡尔曼误差，并在关联代价中加入运动方向的一致性，适合非线性运动的场景\scite{vision-ocsort2023}MOTR不再先检测后关联：它在DETR式的Transformer中为每个已跟踪对象保留一个跟踪查询并逐帧更新，查询本身就承载身份，可与ByteTrack“检测器加关联规则”的分工对照\scite{vision-motr2022}单目标跟踪方面，SiamRPN++先分别提取模板和搜索区域的特征再做相关，OSTrack则把两者的图像块一起送入同一个ViT，让特征提取与模板—搜索关系的建模同时进行，并在中间层提前剔除可能属于背景的搜索块\scite{vision-ostrack2022}视频目标分割方面，STM把过去帧及其掩码编码为键值记忆、由当前帧查询读取，XMem沿用了这种读取方式；Cutie在像素级记忆之上加入对象级的查询与记忆，用前景与背景分开的注意力处理外观相似的干扰物\scite{vision-stm2019,vision-cutie2024}

点跟踪方向，CoTracker用Transformer同时跟踪多个点，让各点轨迹通过注意力互相参考，点被遮挡时可以借助同一物体上其他点的运动推断其位置；它输出跨越多帧的长轨迹，而RAFT估计的是相邻两帧之间的稠密位移\scite{vision-cotracker2024}TAP-Vid提供了在真实视频和合成视频上评价任意点跟踪的基准与指标，同时考察位置精度和遮挡判断\scite{vision-tapvid2022}

活动识别与检索方向，Something-Something数据集的类别由“把某物放到某物上”“假装拿起某物”一类动作模板定义，仅凭单帧外观几乎无法区分，适合检验模型是否真正利用了时间顺序，正对应本章区分拿起与放下的问题\scite{vision-something-something2017}X-CLIP与CLIP4Clip同样从CLIP出发，但同时比较视频与句子、视频与词、帧与句子、帧与词四种粒度的相似度，以减少均值汇总丢掉的细节对应；InternVideo把遮蔽视频重建与视频—文本对比学习结合起来，在识别、检索和时序定位等多类任务上检验同一视频主干的迁移效果\scite{vision-xclip-retrieval2022,vision-internvideo2022}

\chapter{视频复原、增强与压缩}
\label{chap:vision-video-restoration}

一段夜间行车录像中，路牌在第$t$帧被运动模糊和噪声淹没，在第$t-3$帧和第$t+2$帧却各自清楚一些。我们希望得到一段能读出路牌的视频，可能还要把它放大四倍、补成慢动作，最后在有限带宽下传回监控中心。逐帧使用\ref{chap:vision-image-restoration}章的方法，会丢掉邻帧里现成的观测；直接把邻帧叠加到当前帧，又会把移动的车身纹理叠到路牌上，造成比单帧更难察觉的错误。

本章要回答的问题因此是：在规定的等待和成本之内，怎样从其他时刻取得当前画面真正需要的信息，怎样把它变成目标网格上的画面，以及怎样把相邻帧的共同内容变成解码端能够重现的少量比特。三项任务共用同一个前提，即先找到不同时刻画面之间的对应。\ref{chap:vision-video-foundations}章讲过块匹配和光流这两类经典运动估计，\ref{chap:vision-video-recognition}章又用RAFT估计像素位移；本章把这种对应用于读取像素和特征。复原与超分辨率利用对应“取回”邻帧内容，插帧利用对应“推算”未拍到的时刻，压缩则利用对应“少传”可预测的内容；它们对对应错误的承受方式不同，评价时要分别核验。

阅读本章需要\ref{chap:vision-image-restoration}章的退化模型、保真指标和图像压缩，以及\ref{chap:vision-video-foundations}章的运动估计与时间连续性度量。对应错误、遮挡、闪烁、等待和比特分配这几类困难并不是三项任务各自独有的，它们以不同强度同时出现，所以本章先把输入、输出和评价指标写成可检查的规定，再逐个说明这些挑战。后续各节的顺序跟着数据的处理过程走：先从邻帧取得并汇总可用信息，经典的运动补偿时域滤波和VBM4D与学习式的RVRT在这里对照；再把汇总结果落到原网格、更高分辨率网格或中间时刻，分别由RVRT、BasicVSR和FILM承担；编码要把同样的对应变成比特，先讲H.264/HEVC的混合编码框架，再看DVC把其中哪些模块换成网络并联合训练；最后回到评价，统一比较协议并同时核验内容、时间连续性、码率与延迟。图像压缩中的量化、概率模型和熵编码已在\ref{sec:vision-restoration-compression}节讲过，本章只说明视频新增的部分。

\section{视频复原、增强与压缩的问题与挑战}
\label{sec:vision-video-restoration-challenges}
\label{sec:vision-video-restoration-goals}

“让视频更清楚”可能指去噪、提亮、放大、补帧或减少压缩损伤。这些目标允许改变的量不同，所需的输入和参照也不同；同一个多帧网络在一种目标下合格，在另一种目标下可能越界。所以在调用任何模型之前，要先用可检查的方式写下三件事：输出在空间和时间上落在什么网格上，输出某一帧时允许读取哪些帧，以及用什么指标、在什么成本下判断结果合格。这三项规定确定之后，才能说清各种方法在什么地方困难，后面各节的模型也才能在同一口径下比较。

沿用\ref{chap:vision-image-restoration}章的记号，本章以$\bm{Y}_t$表示第$t$帧低质量观测，$\bm{X}_t$表示希望得到的目标帧，$\widehat{\bm{X}}_t$表示算法输出；$t=1,\dots,n$是帧序号，$n$为序列帧数，帧率记为$f$。每帧的高、宽和通道数分别为$H$、$W$和$c$。

\label{sec:vision-video-degradation}
规定输入与输出要回答两个问题：哪些跨帧变化属于应去除的退化，输出与输入在空间和时间网格上是什么关系。前者决定训练配对怎样构造，后者决定哪些输出位置有直接观测、哪些只能估计。

单帧中的噪声、模糊、低照度、分辨率和编码损失在视频中仍然存在，但可能随时间变化。自动曝光使同一路牌亮度改变，运动使车身边缘移动，遮挡使一块背景暂时不可见，切镜则让整个场景更换。只有任务约定要修正的变化才是退化；正常车辆运动不应被算法抹平，切镜也不应被“平滑”掉。

可以写$\bm{Y}_t=\mathcal{D}_t(\bm{X}_t)+\bm{\varepsilon}_t$概括第$t$帧的成像与处理，其中$\mathcal{D}_t$包括时变模糊、采样或编码等作用，$\bm{\varepsilon}_t$表示额外噪声。这只是任务层面的简化：长曝光会积分一段时间内的运动，不能总当作瞬时清晰图像上的固定卷积。多帧方法还要建立$\bm{X}_t$与其他时刻内容的对应，不能默认所有目标帧相同。

夜间录像中，平坦区域的随机颗粒提示噪声，移动路灯的拖尾可能来自曝光期间的运动，整帧明暗跳变可能来自曝光控制，方块边界则可能来自编码。不同成因需要不同的退化配对或数据标定。仅用合成高斯噪声训练的模型，不能据此宣称同时恢复了这些变化。

网格关系把本章的任务分成四类，表\ref{tab:vision-video-task-targets}列出它们的输入、输出与目标。复原和增强保持网格；超分辨率在每帧内增加像素，倍率记为$s$；插帧在两帧之间插入时刻$\tau\in(0,1)$；压缩的目标不是理想清晰帧，而是在比特预算内重现实际输入$\bm{Y}_t$。

\begin{table}[htbp]
  \centering
  \begin{tabularx}{\linewidth}{@{}lXXX@{}}
    \toprule
    任务 & 输入 & 输出网格 & 目标与直接观测\\
    \midrule
    复原与增强 & $\bm{Y}_{1},\dots,\bm{Y}_{n}$，每帧$H\times W$ & 同为$n$帧$H\times W$ & 清晰帧$\bm{X}_t$；每个输出位置都有退化观测\\
    超分辨率 & $n$帧$h\times w$ & $n$帧$sh\times sw$ & 高分辨率帧；新增位置只有邻帧的间接观测\\
    插帧 & 第$t$与$t+1$帧 & 时刻$t+\tau$的一帧 & 中间时刻画面；该时刻没有任何直接观测\\
    压缩 & $n$帧原视频 & 码流与重建帧$\widehat{\bm{Y}}_t$ & 重现$\bm{Y}_t$；解码端只见码流\\
    \bottomrule
  \end{tabularx}
  \caption{视频处理任务的网格关系。右列说明输出位置有无直接观测，它决定输出中哪些内容可以回查原始记录，哪些只是估计。}
  \label{tab:vision-video-task-targets}
\end{table}

\label{sec:vision-video-content-cost}
评价指标要与用途对应。观看可以容忍一定色调变化，车牌读取要求字符保持，速度测量还要求时间戳和几何关系可追溯。模型生成一个看起来清晰的字符，并不说明该字符存在于观测中；超分新增的高频细节、插帧得到的中间时刻，都需要与原始观测区分。保存原视频、处理结果及每帧所用的帧范围，才能事后检查一次改变来自输入证据还是模型先验。若用户要求把红车改成蓝车或改变行驶路线，问题已转为内容创作，见\ref{chap:vision-multimodal-generation}章。在这个前提下，本章使用四类可计算的指标：逐帧质量、时间一致性、码率和延迟。

逐帧质量沿用图像指标，只是要说明怎样在帧之间汇总。峰值信噪比（PSNR）、结构相似性（SSIM）和学习感知图像块相似度（LPIPS）已在\ref{sec:vision-restoration-challenges}节和\ref{sec:vision-restoration-psnr-ssim}节定义：PSNR由均方误差按式\eqref{eq:vision-psnr}换算，单位为分贝；SSIM在局部窗口中比较亮度、对比和结构；LPIPS比较预训练网络的多层特征。视频中，记第$t$帧输出与目标的均方误差为$\mathrm{MSE}_t$，相应峰值信噪比为$\mathrm{PSNR}_t=10\log_{10}(v_{\max}^2/\mathrm{MSE}_t)$，$v_{\max}$在8位图像中为255。视频论文通常先逐帧计算PSNR，再对帧取平均，DVC的补充材料就是这样说明的\scite{vision-dvc2019}。这与先把全部帧的均方误差平均、再换算成PSNR并不等价。教学数值：两帧的均方误差分别为1和100，逐帧PSNR为48.13和28.13 dB，平均38.13 dB；若先平均均方误差得到50.5，再换算则为31.10 dB。差帧的大均方误差在合并平均中占主导，逐帧平均对它的惩罚则轻得多，所以比较时要写明汇总方式。SSIM和LPIPS同样逐帧计算再平均。

逐帧指标看不到闪烁和字符随时间变化。时间一致性、闪烁和扭曲误差已在\ref{sec:vision-videofound-temporal-consistency}节定义（闪烁见式\eqref{eq:vision-videofound-flicker}，扭曲误差见式\eqref{eq:vision-videofound-warping-error}），本章使用其中的扭曲误差，并补充一种比较运动本身的指标。扭曲误差把相邻输出帧按运动对应搬到同一网格，在没有被遮挡的位置比较二者\scite{vision-lai2018blind}。按本章反向取样的约定，写成
\begin{equation}
  E_{\mathrm{warp}}=\frac{1}{n-1}\sum_{t=1}^{n-1}
  \frac{\sum_{\bm{p}}M_t(\bm{p})\,
    \bigl\lVert\widehat{\bm{X}}_t(\bm{p})
      -\mathcal{W}(\widehat{\bm{X}}_{t+1},\bm{u}_{t\to t+1})(\bm{p})\bigr\rVert_2^2}
  {\sum_{\bm{p}}M_t(\bm{p})}\text{，}
  \label{eq:vision-video-temporal-error}
\end{equation}
其中$\mathcal{W}(\cdot,\bm{u})$表示按位移场$\bm{u}$反向取样（\ref{sec:vision-video-warping}节给出计算），$\bm{u}_{t\to t+1}$是第$t$帧指向第$t+1$帧的光流，$M_t(\bm{p})\in\{0,1\}$是非遮挡掩码，取1表示位置$\bm{p}$在两帧中都可见。分母非零时式子才有定义。它就是式\eqref{eq:vision-videofound-warping-error}在整段上的平均，只是改用本章的输出记号与反向取样算子；若改用逐通道绝对误差，即得常见的L1变体，报告时应注明所用范数。这里有一个容易忽略的条件：对应和掩码必须来自与被测方法无关的估计，例如在原始输入上估计光流，否则方法自己的对齐错误会被同一份光流掩盖。教学数值：某可见位置当前输出为0.50，下一帧搬回来的值为0.46，平方差为0.0016；若下一帧在该处闪亮到0.60，平方差变为0.0196，扭曲误差随之增大。

扭曲误差有一个明显的漏洞：每帧都很模糊的输出，相邻帧之间也很接近。TecoGAN的作者指出，基于搬运后帧差的时间指标会被这类模糊结果欺骗，因而改为比较输出与真值“怎样运动”\scite{vision-tecogan2020}。他们定义
\begin{equation}
  \mathrm{tOF}_t=\bigl\lVert\mathrm{OF}(\widehat{\bm{X}}_{t-1},\widehat{\bm{X}}_t)
     -\mathrm{OF}(\bm{X}_{t-1},\bm{X}_t)\bigr\rVert_1\text{，}
  \label{eq:vision-video-tof}
\end{equation}
其中$\mathrm{OF}$是用固定的Farnebäck方法估计的光流，$\lVert\cdot\rVert_1$对全部位置和两个分量求绝对值之和；同文的$\mathrm{tLP}$把光流换成相邻帧之间的LPIPS距离。tOF需要真值视频：输出的运动与真值的运动越接近，值越小；模糊输出的光流往往与真值不同，就不能靠模糊得到低分。例如某位置真值相邻帧位移为$(1.5,0)$，输出估计为$(1.0,0)$，该位置贡献0.5。汇总时是求和还是按像素平均，应随结果写明。

码率衡量传输成本。对帧宽$W$、高$H$、共$n$帧、帧率$f$的视频，完整码流为$B$字节时，
\begin{equation}
  \text{平均比特/像素}=\frac{8B}{nWH},
  \qquad
  \text{平均比特率}=\frac{8Bf}{n}\ \mathrm{bit/s}\text{。}
  \label{eq:vision-video-bitrate}
\end{equation}
例如10秒、30帧每秒、$640\times360$的视频共300帧，完整码流恰为300000字节，则平均码率为240000 bit/s，平均约0.0347 bit/像素。这是构造的成本计算，不对应任何模型的实测表现。分母只计像素数，不乘颜色通道数。

一个编码器在不同参数下给出一条码率—质量曲线，单个工作点无法比较两条曲线。比较曲线常用BD码率（BD-rate），其定义见式\eqref{eq:vision-bd-rate}：对两个编码器各取若干工作点，用三次多项式把对数码率拟合成PSNR的函数，在两条曲线共同覆盖的质量区间上求对数码率的平均差，再换回比例\scite{vision-bjontegaard2001}。原提案取四个工作点；四个质量值互不相同时，过四点的三次多项式唯一存在，这是多项式插值的基本结论，可参考\BookChapterRef{chap:signal-analysis}{卷一的信号章}。取对数码率的理由是，原提案观察到在线性码率上积分会让高码率段主导结果，而同样比例的码率变化在对数尺度上等长；只在共同区间积分，是为了不外推任一条曲线。若B在整个区间上码率都是A的0.8倍，BD码率为$-20\%$，表示同等PSNR下平均节省两成码率；交换码率与质量的角色，可类似地得到同等码率下的平均PSNR差。视频编码比较BD码率时还要固定图像组长度、参考结构和测试序列，DVC的补充材料就按序列分别列出相对H.264的BD码率。拟合函数也会影响结果：原提案先讨论过有理函数，后来改用对数码率上的三次多项式，不同软件实现的拟合方式也不完全相同，所以BD码率应注明计算实现。

延迟衡量某一帧从拍到出结果要等多久。设输出第$t$帧需要读到第$t+r$帧，则
\begin{equation}
  L=\frac{r}{f}+T_{\mathrm{comp}}+T_{\mathrm{buf}}\text{，}
  \label{eq:vision-video-latency}
\end{equation}
三项依次是等待后帧、计算和显示或传输缓冲的时间。每秒30帧、等待两帧时，仅第一项就约为66.7毫秒。延迟与吞吐量不同：吞吐量描述每秒处理多少帧，批处理可以让吞吐很高，却让首帧等待很久。

\subsection{跨帧对齐误差带来的伪影}
\label{sec:vision-video-challenge-alignment}

多帧方法的收益来自把别的时刻对同一内容的观测搬到当前帧，搬错位置就会把错误的内容混进来。路牌上一个白色字母在第$t$帧位于横坐标100，在第$t-1$帧位于103；若估计的位移是2个像素，从邻帧取来的就是字母边缘与背景的混合，与当前帧合成后出现两条错开的边缘，即重影或双边缘。位移误差哪怕只有零点几个像素，在细笔画上也会变成可见的模糊；在压缩中，对应误差留给残差，块与块的运动不一致还会在块边界上形成方块感。

经典方法用显式的运动估计回答这个问题。块匹配为每个块搜索一个使匹配误差最小的位移（\ref{sec:vision-videofound-block-matching}节），计算规则、结果可解释，但一个块内只有一个位移，块里同时有车身和背景时必然有一部分对错；Horn--Schunck和Lucas--Kanade在亮度恒定和小运动假设下估计逐像素光流（\ref{sec:vision-videofound-horn-schunck}节、\ref{sec:vision-videofound-lucas-kanade}节），金字塔把大运动化为各层的小运动，可在运动边界、无纹理区域和亮度变化处仍会出错。经典去噪器VBM4D在沿运动轨迹堆叠图块时加入位置惩罚，并在相邻块差异过大时截断轨迹，用来限制错误对应进入滤波\scite{vision-vbm4d2012}。

学习方法从两处改变了这一环节。其一，光流网络按下游任务的损失端到端微调，学到的是对重建有用的对应，而不只是端点误差小的对应；其二，可变形对齐在初始对应附近读取多个位置，再由网络决定权重，使单个位置的误差不至于直接决定结果。代价是对应的含义变得不那么直观，错误也更难从中间结果上看出来。\ref{sec:vision-video-warping}节先讲反向取样和可见性检查，\ref{sec:vision-video-fusion}节讲RVRT怎样在光流附近多位置取样；编码中的对应误差由残差承担，见\ref{sec:vision-video-rate-allocation}节。

\subsection{遮挡与新显露区域带来的信息缺失}
\label{sec:vision-video-challenge-occlusion}

即使运动估计完全正确，邻帧也未必看到了当前帧需要的内容。一辆车从路牌前驶过，第$t$帧里刚被让出来的路牌左半部分，在第$t-1$帧中还被车身挡着；这块新显露的区域在过去帧中不存在可取的观测，任何从$t-1$搬来的值都是车身而不是路牌。插帧更极端：两侧都被遮挡的区域在任何输入中都没有记录。困难在于反向取样不会报错，它总能在邻帧中读出一个数值。

经典方法先检测，再退回。前后向一致性检查（\ref{sec:vision-video-warping}节式\eqref{eq:vision-video-fb-consistency}）把往返位移不能闭合的位置判为不可靠；VBM4D在遮挡或切镜导致相邻块差异超过阈值时中断轨迹，只用剩下的时间段滤波\scite{vision-vbm4d2012}；运动补偿时域滤波在补偿后的差异很大时减小时间权重，退回到当前帧（\ref{sec:vision-video-fusion}节）。在编码器中，新显露区域的块可以改用帧内预测，由率失真代价决定（\ref{sec:vision-video-rate-allocation}节）。

学习方法通常不显式输出遮挡掩码，而是让注意力或融合网络为不可靠的读取分配低权重，双向传播还可以从未来帧取得过去帧没有的内容。对于两侧都看不到的区域，网络只能依据训练数据中的先验补出内容；FILM的作者加入Gram矩阵损失，正是为了让大面积去遮挡区域的补全不至于过度模糊\scite{vision-film2022}。补出来的内容看起来合理，却不是观测，这一区别要在输出中保留。本章\ref{sec:vision-video-warping}节讲可见性检查，\ref{sec:vision-video-super-resolution}节讲双向传播怎样从未来帧取得内容，\ref{sec:vision-video-interpolation}节讲插帧中两侧都不可见区域的合成。

\subsection{时间一致性与逐帧质量的冲突}
\label{sec:vision-video-challenge-consistency}

逐帧做得最好的结果，连起来播放未必最好。若每帧独立去噪或放大，各帧补出的细节彼此独立，路牌上的纹理和字符边缘在相邻帧之间轻微跳动，播放时表现为闪烁；反过来，若强行让相邻输出相近，正常的运动会被拖慢，细节被平均掉，甚至留下前几帧的残影。前者在逐帧PSNR上可能更高，后者在扭曲误差上更低，二者不能同时由单一指标保证。

经典的运动补偿时域滤波用一个权重在两者之间折中：当前输出取当前观测与上一输出（经运动补偿）的加权和。\ref{sec:vision-video-fusion}节会推导，在对应正确、噪声独立的条件下，时间权重越大，噪声方差降得越多，但一次错误对应的残留也衰减得越慢。这个权重通常按补偿后的差异逐像素调整，本质上是在逐帧保真与时间平滑之间按局部可靠性分配。

学习方法把这种折中交给网络：循环传播让相邻输出共享同一个状态，输出自然更连贯；注意力权重相当于按内容自适应的时间权重。但网络的训练损失大多仍是逐帧重建误差，时间一致性只是间接得到；状态传播在带来连贯的同时也会传播错误。\ref{sec:vision-video-fusion}节讲经典滤波与RVRT的聚合和传播，\ref{sec:vision-video-temporal-quality}节说明怎样同时核验逐帧质量和时间连续性。

\subsection{等待范围与延迟约束}
\label{sec:vision-video-restoration-availability}

能读的帧越多，可用信息越多，但读未来的帧就要等待，读更多的帧就要缓存。输出第$t$帧时，逐帧处理只读当前帧；局部多帧处理读窗口$[t-l,t+r]$；循环传播把已处理的帧压缩进状态，单向时只含过去，双向时还含未来。按式\eqref{eq:vision-video-latency}，右侧等待$r$帧就至少多出$r/f$秒。离线方法若对整段做后向传播，右侧等待就是整段长度，应记录片段长度与缓存成本；把长视频切成有限片段处理，又会改变片段边界附近可用的帧。观察范围与信息到达条件的一般讨论见\ref{sec:vision-video-availability}节。

循环结构与在线因果是两回事。循环描述计算怎样复用：同一模块逐帧更新状态。因果描述输入何时可见：输出第$t$帧时是否读过$t$之后的帧。双向循环网络是循环的，却不是因果的；一个只读过去五帧的窗口网络不循环，却是因果的。把“循环”误读成“可在线”，是本章模型最常见的误用。

经典编码器早就面对同样的取舍。双向预测帧要参考显示顺序上更晚的帧，编码器必须先编码那一帧，于是解码顺序与显示顺序不同，解码端要多等若干帧才能输出；对延迟敏感的场合可以只用前向预测，DVC的单参考设置也属于这一类。学习式复原模型多为离线设计：BasicVSR和RVRT都要读到整段的末尾，FILM插中点必须等后一帧到达。\ref{sec:vision-video-neighbor-selection}节比较各种读取范围，\ref{sec:vision-video-codec-reference}节讲编码中的参考顺序，\ref{sec:vision-video-quality-rate-latency}节把质量与等待放在同一条曲线上比较。

\subsection{帧间冗余与码率分配}
\label{sec:vision-video-challenge-rate}

相邻帧的大部分内容相同，逐帧独立压缩会反复传输它们。帧间编码只传“怎样从已解码的帧预测当前帧”和“预测之后还差多少”两部分，但比特在两部分之间怎样分配并不显然：运动描述得越细，预测越准、残差越少，运动本身却越贵；新显露区域、切镜和曝光变化处，帧间预测可能还不如帧内预测。再加上首帧和随机访问点必须独立编码，一段视频的比特要在关键帧、运动、残差和模式选择之间分配。DVC的实验给出了一个具体量级：在其报告的三个工作点上，运动信息约占总码率的27\%到36\%，且质量要求越高，运动所占比例越低\scite{vision-dvc2019}。

经典编码标准用混合编码框架回答：块运动补偿负责帧间冗余，变换与量化负责空间冗余，熵编码负责统计冗余，编码器再用拉格朗日率失真代价为每个块选择预测模式和划分\scite{vision-rdo-video1998}。各模块按手工设计分别调优，标准只规定码流语法和解码过程。学习方法保留了这一结构，但把运动估计、运动压缩、运动补偿和残差变换换成网络，用一个率失真目标联合训练，使光流网络学会估计“便于压缩”的运动。\ref{sec:vision-video-codec}节按参考依赖、运动与残差编码、解码重建三步，先讲H.264/HEVC再讲DVC。

这五项挑战沿着数据流依次出现，本章也按这个顺序回答。前三项发生在“从邻帧取得信息”这一步：先决定读多远，再在对齐时处理位移误差与可见性，最后在聚合与传播时权衡时间平滑与逐帧细节，经典滤波与RVRT在这里对照。取得的信息随后要落到某个输出网格上，网格越远离原观测，遮挡和先验补全的问题越突出，于是从原网格恢复依次过渡到超分辨率和插帧。等待范围在每个模型中都要交代，在编码中则体现为参考顺序。码率分配是编码独有的问题，它要求对应能被传输和重现，所以放在复原之后。最后的核验把逐帧质量、时间一致性、码率和延迟放回同一协议下比较。

\section{跨帧有效信息的取得}
\label{sec:vision-video-information-fusion}

本节要解决的问题是：给定低质量序列$\bm{Y}_1,\dots,\bm{Y}_n$和目标时刻$t$，为第$t$帧的每个位置取得其他时刻对同一内容的可靠观测，输出一份与第$t$帧网格对齐的特征或画面。难点正是上一节的前三项挑战：同一内容在不同帧的位置不同，需要对应，对应有误差就产生伪影；对应正确时内容仍可能被遮挡或曝光改变，需要判断可用性；为了看得更远而传递的状态，又会在平滑与细节之间取舍，并积累错误。

这些难点按处理顺序依次出现，所以本节按步骤展开：先决定读哪些帧，再把它们对齐到目标网格，最后聚合并向后传播。对同一难点，不同方法给出不同回答。读取范围上，滑动窗口在每个输出时刻重新读取固定邻帧，范围有限但互不干扰；循环方法把历史压缩进状态，看得远但串行且会传播错误；RVRT在片段内并行联合处理几帧，片段之间再循环传递，是两者的折中。对齐上，块匹配为每个块取一个位移，光流为每个像素取一个位移，可变形对齐在光流附近取多个位置再加权；越往后，对单个位移误差的容忍越强，需要学习的量也越多。聚合上，经典的运动补偿时域滤波用一个按局部差异调整的权重递归平均，VBM4D沿运动轨迹把相似图块堆成组再协同收缩，RVRT则用注意力在相邻片段之间读取并传播特征。RVRT放在最后一步完整讲解，因为它的主要改动落在对齐后怎样聚合与传播；以双向循环为核心的BasicVSR放在\ref{sec:vision-video-basicvsr}节，与超分辨率任务一起讲。

\subsection{邻帧范围的选择}
\label{sec:vision-video-neighbor-selection}

选帧要回答：哪些时刻可能观察到第$t$帧缺少的内容，读取它们要付出多少等待和计算。邻帧的价值不在“多一张图”，而在它是否看到了目标位置需要的内容。

固定窗口直接读取当前帧前后各$k$帧，计算范围明确，各输出时刻相互独立，可以并行；代价是每帧特征被相邻$2k+1$个输出重复读取，且窗口外的清晰观测无法利用。单向循环让历史信息沿时间累积，同一模块复用于各帧，每帧只处理一次；但状态容量有限，会压缩细节并传播错误，序列开头的帧也几乎得不到上下文。双向循环再加一条从后向前的状态，使每帧同时获得整段的前后信息，代价是等待整段输入并保存中间状态。片段循环先把相邻$N$帧组成片段并行处理，再在片段之间传递状态，介于两者之间。表\ref{tab:vision-video-range-strategies}按同一组维度比较这四种范围。

\begin{table}[htbp]
  \centering
  \begin{tabularx}{\linewidth}{@{}lXXXX@{}}
    \toprule
    范围策略 & 可见范围 & 右侧等待 & 每帧特征计算次数 & 主要风险\\
    \midrule
    固定窗口 & $[t-k,t+k]$ & $k$帧 & 约$2k+1$次 & 错过窗口外的清晰观测\\
    单向循环 & $[1,t]$ & 0帧 & 1次 & 状态遗忘与错误传播；开头帧上下文少\\
    双向循环 & $[1,n]$ & 至整段末尾 & 每个方向1次 & 等待整段；缓存全部中间状态\\
    片段循环 & 片段内联合，片段间传播 & 依传播方向而定 & 每层1次 & 片段内运动过大时联合处理失效\\
    \bottomrule
  \end{tabularx}
  \caption{邻帧范围的四种策略。“每帧特征计算次数”指一帧的特征在整段处理中被重新提取的次数，固定窗口可以缓存特征来减少重复，但会增加内存。片段循环若交替方向，也会读到未来帧。}
  \label{tab:vision-video-range-strategies}
\end{table}

用路牌序列看范围的影响。某块车牌在$t-3$帧清晰、$t$帧被雨滴遮住、$t+2$帧再次出现，半径$k=1$的窗口看不到两次清晰观测；单向循环可能把$t-3$帧的特征带来，却不能保证细笔画在三次更新后仍保存在状态中；双向循环能同时利用$t-3$与$t+2$，但必须等到$t+2$帧到达。若中间发生切镜，继续传播旧场景特征反而有害，系统可按\ref{sec:vision-videofound-shot-boundary}节的镜头边界检测重置状态。BasicVSR的作者把测试序列切成更短片段处理，观察到可用时间范围缩短时重建质量下降，据此认为远处帧的信息不应忽略\scite{vision-basicvsr2021}。这说明范围有收益，但收益仍要和遮挡、切镜与等待一起检查。

\subsection{邻帧内容的对齐}
\label{sec:vision-video-warping}

选定邻帧$s$后，要回答“目标帧位置$\bm{p}$的内容在邻帧的哪里”。输入是第$t$帧与第$s$帧（或它们的特征），输出是把邻帧内容搬到第$t$帧网格后的结果。难点是位移不是整数，估计可能有误，被搬来的位置也可能根本不可见。对齐有三条由粗到细的路线：按块取一个位移，按像素沿光流读取单个位置，或在光流附近学习多个取样位置再加权；越往后越能容忍局部误差，但需要估计或学习的量也越多。

最早被大量使用的是块匹配。它把当前帧切成固定大小的块，在邻帧的搜索窗口内找使绝对差之和最小的位置，搜索方法与代价见\ref{sec:vision-videofound-block-matching}节。块匹配的结果是一组块位移，同一块内所有像素共用一个位移，这正是视频编码标准采用的运动表示（\ref{sec:vision-video-rate-allocation}节）。它在块内运动一致时有效，块内同时包含前景和背景、或发生旋转和缩放时，总有一部分像素对错。逐像素的光流消除了块的限制：Horn--Schunck用全局平滑项、Lucas--Kanade用局部窗口最小二乘，在亮度恒定和小位移的假设下求每个像素的位移，配合图像金字塔处理大运动（\ref{sec:vision-videofound-horn-schunck}节、\ref{sec:vision-videofound-lucas-kanade}节）。

本章三个学习模型RVRT、BasicVSR和DVC都用同一个光流网络SPyNet初始化对应，它把Lucas--Kanade式的由粗到细金字塔直接搬进了网络\scite{vision-spynet2017}。SPyNet先把两帧$\bm{I}_1,\bm{I}_2$各建成5级图像金字塔，第$k=0$级最粗，每级分辨率是下一级的一半。在第$k$级，它把上一级的光流放大两倍得到$u(\bm{V}_{k-1})$，用它把第二帧反向取样到第一帧的网格，再由一个小卷积网络$G_k$估计剩余的位移：
\begin{equation}
  \bm{v}_k=G_k\bigl(\bm{I}^{k}_1,\ \mathcal{W}(\bm{I}^{k}_2,u(\bm{V}_{k-1})),\ u(\bm{V}_{k-1})\bigr),\qquad
  \bm{V}_k=u(\bm{V}_{k-1})+\bm{v}_k\text{，}
  \label{eq:vision-video-spynet}
\end{equation}
最粗一级的初始光流取零。$G_k$的输入是两幅RGB图像和两通道光流拼成的8通道张量，由5层$7\times7$卷积组成，输出通道依次为32、64、32、16、2，除最后一层外每层后接ReLU。原文用Flying Chairs等合成数据上的端点误差逐级训练这5个网络，每级以上一级的网络初始化。由于每一级只需估计很小的修正（原文称通常不到一个像素），$G_k$可以很小；这与经典金字塔光流“粗层看大运动、细层补小运动”的分工相同，区别在于每一级的更新规则由数据学到，而不是由亮度恒定方程推出。

无论位移来自块匹配、经典光流还是网络，取值方式都一样。设目标帧位置为$\bm{p}=(x,y)^\top$，由目标帧$t$指向邻帧$s$的光流为$\bm{u}_{t\to s}(\bm{p})$，它可以由SPyNet、\ref{sec:vision-video-raft}节的RAFT或其他方法估计。为了在目标网格中重建同一内容，应到邻帧的位置$\bm{p}+\bm{u}_{t\to s}(\bm{p})$取值。方向约定很重要：若误用$s\to t$的位移而不转换，通常不能靠简单取负得到正确采样位置，因为两者定义在不同帧的网格上。

\term{反向取样}（backward warping）逐个遍历输出位置，到源图像查询内容。与之相对的前向搬运把源像素推到目标位置，会形成空洞或多对一写入冲突；反向取样保证每个输出位置恰好得到一个值。查询位置通常不是整数。若横纵坐标分别为$j+\alpha,i+\beta$，其中$0\leq\alpha,\beta\leq1$，双线性插值用周围四个像素计算某个通道的值：
\begin{equation}
  \begin{aligned}
    \widetilde y={}&(1-\alpha)(1-\beta)y_{i,j}
       +\alpha(1-\beta)y_{i,j+1}\\
       &+(1-\alpha)\beta y_{i+1,j}
       +\alpha\beta y_{i+1,j+1}\text{。}
  \end{aligned}
  \label{eq:vision-video-bilinear}
\end{equation}
这里的四个$y$都来自邻帧$s$同一通道，省略了时间与通道标记。插值核的一般讨论见\ref{sec:vision-formation-interpolation}节。权重随距离连续变化且和为1，因此结果对位移可导，光流网络可以通过取样结果接收梯度；但插值只说明怎样读数，不保证读到的位置确实可见。

在教学网格中，四角值从左上到右下为$0.2,0.6,0.4,1.0$，取样位置为$\alpha=0.25,\beta=0.5$。对应权重为$0.375,0.125,0.375,0.125$，得到
\[
  \widetilde y=0.375(0.2)+0.125(0.6)
       +0.375(0.4)+0.125(1.0)=0.425\text{。}
\]
图\ref{fig:vision-bilinear-warp}把位置与权重放在同一网格中。若移动字母的位移估计少了一个像素，取来的可能是字母边缘与背景的混合，和当前帧融合后就会产生双边缘。

\begin{figure}[htbp]
  \centering
  % 四角与查询点均为正文可复算的教学数值。
\begin{tikzpicture}[x=1cm,y=1cm,font=\figurefont\small]
  \draw[BookRule,line width=0.6pt] (1,0) rectangle (4.5,-2.8);
  \foreach \x/\y/\val/\wt/\anchor in {
    1/0/0.2/0.375/south west,
    4.5/0/0.6/0.125/south east,
    1/-2.8/0.4/0.375/north west,
    4.5/-2.8/1.0/0.125/north east}{
    \fill[BookBlue] (\x,\y) circle (2pt);
    \node[anchor=\anchor] at (\x,\y) {\val\ (\wt)};
  }
  \draw[BookMuted,dashed,line width=0.6pt] (1.875,0) -- (1.875,-2.8);
  \draw[BookMuted,dashed,line width=0.6pt] (1,-1.4) -- (4.5,-1.4);
  \fill[BookTeal] (1.875,-1.4) circle (2.5pt);
  \node[above right] at (1.875,-1.4) {取样点};
  \draw[-{Stealth[length=2mm]},BookInk,line width=0.8pt] (0.6,0.65) -- (4.8,0.65);
  \node[anchor=west] at (4.8,0.65) {$x$};
  \draw[-{Stealth[length=2mm]},BookInk,line width=0.8pt] (0.5,0.45) -- (0.5,-3.2);
  \node[below] at (0.5,-3.2) {$y$};
  \node[anchor=west,align=left] at (5.7,-0.6) {$\alpha=0.25$\\$\beta=0.5$};
  \node[anchor=west,align=left] at (5.7,-2.0) {四角加权\\$\widetilde y=0.425$};
\end{tikzpicture}

  \caption{双线性取样的教学网格。横坐标向右、纵坐标向下；查询点位于两行的中间、偏左四分之一处。四角旁括号内是插值权重，权重和为1，输出为0.425。图中线段表示插值邻域，不表示物体真实运动轨迹。}
  \label{fig:vision-bilinear-warp}
\end{figure}

图像对齐移动像素值，特征对齐移动网络提取的表示。特征对齐的好处是取样误差造成的模糊可以由后续网络部分补偿；BasicVSR报告，把特征对齐换成图像对齐会使PSNR下降0.17 dB\scite{vision-basicvsr2021}。两者都要统一坐标和尺度：在下采样$2^\ell$倍的特征上取样时，位移数值也要除以$2^\ell$。越界位置需有明确的填充及有效性标记。

可见性要单独判断。邻帧中的位置可能已被前景挡住，插值照样给出一个数值。常用的筛选线索是前后向一致性：从$\bm{p}$沿$\bm{u}_{t\to s}$到达邻帧，再沿邻帧到目标帧的光流$\bm{u}_{s\to t}$回来，若内容可见，应大致回到出发点，即
\begin{equation}
  \bigl\lVert\bm{u}_{t\to s}(\bm{p})
   +\bm{u}_{s\to t}\bigl(\bm{p}+\bm{u}_{t\to s}(\bm{p})\bigr)\bigr\rVert_2
  \leq\delta\text{，}
  \label{eq:vision-video-fb-consistency}
\end{equation}
其中$\delta$是以像素计的容许阈值，$\bm{u}_{s\to t}$在非整数位置同样用双线性插值读取。教学数值：若$\bm{u}_{t\to s}(\bm{p})=(3,0)^\top$，回程读到$(-2.6,0.1)^\top$，和为$(0.4,0.1)^\top$，长度约0.41像素；取$\delta=1$时判为一致。若该位置在邻帧被遮挡，回程光流描述的是遮挡物的运动，和可能达到数个像素，从而被排除。这个检查依赖两份光流本身，两者同时出错时仍会漏判。

即使对应和可见性都成立，曝光差异和反光也可能改变亮度，光度校正或可靠性权重属于额外判断。单一位置取样对光流误差十分敏感，于是出现了第三条路线：\term{可变形对齐}（deformable alignment）在初始对应附近由网络预测若干偏移，到多个位置取样，再加权合成，让网络自己在附近寻找正确内容。它的代价是偏移由谁预测、权重怎样形成都必须交代清楚；RVRT的引导可变形注意力是这一路线的具体实现，见下一小节。

\subsection{对齐信息的聚合与传播}
\label{sec:vision-video-fusion}
\label{sec:vision-video-rvrt}

对齐之后，要把当前帧与多份对齐结果合成一份输出，并决定把什么传给后续帧。输入是当前帧（或其特征）和若干对齐后的邻帧，输出是更新后的当前结果与传递状态。难点在于，错误对应即使只占少数，也可能主导合成结果；而传递出去的错误会在后续每一帧重复出现。

先用一个教学例子说明为何不能简单平均。当前观测值为0.40，对齐后的邻帧值为0.425，假设可靠性分别为1和3，加权合成为$(0.40+3\times0.425)/4=0.41875$。若错误对应带来0.90且仍赋予权重3，结果就变成0.775。分母归一化不能消除错配，多帧汇总是否有益取决于对应及权重是否可靠。下面三种方法都在回答“权重从哪里来”：运动补偿时域滤波用补偿后的差异逐像素决定，VBM4D用图块之间的相似度决定，RVRT用学到的注意力决定。

沿运动轨迹做时间滤波是经典视频降噪的一大类方法\scite{vision-brailean1995}。最简单的形式是递归滤波：把上一帧的输出按运动补偿到当前网格，记为$\widetilde{\bm{X}}_t=\mathcal{W}(\widehat{\bm{X}}_{t-1},\bm{u}_{t\to t-1})$，再与当前观测加权平均，
\begin{equation}
  \widehat{\bm{X}}_t(\bm{p})=\bigl(1-a_t(\bm{p})\bigr)\bm{Y}_t(\bm{p})
    +a_t(\bm{p})\,\widetilde{\bm{X}}_t(\bm{p})\text{，}
  \label{eq:vision-video-mctf}
\end{equation}
其中$a_t(\bm{p})\in[0,1)$是时间权重。递归的好处是只需保存一帧输出，却相当于对整段历史做了指数加权平均。

时间权重怎样影响降噪，可以在理想条件下算出来。设场景静止或运动补偿完全正确，真值在该位置为常数$x$，各帧噪声独立、均值为零、方差为$\sigma^2$，权重取常数$a$。记误差$e_t=\widehat{X}_t(\bm{p})-x$，由式\eqref{eq:vision-video-mctf}得$e_t=(1-a)\varepsilon_t+a\,e_{t-1}$，其中$\varepsilon_t$是第$t$帧的噪声。$e_{t-1}$只依赖第$t$帧之前的噪声，与$\varepsilon_t$独立，所以方差相加：
\[
  v_t=(1-a)^2\sigma^2+a^2v_{t-1}\text{。}
\]
当$|a|<1$时，这个线性递推收敛到不动点$v=(1-a)^2\sigma^2+a^2v$，解得
\begin{equation}
  v=\frac{(1-a)^2}{1-a^2}\sigma^2=\frac{1-a}{1+a}\sigma^2\text{。}
  \label{eq:vision-video-mctf-variance}
\end{equation}
第二个等号用了$1-a^2=(1-a)(1+a)$。$a=0.5$时方差降到$\sigma^2/3$，$a=0.8$时降到$\sigma^2/9$，与直接平均9帧相当。代价体现在错误的残留上：若某一帧因对应错误混入了偏差$\delta$，它在$k$帧之后仍以$a^k\delta$的大小留在输出里。$a=0.8$时，三帧后还剩$0.512\delta$；$a=0.5$时只剩$0.125\delta$。时间权重越大，降噪越强，残影也越久，这正是上一节第三项挑战的定量形式。

因此实际的滤波器让权重随局部差异变化：补偿后的上一输出与当前观测差得越多，越可能是对应错误或遮挡，权重就越小。算法\ref{alg:vision-video-mctf}给出一个教学版本，权重函数$a=a_{\max}\max(0,1-d/\kappa)$是为说明机制而选的简单形式，不对应某一篇论文的具体设计。

\begin{algorithm}[运动补偿时域递归滤波（教学版本）]
  \label{alg:vision-video-mctf}
  \begin{algorithmic}[1]
    \Require 观测序列$\bm{Y}_1,\dots,\bm{Y}_n$；最大时间权重$a_{\max}\in[0,1)$；差异尺度$\kappa>0$；光流估计方法
    \Ensure 滤波结果$\widehat{\bm{X}}_1,\dots,\widehat{\bm{X}}_n$
    \State $\widehat{\bm{X}}_1\gets\bm{Y}_1$
    \For{$t=2,\dots,n$}
      \State 估计光流$\bm{u}_{t\to t-1}$，计算$\widetilde{\bm{X}}_t\gets\mathcal{W}(\widehat{\bm{X}}_{t-1},\bm{u}_{t\to t-1})$
      \For{每个位置$\bm{p}$}
        \State $d\gets\lvert\bm{Y}_t(\bm{p})-\widetilde{\bm{X}}_t(\bm{p})\rvert$
        \State $a\gets a_{\max}\max(0,1-d/\kappa)$
        \State $\widehat{\bm{X}}_t(\bm{p})\gets(1-a)\bm{Y}_t(\bm{p})+a\,\widetilde{\bm{X}}_t(\bm{p})$
      \EndFor
    \EndFor
  \end{algorithmic}
\end{algorithm}

用一个像素复算。取$a_{\max}=0.8$、$\kappa=0.2$，路牌上某个已对准的像素真值约为0.5，观测依次为0.50、0.42、0.55。第2帧差异$d=0.08$，权重$a=0.8\times0.6=0.48$，输出$0.52\times0.42+0.48\times0.50=0.4584$；第3帧$d=|0.55-0.4584|=0.0916$，$a=0.8\times(1-0.458)=0.4336$，输出$0.5664\times0.55+0.4336\times0.4584\approx0.5103$。两次输出都比观测更接近0.5。第4帧一辆车驶入该位置，观测跳到0.90，$d=0.39>\kappa$，权重变为0，输出直接取0.90，不会把路牌的值拖成残影。反过来，若这次跳变其实是一个脉冲噪声，滤波器同样会把它原样保留：固定阈值区分不了“新内容”和“大噪声”。

VBM4D把非局部均值（\ref{sec:vision-transforms-nonlinear-denoising}节）“在别处找相似图块”的思路推广到视频\scite{vision-vbm4d2012}。它先沿运动轨迹逐帧跟踪一个图块，得到一个三维时空体；再在附近找与之相似的其他时空体，沿第四个维度堆叠成“组”。组内有三种相关：块内的局部相关、沿轨迹的时间相关、不同体之间的非局部相似。对组做可分离的四维去相关变换，收缩变换系数，再逆变换，把每个体的估计放回原位置并加权聚合，这一步称为协同滤波。整个过程分两级：第一级用硬阈值收缩得到初步估计，第二级以初步估计为引导做经验维纳滤波。轨迹在相邻块差异超过阈值时中断，以应对遮挡和切镜。同一框架还用于去块效应和增强。

两种经典方法的设计思想都在学习方法中延续了下来。按局部差异调节的时间权重，对应注意力中由查询与键的相似度决定的权重；把多个相似候选放在一起再加权，对应在多个取样位置上的注意力聚合；轨迹中断，对应注意力为不可靠的候选分配低权重。差别在于，经典方法的相似度量、阈值和变换是手工设定的，学习方法从配对数据中学到这些量，可以处理更复杂的退化，代价是需要配对数据、计算量大得多，而且行为更难逐项解释。

学习方法在组织聚合与传播时也有两种极端。并行方法把整段视频的所有帧一起送入网络，用注意力在帧间交换信息；它的帧间融合充分，但模型大、对长序列内存消耗巨大。循环方法逐帧更新共享参数的状态，模型小；但它每一步只能把一帧信息压进状态，RVRT的作者指出这种串行处理会造成信息损失和噪声放大\scite{vision-rvrt2022}。\ref{sec:vision-video-basicvsr}节的BasicVSR属于后一种。

RVRT（recurrent video restoration Transformer）的主要想法是：把视频切成每段$N$帧的相邻片段，片段内部的几帧并行联合更新，片段之间按循环方式传递特征；跨片段对齐使用引导可变形注意力，在光流给出的位置附近为每帧读取多个候选。$N=1$时它退化为逐帧循环，$N=n$时它变成整段并行，因此可以在两种极端之间选择折中。片段内联合处理让相邻帧直接交换信息，缓解了逐帧压缩造成的损失；片段间循环又避免了整段并行的内存开销。

RVRT的主干由Swin Transformer层搭成。一个Swin层把特征划分成不重叠的$M\times M$窗口，在每个窗口内做多头自注意力，相邻层交替使用规则划分和平移半个窗口的划分，使信息能跨窗口流动；窗口注意力的一般计算可参考\BookChapterRef{chap:transformer}{卷五的Transformer章}。设窗口内的词元矩阵为$\bm{Z}$，一层的计算顺序是
\begin{equation}
  \bm{Z}'=\bm{Z}+\operatorname{MSA}\bigl(\operatorname{LN}(\bm{Z})\bigr),\qquad
  \bm{Z}''=\bm{Z}'+\operatorname{MLP}\bigl(\operatorname{LN}(\bm{Z}')\bigr)\text{，}
  \label{eq:vision-video-swin-layer}
\end{equation}
其中$\operatorname{LN}$是层归一化，$\operatorname{MSA}$是窗口内的多头自注意力（带可学习的相对位置偏置），$\operatorname{MLP}$是中间夹GELU的两层全连接网络。若干个Swin层串联，末端接一层卷积，再与块的输入相加，就是\term{残差Swin Transformer块}（residual Swin Transformer block, RSTB），它来自图像复原模型SwinIR\scite{vision-swinir2021}：$\bm{F}_{\mathrm{out}}=\operatorname{Conv}\bigl(\mathrm{STL}_L\circ\cdots\circ\mathrm{STL}_1(\bm{F}_{\mathrm{in}})\bigr)+\bm{F}_{\mathrm{in}}$，$\mathrm{STL}$表示一个Swin层。卷积为注意力补充平移等变的局部处理，残差连接让不同深度的特征直接通到后面。

RVRT的输入是$n$帧低质量视频，形状为$n\times H\times W\times3$，整体分三段。浅层特征提取先用一层卷积，再接1个含2个Swin层的RSTB，得到$n\times H\times W\times C$的特征；超分辨率配置取$C=144$，去模糊和去噪配置取$C=192$。去模糊、去噪配置的输出与输入同尺寸，原实现另加两层步长卷积先降低特征的空间尺寸以控制计算，最后再恢复。随后特征按$N$帧一组切成片段，原文取$N=2$。中段是$L=4$个循环特征细化模块，最后经1个含2个Swin层的RSTB整理特征，由像素重排层（见\ref{sec:vision-video-basicvsr}节）把通道转换为$s$倍空间尺寸，输出$n\times sH\times sW\times3$的视频。表\ref{tab:vision-video-rvrt-config}列出各部分的配置。

\begin{table*}[tbp]
  \centering
  \small
  \begin{tabularx}{\linewidth}{@{}lXrrrrl@{}}
    \toprule
    模块 & 窗口或取样范围 & 每窗口词元数 & 层数 & 宽度 & 头数 & MLP倍率\\
    \midrule
    浅层特征提取 & 1个RSTB，$8\times8$空间窗口 & 64 & 2 & $C$ & 6 & 论文未给出\\
    循环特征细化（每个模块） & 1层融合卷积，2个修改的RSTB，$N\times8\times8$时空窗口 & 128 & 4 & $C$ & 6 & 论文未给出\\
    引导可变形注意力（每个模块） & 前一片段2帧，每帧9个候选位置 & 18 & 1 & 查询投影到$2C$ & 12组，12头 & 论文未给出\\
    重建 & 1个RSTB，$8\times8$空间窗口，像素重排 & 64 & 2 & $C$ & 6 & 论文未给出\\
    \midrule
    全模型 & 4个细化模块，合计16个修改的Swin层 & — & 20 & 144或192 & — & —\\
    \bottomrule
  \end{tabularx}
  \caption{RVRT的配置，取自原论文的实验设置，片段长度$N=2$。宽度$C$在超分辨率配置中为144，在去模糊和去噪配置中为192。“层数”指Swin层的个数，全模型一行把浅层、细化和重建三段相加。原论文报告的参数量为：超分辨率1080万（低质量输入$320\times180$时运行时间183毫秒），去模糊1360万，去噪1280万。Swin层和引导可变形注意力中多层感知机的隐藏层倍率，论文没有给出。}
  \label{tab:vision-video-rvrt-config}
\end{table*}

核心是4层循环特征细化。每层按时间顺序依次处理各片段；第$k$个片段在第$i$层的计算分三步。第一步，用引导可变形注意力把本层已经算好的前一片段特征对齐到当前片段。第二步，把对齐结果与当前片段在第$0$到$i-1$层的特征沿通道拼接，用一层卷积融合回$C$个通道。第三步，用2个修改过的RSTB细化，每个含2个Swin层：它们把二维$8\times8$窗口扩展为$N\times8\times8$的时空窗口，使片段内各帧在同一窗口中互相注意；$N=2$时每个窗口含128个位置，每层6个头。偶数层把视频序列反转后再处理，于是各层交替从前向后、从后向前传播。

\term{引导可变形注意力}（guided deformable attention, GDA）负责第一步的对齐。它要在前一片段的$N$帧中，为当前帧位置$\bm{p}$找到相关内容。SPyNet先由低质量输入估计光流$\bm{u}_{t\to s}$，用它把支持帧特征预对齐；一个由若干卷积层和ReLU组成的小网络读取当前特征、预对齐的支持特征和光流三者的拼接，为每个支持帧$s$预测$M$个偏移$\Delta\bm{p}_{s,m}$，原实现取$M=9$。第$m$个取样位置为
\[
  \bm{p}_{s,m}=\bm{p}+\bm{u}_{t\to s}(\bm{p})
                   +\Delta\bm{p}_{s,m}\text{，}
\]
非整数处沿用式\eqref{eq:vision-video-bilinear}取样。光流把搜索范围缩小到正确位置附近，偏移在附近微调取样点，注意力再决定每个取样值占多大比重，三者分工不同。

令当前查询为$\bm{q}$，第$(s,m)$个取样位置上的键和值为$\bm{k}_{s,m}$、$\bm{v}_{s,m}$，单个注意力头计算
\begin{equation}
  \begin{aligned}
    a_{s,m}&=
      \frac{\exp(\langle\bm{q},\bm{k}_{s,m}\rangle/\sqrt d)}
      {\sum_{s',m'}\exp(\langle\bm{q},\bm{k}_{s',m'}\rangle/\sqrt d)},\\
    \bm{r}&=\sum_{s,m}a_{s,m}\bm{v}_{s,m}\text{，}
  \end{aligned}
  \label{eq:vision-video-gda}
\end{equation}
其中$d$是查询和键的通道维度，求和遍历前一片段的$N$帧和每帧$M$个位置，$N=2$、$M=9$时共18个候选。查询由当前片段在上一层的特征投影得到，键从前一片段在上一层的特征中取样，值从前一片段在本层已更新的特征中取样；比较依据来自同一层次的特征，待传递的内容则是最新结果。原实现先投影再取样，以减少重复计算，并把查询投影到$2C$维，作者称这能略微提升效果。通道被分成12个可变形组分别取样，并用12个注意力头分别计算，所有组与头的结果拼接后，再经带残差连接的两层感知机交换通道信息，因为式\eqref{eq:vision-video-gda}只在空间位置之间加权。每层结束时，光流加上$M$个偏移的平均值，作为下一层的引导，使引导位置逐层修正。

回到车牌例子。光流给出的初始位置可能偏向笔画一侧，多个取样点可以同时覆盖笔画两侧，再按与当前特征的相似度聚合。若所有取样点都落在雨滴上，注意力也没有真实字符可读，它只能在给定候选中分配权重。交替方向的多层传播还意味着第$t$帧的输出读过$t$之后的帧，RVRT虽然有循环结构，仍属于离线方法。

RVRT用低质量序列与高质量目标配对训练，对所有输出帧施加逐元素的Charbonnier损失：
\begin{equation}
  L_{\mathrm{rec}}=\frac{1}{|\Omega|}\sum_{(t,i,j,k)\in\Omega}
    \sqrt{\bigl(\widehat{x}_{t,i,j,k}-x_{t,i,j,k}\bigr)^2+\epsilon^2}\text{，}
  \label{eq:vision-video-charbonnier}
\end{equation}
其中$\Omega$是所有帧、空间位置和通道的下标集合，$\epsilon>0$是小常数，RVRT取$10^{-3}$。误差远大于$\epsilon$时，每项约等于绝对误差，对少数大误差不像平方误差那样过度敏感；误差接近零时，函数仍光滑可导。原实现从$256\times256$的高质量块随机裁剪，在REDS上每个样本取30帧，在Vimeo-90K上取14帧，用Adam以批量8训练60万步，学习率从$4\times10^{-4}$按余弦退火下降。SPyNet以预训练权重初始化，前3万步冻结，之后学习率降为其余参数的四分之一；冻结是为了避免随机初始化的其他模块在训练初期把光流破坏。去噪模型按非盲设定训练，噪声标准差在0到50之间均匀抽取。推理时整段输入，各层逐片段运行，不针对每段视频重新拟合。

原文的消融实验在缩减一半修改RSTB、用12帧训练的简化模型上进行，可以看出各设计的作用：片段长度从1增至2，REDS上的PSNR从31.98升到32.10 dB，增至3时为32.07 dB，作者将不再提升归因于片段内运动较大和光流不准；去掉引导可变形注意力中的光流引导，PSNR降到30.99 dB。原论文在超分辨率、去模糊和去噪上验证了这一设计，这些结果不直接证明低照度增强效果。对切镜或错误传播，应分别查看帧内清晰度与后续连续性；重复传递的错误特征来自同一错误，不会因为经过多个片段就变成多份独立证据。

表\ref{tab:vision-video-aggregation-methods}在同一去噪任务下比较三种聚合方式。从运动补偿时域滤波到VBM4D，改变的是候选的组织方式：从“上一帧同一位置的一个值”扩大到“多条轨迹上的多个相似图块”，换来更强的降噪，代价是块匹配与变换的计算量。从VBM4D到RVRT，改变的是权重的来源：从手工的相似度和阈值换成学到的注意力，换来对复杂退化的适应能力，代价是需要配对训练数据和大得多的模型，以及输出依赖于训练分布。

\begin{table}[htbp]
  \centering
  \small
  \begin{tabularx}{\linewidth}{@{}lXXX@{}}
    \toprule
    比较维度 & 运动补偿时域滤波 & VBM4D & RVRT\\
    \midrule
    对应来源 & 每帧一份光流或块位移 & 带位置惩罚的块轨迹跟踪 & SPyNet光流加学到的多位置偏移\\
    聚合权重 & 按补偿后差异设定的时间权重 & 相似图块分组后的变换域收缩 & 查询与键相似度的注意力权重\\
    可见性处理 & 差异过大时权重降为0 & 差异超过阈值时截断轨迹 & 隐式地为不可靠候选分配低权重\\
    时间范围 & 递归，等效于指数加权的全部过去 & 轨迹覆盖的有限前后帧 & 交替方向传播，覆盖整段\\
    需要训练 & 否 & 否 & 是，需要配对数据\\
    \bottomrule
  \end{tabularx}
  \caption{三种跨帧聚合方式的比较。比较的是机制，不是在某一数据集上的质量排序；运动补偿时域滤波的设定为算法\ref{alg:vision-video-mctf}的教学版本。}
  \label{tab:vision-video-aggregation-methods}
\end{table}

\section{目标画面的恢复与采样扩展}
\label{sec:vision-video-restoration-output}

取得跨帧特征后，要把它转换成目标网格上的画面。输入是对齐聚合后的特征，或两帧原图；输出由\ref{sec:vision-video-degradation}节规定的网格决定。难点随网格而变：原网格上每个位置都有退化观测，主要难在区分细节与噪声；更高分辨率网格上新增的位置只有邻帧的间接观测；中间时刻则没有任何直接观测，只能根据两侧推算运动后再合成，遮挡带来的信息缺失在这里最突出。

这三类输出由同一组模块的不同组合实现。原网格恢复在聚合特征后接重建头即可，以RVRT的去噪、去模糊配置为例，并与VBM4D这类经典多帧滤波对照；更高分辨率恢复还需要上采样，并且更依赖长范围的亚像素观测，BasicVSR用双向循环传播和像素重排回答这一点；中间时刻恢复则把对齐方向反过来，从待生成的时刻出发去两侧读取，FILM在多尺度上学习这种读取位移，不需要光流标注。三者都使用反向取样，差别在于取样的起点、位移的来源以及输出位置有无直接观测。

\subsection{原采样画面的恢复}
\label{sec:vision-video-native-restoration}

原采样恢复要在帧时刻和空间尺寸都不变的条件下去除噪声、模糊或改善亮度。邻帧只提供补充证据，当前帧本身仍是主要观测。图像去噪的退化模型、条件均值解释和残差学习已在\ref{sec:vision-restoration-denoising}节讲过；视频版本的新增问题在于训练配对和时间稳定。

训练样本应说明目标是怎样取得的，例如对清晰序列合成规定噪声，或使用已对齐的模糊与清晰拍摄配对。若目标帧与输入帧错位一两个像素，网络会受到互相冲突的监督，学会的往往是模糊折中；合成退化与真实夜间成像不一致，也会影响迁移。

多帧去噪的收益来自邻帧对同一内容的独立噪声观测：在对应正确、噪声独立的理想条件下，平均$K$个观测可以把噪声方差降到$1/K$，式\eqref{eq:vision-video-mctf-variance}的递归滤波是同一原理的指数加权版本；对应错误或噪声在帧间相关时，这一收益随之减少。经典方法中，VBM4D不需要训练，只要给定噪声标准差就能运行，并且同一框架可以去除压缩造成的块效应\scite{vision-vbm4d2012}；它的局限是退化模型固定为加性噪声或块效应，真实夜间录像中随亮度变化的噪声、运动模糊和曝光跳变不在其模型之内。

RVRT的去噪和去模糊配置沿用\ref{sec:vision-video-rvrt}节的跨片段机制，宽度取$C=192$，浅层用两层步长卷积降低特征尺寸，重建时再恢复到原尺寸；输出尺度为1，并不意味着中间特征都保持原尺寸。原文报告，去噪模型约有1280万个参数，在$1280\times720$的输入上运行约0.2秒，去模糊模型约有1360万个参数、约0.3秒\scite{vision-rvrt2022}。去噪模型按非盲设定训练，即把噪声标准差作为已知量提供给模型；这与VBM4D需要噪声水平的条件相同，真实场景中仍要先估计它。

夜间路牌提亮后，要检查的不只是平均亮度，还包括相邻帧是否明暗跳变、白色字符是否偏色、运动边缘是否拖尾。降噪抹掉细笔画、锐化改变字符、局部增强放大闪烁，都是网格不变时仍可能出现的失真。增强服务观看时可以接受的改变，用于文字读取或亮度测量时可能不可接受。

\subsection{更高分辨率画面的恢复}
\label{sec:vision-video-super-resolution}
\label{sec:vision-video-basicvsr}

视频超分辨率的输入是$n$帧$h\times w$低分辨率画面，输出是$n$帧$sh\times sw$画面。难点在于新增像素没有直接观测，单帧放大只能依靠先验补出细节；多帧方法则希望从邻帧取回真实拍到的信息。这要求两件事同时成立：邻帧确实提供了不同的采样位置，且网络能把它们准确对齐到当前帧。

先看邻帧为什么能提供新信息。考虑一个简化的一维清晰信号$\bm{x}=(x_1,x_2,x_3,x_4)^\top$。第一帧只采到$x_1,x_3$，第二帧因半个低分辨率像素的位移采到$x_2,x_4$。在位置精确、没有模糊和噪声的教学条件下，两帧合起来给出全部四个分量，而单帧只给出两个。真实相机还进行光学模糊和像素积分，观测不是任意选出的几个高分辨率点；遮挡、光流偏差和颜色变化还会破坏互补条件。多帧超分可以利用这种观测差异，却不能宣称所有新增像素都被真实拍到。

在BasicVSR之前，视频超分常用滑动窗口：每个输出帧读取前后几帧，用复杂的对齐与融合模块处理。BasicVSR的作者把视频超分拆成传播、对齐、聚合、上采样四个功能，逐项检查哪种最简单的选择已经足够\scite{vision-basicvsr2021}。他们观察到两点：窗口方法看不到远处帧，而缩短可用范围会降低质量；单向循环方法中，序列前部的帧得到的信息少，质量随之下降。这两点分别指向“看得更远”和“前后都看”，所以BasicVSR的主要想法是用双向循环传播取得整段信息，再配以光流特征对齐、特征拼接聚合和像素重排上采样这几项简单选择。作者的结论是传播和对齐对质量与效率影响最大。

BasicVSR的输入是低分辨率序列$\bm{Y}_1,\dots,\bm{Y}_n$，每帧$h\times w\times3$。光流网络SPyNet（结构见\ref{sec:vision-video-warping}节）以预训练权重初始化，估计每对相邻帧之间两个方向的光流，每份形状为$h\times w\times2$。以前向分支为例，在时刻$t$，它把前一状态$\bm{h}^{\to}_{t-1}$按光流$\bm{u}_{t\to t-1}$反向取样到当前坐标，再与当前低分辨率帧一起送入由30个残差块组成的更新网络，得到$\bm{h}^{\to}_t$；状态有64个通道，形状为$h\times w\times64$。后向分支对称地读取$t+1$。用$\mathcal{W}$表示反向取样、$R_{\to}$和$R_{\leftarrow}$表示两路更新网络：
\begin{equation}
  \begin{aligned}
    \bm{h}^{\to}_t&=
      R_{\to}\!\left(\bm{Y}_t,\,
      \mathcal{W}(\bm{h}^{\to}_{t-1},\bm{u}_{t\to t-1})\right),\\
    \bm{h}^{\leftarrow}_t&=
      R_{\leftarrow}\!\left(\bm{Y}_t,\,
      \mathcal{W}(\bm{h}^{\leftarrow}_{t+1},\bm{u}_{t\to t+1})\right)\text{。}
  \end{aligned}
  \label{eq:vision-video-basicvsr}
\end{equation}
序列两端缺少的状态初始化为零。两路的参数各自在所有时刻共享，循环状态的一般性质可参考\BookChapterRef{chap:recurrent-networks}{卷五的循环网络章}。这里的状态不经门控，直接由残差块覆盖更新；能保留多少历史取决于训练学到的变换。原文只把更新网络称为“残差块的堆叠”，没有给出当前帧与取样状态在入口处怎样合并，也没有给出块内的层序与卷积核大小；读者复现时应以作者公开的实现为准，本书不补写这些细节。

聚合与上采样在每个时刻独立进行。把$\bm{h}^{\to}_t$与$\bm{h}^{\leftarrow}_t$沿通道拼接得到$h\times w\times128$的特征，再由若干卷积和\term{像素重排}（pixel shuffle）放大。像素重排是一种无参数的重新排列：把$h\times w\times(r^2c)$的特征中每个位置的$r^2$组通道，分别放到放大后$r\times r$小块的各个位置，得到$rh\times rw\times c$。例如$r=2$、$c=1$时，某位置的4个通道值$(v_1,v_2,v_3,v_4)$依次成为放大网格中该$2\times2$小块的左上、右上、左下、右下四个像素。它只改变排列，新增高分辨率值由前面卷积学到的通道内容决定。BasicVSR做4倍放大，原文说上采样模块由多个卷积和像素重排组成，输出$4h\times4w\times3$的画面，具体层数与通道没有在论文中列出。表\ref{tab:vision-video-basicvsr-layers}汇总论文给出的结构。

\begin{table*}[tbp]
  \centering
  \small
  \begin{tabularx}{\linewidth}{@{}>{\raggedright\arraybackslash}p{22mm}>{\raggedright\arraybackslash}X>{\raggedright\arraybackslash}X>{\raggedright\arraybackslash}p{17mm}>{\raggedright\arraybackslash}X>{\raggedright\arraybackslash}X@{}}
    \toprule
    阶段或层组 & 层类型 & 卷积核与步长 & 输出通道 & 重复次数 & 输出尺寸\\
    \midrule
    输入 & — & — & 3 & — & $h\times w$，每帧\\
    光流估计（SPyNet，每级） & 卷积，除末层外接ReLU & $7\times7$，步长1 & 32, 64, 32, 16, 2 & 每级5层，共5级 & $h\times w\times2$，每对相邻帧两个方向\\
    前向传播：对齐 & 双线性反向取样 & — & 64 & 每帧1次 & $h\times w\times64$\\
    前向传播：更新 & 残差块 & 论文未给出 & 64 & 30 & $h\times w\times64$\\
    后向传播：对齐与更新 & 同前向，参数独立 & 论文未给出 & 64 & 30 & $h\times w\times64$\\
    聚合 & 两路状态沿通道拼接 & — & 128 & 1 & $h\times w\times128$\\
    上采样 & 卷积与像素重排 & 论文未给出 & 3 & 论文未给出 & $4h\times4w\times3$\\
    \bottomrule
  \end{tabularx}
  \caption{BasicVSR的结构，按原论文正文与附录整理，输入为4倍超分辨率的低分辨率帧。SPyNet的层配置取自SPyNet原论文，BasicVSR未说明对它作了改动。原文给出的参数量为：光流网络140万，其余部分490万，合计630万；在$180\times320$的低分辨率输入上，表中报告的运行时间为63毫秒。}
  \label{tab:vision-video-basicvsr-layers}
\end{table*}

训练使用低分辨率序列与高分辨率目标，用式\eqref{eq:vision-video-charbonnier}的Charbonnier损失约束全部输出帧，BasicVSR取$\epsilon=10^{-8}$。低分辨率输入由两种规定的退化生成：双三次下采样，以及先做标准差1.6的高斯模糊再每隔四个像素取样；两种退化分别训练和评价。在REDS上每个样本取15帧并对15个输出都计算损失，在Vimeo-90K上把序列翻转拼接成14帧以加长传播；低分辨率块$64\times64$，批量8，用Adam和余弦退火训练30万步，主体学习率$2\times10^{-4}$。光流网络以预训练权重初始化，前5000步冻结，之后以$2.5\times10^{-5}$的较小学习率一起更新；它不使用光流真值，而是通过重建损失学习对超分有用的对应。推理时把整段视频一次输入。同一论文还提出IconVSR，在每隔五帧的关键帧上用额外特征提取器补充信息，并把后向分支特征送入前向分支；这些是对BasicVSR的扩展，不属于上述基本结构。

三帧车牌例子可以检查每个组件的作用。第一帧清楚的左侧笔画经前向状态传到第二帧，第三帧清楚的右侧笔画经后向状态补充第二帧。去掉对齐，同一位置可能来自不同笔画；去掉传播，只剩当前帧；去掉上采样，输出网格不会增加。拼接只让两路信息同时可用，是否被正确利用仍由训练和输入条件决定。图\ref{fig:vision-basicvsr-propagation}把两条时间依赖与一个前向更新分别画出，可以沿取样、合并和更新检查历史怎样进入当前状态。

\begin{figure}[htbp]
  \centering
  % 上半图省略每列共有的输入/重建边，以突出两条时间依赖。
\begin{tikzpicture}[x=1cm,y=1cm,font=\figurefont\small,
  state/.style={draw=BookTeal,fill=BookTeal!10,minimum width=17mm,
    minimum height=7mm,line width=0.8pt},
  block/.style={draw=BookInk,fill=BookPaper,minimum width=18mm,
    minimum height=8mm,align=center,line width=0.8pt},
  arr/.style={-{Stealth[length=2mm]},BookInk,line width=0.8pt}]
  \node[anchor=west] at (0,4.8) {两条独立的时间传播};
  \foreach \x/\i in {1.3/1,5/2,8.7/3}{
    \node at (\x,4.2) {时刻$\i$};
    \node[state] (f\i) at (\x,3.5) {$\bm h^{\to}_{\i}$};
    \node[state] (b\i) at (\x,2.4) {$\bm h^{\leftarrow}_{\i}$};
  }
  \draw[arr] (f1) -- (f2);
  \draw[arr] (f2) -- (f3);
  \draw[arr] (b3) -- (b2);
  \draw[arr] (b2) -- (b1);
  \node[anchor=west] at (0,1.45) {一个前向更新};
  \node[block] (old) at (1.1,0.6) {历史状态\\$\bm h^\to_{t-1}$};
  \node[block] (warp) at (4.1,0.6) {取样对齐};
  \node[block] (concat) at (7.0,0.6) {与当前帧合并};
  \node[state] (new) at (9.4,0.6) {$\bm h^\to_t$};
  \draw[arr] (old) -- (warp);
  \draw[arr] (warp) -- (concat);
  \draw[arr] (concat) -- node[above] {$R_\to$} (new);
  \node (flow) at (4.1,-0.6) {$\bm u_{t\to t-1}$};
  \node (input) at (7.0,-0.6) {$\bm Y_t$};
  \draw[arr] (flow) -- (warp);
  \draw[arr] (input) -- (concat);
\end{tikzpicture}

  \caption{BasicVSR的时间依赖与单步更新。上半图只画两路状态传播，每列都使用该时刻输入帧，同一列两路状态共同用于重建。下半图展开一个前向更新：历史状态按光流取样，再与当前帧一起送入残差块更新。后向分支必须等未来帧可用。}
  \label{fig:vision-basicvsr-propagation}
\end{figure}
\FloatBarrier

双向方法依赖未来帧。若改成单向在线运行，结构、训练和边界条件都发生变化，应重新训练和验证，不能只删除后向结果就沿用原成绩。大运动、遮挡和切镜可能使历史特征被拖入错误区域；车牌即使看上去变清晰，也须回查字符是否有观测支持。

BasicVSR与RVRT可以在同一任务（4倍视频超分、相同退化与评价协议）下比较，表\ref{tab:vision-video-rvrt-basicvsr}列出二者的主要设计。RVRT保留了BasicVSR的光流引导与循环传播，改变了三处：把逐帧更新换成$N$帧片段的联合更新，把单位置特征取样换成光流引导下的多位置注意力取样，并用多层交替方向代替两条独立分支。换来的是片段内帧间的直接信息交换和对光流误差更强的容忍；付出的代价是每步计算更重，参数量从630万增至1080万，训练所用帧数和步数也更多。RVRT原文在$320\times180$输入上的对比中，还列出了BasicVSR的后续版本BasicVSR++（730万参数、77毫秒），RVRT为183毫秒；作者认为运行时间的差距主要来自卷积层在现有深度学习框架中优化得更充分。两者都要等待整段输入，都不是在线方法。

\begin{table}[htbp]
  \centering
  \begin{tabularx}{\linewidth}{@{}lXX@{}}
    \toprule
    设计维度 & BasicVSR & RVRT\\
    \midrule
    传播单位 & 单帧状态，两条独立方向分支 & $N$帧片段，$L$层交替方向\\
    对齐方式 & 沿光流单位置反向取样特征 & 光流引导、每帧$M$个偏移位置的注意力取样\\
    帧内细化 & 卷积残差块 & 时空窗口内的Transformer块\\
    聚合方式 & 两路状态沿通道拼接 & 对齐特征与历史各层特征拼接后由注意力细化\\
    参数量 & 630万 & 1080万\\
    在线性 & 需整段，非因果 & 需整段，非因果\\
    \bottomrule
  \end{tabularx}
  \caption{同一视频超分任务下BasicVSR与RVRT的设计差别。两者都用预训练的SPyNet初始化并以Charbonnier损失端到端训练；参数量取自各自原论文，表中只列结构差别，不表示在任意数据上的质量排序。}
  \label{tab:vision-video-rvrt-basicvsr}
\end{table}

\subsection{中间时刻画面的恢复}
\label{sec:vision-video-interpolation}
\label{sec:vision-video-film}

\term{视频插帧}（video frame interpolation）根据两个时刻的观测，估计两者之间某一时刻的画面。输入是两帧$\bm{I}_0,\bm{I}_1$，形状$H\times W\times3$；输出是时刻$\tau\in(0,1)$的画面$\widehat{\bm{I}}_\tau$。它有两侧约束；只依据历史预测未来缺少后一侧约束，属于不同设置。缺帧补全也要注明缺失长度和实际可用的邻帧，缺一帧与缺半秒难度不同。

插帧的难点在于目标时刻没有观测，网络既要推算每个像素在$\tau$时刻的位置，又要在两侧都看不到的区域补出内容。经典做法沿用运动补偿的思路：估计两帧之间的光流，假设匀速运动把它按$\tau$缩放，从两侧反向取样后按时间距离加权平均。这一做法在运动边界和遮挡处会产生重影，因为缩放后的光流定义在错误的网格上，被遮挡的一侧也照样参与平均。基于学习的光流路线先估计运动，再从两侧反向取样并由网络合成；此前常见做法用一个预训练光流或深度网络估计运动，再用另一个网络合成，这需要稀缺的光流或深度真值，并且两辆车在两帧之间移动很多像素时，单一尺度的局部匹配可能找不到同一边缘\scite{vision-film2022}。

FILM（frame interpolation for large motion）的主要想法是：用跨尺度共享权重的特征提取器，让粗尺度上的大运动与细尺度上的小运动用同一套特征表示；在此基础上从粗到细直接预测从中间时刻指向两侧输入的位移，整个网络只用帧三元组训练。作者指出，标准的特征金字塔有两个困难：快速移动的小物体在粗层消失，粗层像素少、能提供的大运动监督也少。共享权重之所以有帮助，是因为大位移在粗尺度上变成小位移，而训练数据中小位移样本充足；同一组卷积学会处理小位移后，在粗尺度上就能处理原图中的大位移。

FILM的特征提取先把每张输入建成7层图像金字塔，第$l$层分辨率是原图的$2^{-(l-1)}$。每层图像都送入同一个编码器，编码器第$d$个深度的输出记为$\bm{f}^{l,d}$，原实现共4个深度；每个深度是两层$3\times3$卷积加ReLU，之后接一个大小和步长都为2的平均池化，相同深度在所有金字塔层共享卷积权重。于是第$l$层图像在深度2的特征与第$l+1$层图像在深度1的特征具有相同空间分辨率。FILM把相同分辨率、来自不同深度和不同层的特征沿通道拼接，记为$\bm{F}^{l}=(\bm{f}^{l-2,3},\bm{f}^{l-1,2},\bm{f}^{l,1})$；最细的两层可用的项较少，分别只拼接一份和两份。这样，每个尺度上的特征都混合了“浅层看细节”和“深层看大范围”的信息，含义在各尺度之间可比，原文称之为尺度无关的特征。

运动估计从最粗层开始。第$l$层要预测从中间时刻指向第1帧的位移$\bm{w}^{l}_{\tau\to1}$：先把上一粗层结果放大到本层网格，记为$(\bm{w}^{l+1}_{\tau\to1})_{\times2}$，用它把第1帧的特征$\bm{F}_1^{l}$反向取样，再由一个小卷积网络$G^l$读取第0帧特征和取样结果，预测残差：
\begin{equation}
  \bm{w}^{l}_{\tau\to1}
   =(\bm{w}^{l+1}_{\tau\to1})_{\times2}
   +G^{l}\Bigl(\bm{F}_0^{l},\,
     \mathcal{W}\bigl(\bm{F}_1^{l},(\bm{w}^{l+1}_{\tau\to1})_{\times2}\bigr)\Bigr)\text{。}
  \label{eq:vision-video-film-flow}
\end{equation}
这与式\eqref{eq:vision-video-spynet}的SPyNet更新形式相同，差别在于FILM在共享特征上而不是在图像上做残差估计，位移的起点也换成了中间时刻。原文把$(\cdot)_{\times2}$写作两倍双线性上采样；位移若以当前网格的像素为单位，网格放大一倍时数值也应随之加倍，读者复现时要确认实现采用的单位约定。指向第0帧的位移$\bm{w}^{l}_{\tau\to0}$对称计算。原文结构图中，$G^l$为三层$3\times3$卷积加ReLU，在较粗的第3至7层共享权重，最细的两层各用自己的参数。粗尺度上原图40像素的运动在第4层只有5像素，较容易进入局部比较范围，再逐层加上细的修正。这些位移由最终中间帧的监督间接决定，原文称为面向任务的位移，它不必等于真实光流，也不能当作独立验证的光流使用。

以教学数值说明两路取样方向。设$\tau=1/2$，一辆车在两端横坐标20和28之间匀速平移，中间时刻位置为24。输出位置24应分别向第0帧的20和第1帧的28取样，位移为$-4$和$+4$。真实模型还要处理加速、遮挡和新显露的背景，不是对两个输入坐标机械取平均；两侧都没见过的区域只能由合成网络依先验估计。

合成阶段先用各层的位移把两侧的特征与该层图像一起反向取样到中间时刻，再在每个金字塔层把两侧取样结果和两份位移拼接，送入类似U-Net的解码器。解码器从粗到细逐层上采样，与本层的拼接输入合并，经两层卷积加ReLU，最后一层卷积输出$H\times W\times3$的中间帧。表\ref{tab:vision-video-film-layers}汇总论文给出的结构。原文报告，在V100 GPU上插一帧720p画面约0.39秒，峰值显存约4.5 GB；论文没有给出参数量。

\begin{table*}[tbp]
  \centering
  \small
  \begin{tabularx}{\linewidth}{@{}>{\raggedright\arraybackslash}p{22mm}>{\raggedright\arraybackslash}X>{\raggedright\arraybackslash}X>{\raggedright\arraybackslash}p{17mm}>{\raggedright\arraybackslash}X>{\raggedright\arraybackslash}X@{}}
    \toprule
    阶段或层组 & 层类型 & 卷积核与步长 & 输出通道 & 重复次数 & 输出尺寸\\
    \midrule
    输入 & 两帧 & — & 3 & 2 & $H\times W$\\
    图像金字塔 & 逐级缩小 & — & 3 & 7级 & 第$l$级$H/2^{l-1}\times W/2^{l-1}$\\
    共享编码器（每个深度） & 卷积加ReLU，再平均池化 & $3\times3$，步长1；池化$2\times2$，步长2 & 第1深度64，其余论文未给出 & 每深度2层卷积，共4个深度 & 每个深度减半\\
    尺度无关特征 & 同分辨率特征拼接 & — & 三个深度通道之和 & 每级1次 & 第$l$级分辨率\\
    残差位移估计$G^l$ & 卷积加ReLU & $3\times3$，步长1 & 2（每个方向） & 3层；第3至7级共享 & 第$l$级分辨率\\
    融合解码器（每级） & 上采样、拼接、卷积加ReLU & 论文未给出 & 论文未给出 & 每级2层 & 逐级放大\\
    输出层 & 卷积 & 论文未给出 & 3 & 1 & $H\times W\times3$\\
    \bottomrule
  \end{tabularx}
  \caption{FILM的结构，按原论文正文和结构图整理。论文正文说明所用模型的编码器从64个滤波器开始，结构图以32起始作示意，未逐项列出更深层的通道数；结构图标注的$G^l$为三层带ReLU的卷积，末层输出两通道位移。论文未给出参数量。}
  \label{tab:vision-video-film-layers}
\end{table*}

训练样本是两端与中间真实帧组成的三元组，只在$\tau=1/2$监督。损失作用在最终输出上：
\begin{equation}
  L_{\mathrm{FILM}}
  =w_{1}\lVert\widehat{\bm{I}}_{1/2}-\bm{I}_{1/2}\rVert_1
   +w_{\mathrm{VGG}}L_{\mathrm{VGG}}
   +w_{\mathrm{Gram}}L_{\mathrm{Gram}}\text{，}
  \label{eq:vision-video-film-loss}
\end{equation}
第一项是像素绝对误差，约束颜色和位置；$L_{\mathrm{VGG}}$比较预训练VGG-19网络各层特征的绝对误差，约束局部结构；$L_{\mathrm{Gram}}$比较这些特征的Gram矩阵（通道两两相关），约束纹理统计，使大运动区域不至于被平均得过于模糊。后两项的定义和作用见\ref{sec:vision-generation-perceptual-style}节。原实现在前150万步取$(w_1,w_{\mathrm{VGG}},w_{\mathrm{Gram}})=(1,1,0)$，后150万步转为$(1,0.25,40)$，使三项贡献大致相当；原文还训练了只用第一项和只加VGG项的版本：各版本在像素指标上差距很小，但原文展示只用第一项时画面偏模糊，加入Gram项后明显变锐，报告结果时应注明使用哪一版本。训练使用Vimeo-90K三元组，$256\times256$随机裁剪，批量8，用Adam优化共300万步，学习率从$10^{-4}$开始每75万步按0.464的比例衰减，并以旋转、翻转和三元组时间反转增强。原文还从视频中挖掘不同运动幅度范围的三元组，报告用较大运动范围的数据训练有助于大运动测试集。

需要更密的时间网格时，FILM递归插入中点，得到$1/4,1/2,3/4$等位置。新一轮可能把先前生成的帧作为输入，误差会继承，生成帧不能当作新增的实拍约束。在线插帧必须等待后一端：30帧每秒的输入估计中点时，中点内容至少要在后一帧到达后才能生成，相对该中点的观测等待约为16.7毫秒，还要加上计算耗时和播放缓冲。平滑慢动作适合呈现，却不能证明插出的碰撞、接触或遮挡细节真实发生过。

FILM与BasicVSR都用反向取样和多尺度运动，但目的不同，可以从三个方面比较。第一，取样起点：BasicVSR从已有帧的网格出发读取邻帧，FILM从一个尚无观测的时刻出发读取两侧。第二，运动来源：BasicVSR用预训练的SPyNet初始化后微调，FILM的运动估计完全由帧监督学到，不依赖光流或深度真值。第三，输出依据：BasicVSR的每个输出位置至少有本帧的低分辨率观测，FILM的输出完全是由两侧推算的估计。FILM因此省去了光流数据，但它的位移和输出都缺少独立的观测检查，只能靠两侧一致性和下游用途核验。

\section{有限比特下的视频编码与解码}
\label{sec:vision-video-codec}

前两节利用帧间对应改善画面，本节利用同一对应减少传输。问题是：输入原视频$\bm{Y}_1,\dots,\bm{Y}_n$，输出一条码流，使解码端只凭码流就能重建$\widehat{\bm{Y}}_t$，并在给定比特预算下让失真尽量小。相邻帧有大量相同内容，逐帧独立编码会反复传输它们；帧间编码改为用已解码的帧预测当前帧，只编码预测解释不了的部分。这里待编码的输入仍记为$\bm{Y}_t$，压缩目标是保持这个实际输入，不默认它已是前两节追求的无退化真值；压缩也可以直接处理原视频，不要求先复原。

难点在于编码器与解码器看到的信息不同。编码器能看原视频，解码器只能读码流和自己已经重建的帧；任何预测都必须只用两端共有的信息，否则两端的重建会逐帧分叉。此外，运动和残差都要花比特，减少一方可能增加另一方，必须联合权衡。

视频编码标准对这些难点的回答是\term{混合视频编码}（hybrid video coding）：用帧间或帧内预测去掉时间和空间上可预测的部分，再对预测残差做二维变换编码。从H.261起的各代视频编码标准都采用这一结构，H.264/AVC和HEVC也不例外\scite{vision-sullivan2012hevc}。图\ref{fig:vision-video-hybrid-codec}按编码一帧的顺序画出主要模块。编码器把图像划分成块；序列首帧和每个随机访问点只用帧内预测，其余帧的多数块用帧间预测，编码器为每块选定参考帧和运动矢量，作为辅助信息传给解码器。预测残差经变换、缩放、量化和熵编码写入码流。编码器内部复制一遍解码器的处理：把量化后的系数反量化、反变换，加回预测，经环路滤波后存入已解码帧缓存，供后续帧预测使用；这样两端才能产生相同的预测。标准只规定码流语法和解码过程，编码器怎样搜索运动、怎样选择模式都由实现自行决定\scite{vision-wiegand2003h264}。

\begin{figure*}[tbp]
  \centering
  % H.264/HEVC式混合编码器的主要模块，按两篇标准概述的编码流程重绘；宽约150mm，跨栏使用。
% 灰底模块组成编码器内部的解码环路，与独立解码器执行相同的计算；白底模块只在编码端出现。
\begin{tikzpicture}[x=1cm,y=1cm,font=\figurefont\small,
  enc/.style={draw=FigureBlue,fill=white,rounded corners=1.8pt,line width=0.8pt,
    minimum width=21mm,minimum height=8mm,align=center,inner sep=1.2mm},
  dec/.style={enc,draw=FigureStroke,fill=FigurePaper},
  sum/.style={circle,draw=FigureStroke,fill=white,line width=0.8pt,
    minimum size=5.5mm,inner sep=0pt},
  arr/.style={-{Stealth[length=1.8mm]},draw=FigureStroke,line width=0.8pt,rounded corners=1mm},
  side/.style={arr,dashed},
  note/.style={text=FigureMuted,font=\figurefont\footnotesize}]
  \node[enc] (in) at (0,3) {当前帧\\$\bm Y_t$};
  \node[sum] (sub) at (2.4,3) {$-$};
  \node[enc] (tr) at (4.6,3) {整数变换};
  \node[enc] (q) at (7.2,3) {量化\\（量化参数）};
  \node[enc] (ent) at (10.6,3) {熵编码\\CAVLC/CABAC};
  \node[enc,minimum width=14mm] (bs) at (14.0,3) {码流};
  \node[dec] (iq) at (7.2,1.5) {反量化\\反变换};
  \node[sum] (add) at (7.2,0.2) {$+$};
  \node[dec] (lf) at (9.8,0.2) {环路滤波\\去块、SAO};
  \node[dec] (buf) at (12.6,0.2) {已解码帧\\缓存};
  \node[enc] (me) at (12.6,1.7) {运动估计};
  \node[dec] (pred) at (2.4,-1.3) {帧内预测或\\运动补偿};
  \draw[arr] (in) -- (sub);
  \draw[arr] (sub) -- node[above,note] {残差} (tr);
  \draw[arr] (tr) -- (q);
  \draw[arr] (q) -- node[above,note] {系数} (ent);
  \draw[arr] (ent) -- (bs);
  \draw[arr] (q) -- (iq);
  \draw[arr] (iq) -- (add);
  \draw[arr] (add) -- (lf);
  \draw[arr] (lf) -- (buf);
  \draw[arr] (buf) -- (me);
  \draw[arr] (in.north) -- (0,4.1) -| (me.north);
  \draw[side] (me.west) -| node[pos=0.75,left,note] {运动矢量、模式} (ent.south);
  \draw[arr] (buf.south) |- node[pos=0.75,above,note] {参考帧} (pred.east);
  \draw[arr] (pred.north) -- node[left,note] {预测$\overline{\bm Y}_t$} (sub.south);
  \draw[arr] (2.4,0.2) -- (add.west);
  \draw[side] (me.east) -- (14.1,1.7) |- (2.4,-2.3) -- (pred.south);
  \node[note,anchor=west] at (3.0,-2.6) {灰底：编码器内部的解码环路，与解码器逐位一致};
\end{tikzpicture}

  \caption{H.264/HEVC式混合编码器的主要模块，根据两篇标准概述的编码流程重绘。实线为数据流，虚线为运动矢量和模式等辅助信息。灰底模块构成编码器内部的解码环路，独立解码器执行同样的计算；运动估计和模式选择只在编码端进行，不属于标准规定的范围。}
  \label{fig:vision-video-hybrid-codec}
\end{figure*}

这套框架的每个模块都按手工设计分别调优：块匹配找运动，固定的整数变换去相关，量化参数控制精度，熵编码器按上下文切换概率，编码器再用率失真代价为每个块选择模式。图像压缩的两条路线已在\ref{sec:vision-restoration-compression}节讲过：JPEG式的经典编码（\ref{sec:vision-restoration-classic-codec}节）用块变换、量化表和熵编码处理单幅图像，学习式编码（\ref{sec:vision-restoration-learned-codec}节）证明了变换、量化和概率模型可以联合训练，但二者都不处理帧间依赖。本节只讲视频新增的部分：帧间预测与参考结构、整数变换与量化参数、CABAC、环路滤波和模式选择。DVC沿用传统的“运动补偿—残差编码”框架，把每个模块换成网络，用一个率失真目标联合训练全部模块\scite{vision-dvc2019}。表\ref{tab:vision-video-codec-modules}按模块对照两者。编码流程中的三个步骤分别面对上面的难点：确定参考依赖解决两端信息不对称，运动与残差的编码解决比特分配，解码重建检验依赖能否在接收端独立还原。每一步先讲H.264/HEVC的做法，再讲DVC改了什么。

\begin{table}[htbp]
  \centering
  \small
  \begin{tabularx}{\linewidth}{@{}lXX@{}}
    \toprule
    模块 & H.264/HEVC & DVC\\
    \midrule
    运动估计 & 编码器自选的块搜索；可变块大小，四分之一像素精度 & 预训练后联合微调的SPyNet，逐像素光流\\
    运动信息编码 & 运动矢量由相邻块预测，差值熵编码 & 卷积自编码器压缩光流，学习的概率模型估计码率\\
    运动补偿 & 按块从参考帧复制，分数位置插值 & 双线性反向取样后由补偿网络修正\\
    残差变换 & 近似离散余弦变换的整数变换 & 尺度超先验方法的非线性变换\\
    量化 & 由量化参数控制的均匀量化 & 取整，训练时以加均匀噪声代替\\
    熵编码 & CAVLC或CABAC & 按学习的符号概率熵编码\\
    环路滤波 & 去块滤波；HEVC另有样点自适应偏移 & 不设，原文认为逐像素补偿不产生块效应\\
    模式与比特分配 & 编码时逐块最小化拉格朗日代价 & 训练时最小化率失真损失，每个$\lambda$一个模型\\
    \bottomrule
  \end{tabularx}
  \caption{混合编码框架与DVC的模块对照。两者的数据流相同，差别在于各模块由手工设计还是由网络实现，以及率失真权衡发生在编码时还是训练时。H.264/HEVC一列根据两篇标准概述整理，只列与本章相关的工具。}
  \label{tab:vision-video-codec-modules}
\end{table}

\subsection{参考依赖的确定}
\label{sec:vision-video-codec-reference}
\label{sec:vision-video-dvc}

编码第一步要决定每帧依赖什么。\term{帧内编码}（intra-frame coding）只用当前帧自身的信息，可以独立解码，作为序列起点或错误恢复点；\term{帧间编码}（inter-frame coding）依赖其他已重建帧，比特更少，却必须先解码它所依赖的帧。从一个帧内编码帧开始、到下一个帧内编码帧之前的一组帧称为图像组（group of pictures, GOP，见\ref{sec:vision-videofound-formats}节），它规定了参考链的长度。这一步的难点不是选择哪种编码，而是保证参考对两端相同，并控制依赖带来的等待。

H.264把一帧划分为若干条带，按可用的预测方式分为三类\scite{vision-wiegand2003h264}：I条带中的块全部用帧内预测；P条带中的块还可以用帧间预测，每个预测块至多有一个运动补偿预测信号；B条带中的块可以同时用两个运动补偿预测信号，例如一个来自显示顺序上更早的帧、一个来自更晚的帧。H.264还允许从多个已解码帧中选择参考，并把参考顺序与显示顺序解耦，B帧本身也可以作为其他帧的参考。下文按习惯把主要由这三类条带组成的帧称为I帧、P帧和B帧。

参考更晚的帧，意味着解码顺序与显示顺序不同。图\ref{fig:vision-video-gop-order}是一个七帧的教学图像组：P$_3$参考I$_0$，B$_1$与B$_2$同时参考I$_0$和P$_3$。要解码B$_1$，P$_3$必须先到达，所以码流中的顺序是I$_0$、P$_3$、B$_1$、B$_2$、P$_6$、B$_4$、B$_5$。编码端在拍到第3帧之前无法编码B$_1$，按每秒30帧计算，B$_1$至少多等两帧，约66.7毫秒；B帧能借两侧参考减少残差，代价正是\ref{sec:vision-video-restoration-availability}节所说的等待。只用I帧和P帧的结构没有这种重排等待，后文的DVC就属于这一类。

\begin{figure}[htbp]
  \centering
  % 七帧教学图像组：显示顺序与解码顺序；宽约100mm。实际编码器可选择不同的参考结构。
\begin{tikzpicture}[x=1cm,y=1cm,font=\figurefont\small,
  frame/.style={draw=#1,fill=#1Fill,rounded corners=1.8pt,line width=0.8pt,
    minimum width=9mm,minimum height=8mm,inner sep=0.8mm},
  ref/.style={-{Stealth[length=1.6mm]},draw=FigureStroke,line width=0.7pt},
  bref/.style={ref,dashed,draw=FigureMuted},
  note/.style={text=FigureMuted,font=\figurefont\footnotesize}]
  \node[anchor=east] at (-0.6,2.2) {显示顺序};
  \foreach \i/\t/\c in {0/I/FigureBlue,1/B/FigureRed,2/B/FigureRed,3/P/FigureYellow,
                        4/B/FigureRed,5/B/FigureRed,6/P/FigureYellow}{
    \node[frame=\c] (d\i) at ({1.2*\i},2.2) {$\mathrm{\t}_{\i}$};
  }
  \draw[ref] (d0.north) to[bend left=35] node[above,note] {参考} (d3.north);
  \draw[ref] (d3.north) to[bend left=35] (d6.north);
  \draw[bref] (d0.south) to[bend right=40] (d1.south);
  \draw[bref] (d3.south) to[bend left=40] (d1.south);
  \draw[bref] (d3.south) to[bend right=40] (d4.south);
  \draw[bref] (d6.south) to[bend left=40] (d4.south);
  \node[anchor=east] at (-0.6,0) {解码顺序};
  \foreach \k/\i/\t/\c in {0/0/I/FigureBlue,1/3/P/FigureYellow,2/1/B/FigureRed,3/2/B/FigureRed,
                           4/6/P/FigureYellow,5/4/B/FigureRed,6/5/B/FigureRed}{
    \node[frame=\c] (e\k) at ({1.2*\k},0) {$\mathrm{\t}_{\i}$};
  }
  \foreach \k in {0,...,5}{
    \pgfmathtruncatemacro{\n}{\k+1}
    \draw[ref,draw=FigureMuted] (e\k.east) -- (e\n.west);
  }
\end{tikzpicture}

  \caption{七帧教学图像组的参考关系与解码顺序。下标为显示时刻；上方实线箭头从参考帧指向P帧，下方虚线箭头为B$_1$和B$_4$的前后两个参考，B$_2$、B$_5$同理而未画出。解码顺序把每个B帧排在它的两个参考之后。这是教学构造，实际编码器可以选择不同的参考结构。}
  \label{fig:vision-video-gop-order}
\end{figure}

参考对两端相同，不只要求两端用同一帧，还要求两端对这一帧的重建逐位一致。早期标准只在一个误差容限内规定反变换，不同解码器的输出略有差异，这种差异会沿参考链累积，称为漂移；H.264把反变换定义为精确的整数运算，成为第一个使所有解码器输出完全相同的标准\scite{vision-wiegand2003h264}。这个要求同样适用于学习式编码，后文会看到它给神经网络带来的新问题。

DVC在帧间编码中只用上一重建帧作参考，即只有前向预测的低延迟结构。它的动机有两点：传统编码器的块运动估计和离散余弦变换是手工设计，整个系统没有端到端优化；此前的学习式视频压缩或只替换单个模块，或仍依赖传统块运动估计，或只优化失真而不计码率。DVC的主要想法是保留“估计运动、压缩运动、\term{运动补偿}（motion compensation）、压缩残差”这一被长期验证的结构，用网络实现每个模块，再用率失真损失共同训练。运动补偿是指用运动信息把参考帧变换成当前帧的预测；网络化之后，运动可以是逐像素的稠密光流而非块矢量，预测不会出现块边界。

DVC的数据流如下。输入当前帧$\bm{Y}_t$（$H\times W\times3$）和上一重建帧$\widehat{\bm{Y}}_{t-1}$。光流网络估计光流$\bm{v}_t$，形状$H\times W\times2$；运动编码器把它压缩为潜在表示$\bm{m}_t$，量化为$\widehat{\bm{m}}_t$，运动解码器再恢复出$\widehat{\bm{v}}_t$。运动补偿用$\widehat{\bm{v}}_t$对$\widehat{\bm{Y}}_{t-1}$反向取样，再由补偿网络修正，输出预测帧$\overline{\bm{Y}}_t$。残差$\bm{E}_t$经残差编码器变为潜在表示$\bm{z}_t$，量化为$\widehat{\bm{z}}_t$后解码为$\widehat{\bm{E}}_t$：
\begin{equation}
  \bm{E}_t=\bm{Y}_t-\overline{\bm{Y}}_t,\qquad
  \widehat{\bm{Y}}_t=\overline{\bm{Y}}_t+\widehat{\bm{E}}_t\text{。}
  \label{eq:vision-video-dvc-reconstruction}
\end{equation}
$\bm{E}_t$是编码端可算的残差，$\widehat{\bm{E}}_t$是它经编码、量化与解码后的重建。当前重建帧写入参考缓存，成为下一帧的共同条件。各网络的层结构在下一小节给出。

关键的约束是：运动补偿必须使用解码后的$\widehat{\bm{v}}_t$，而不是光流网络直接输出的$\bm{v}_t$。原始光流未经编码量化，解码端拿不到；若编码端用它产生预测，就算出了解码端无法复现的残差，也掩盖了运动压缩误差。图\ref{fig:vision-dvc-reference}按两端职责画出这条依赖，它与图\ref{fig:vision-video-hybrid-codec}中编码器内部的解码环路是同一个要求。

\begin{figure*}[tbp]
  \centering
  % 原图仅编码端可见；下路也由编码端本地执行，保证共同参考。
\begin{tikzpicture}[x=1cm,y=1cm,font=\figurefont\small,
  block/.style={draw=BookInk,fill=BookPaper,minimum width=24mm,
    minimum height=9mm,align=center,line width=0.8pt},
  arr/.style={-{Stealth[length=2mm]},BookInk,line width=0.8pt,rounded corners=1mm}]
  \node[anchor=west] at (-0.2,4.7) {编码端：从原帧与本地预测形成两类符号};
  \node[block] (srcm) at (1.4,3.6) {$\bm Y_t,\widehat{\bm Y}_{t-1}$};
  \node[block] (encm) at (5.2,3.6) {运动估计\\编码、量化};
  \node[block] (symm) at (9.0,3.6) {$\widehat{\bm m}_t$};
  \node[block] (srcr) at (1.4,2.2) {$\bm E_t=\bm Y_t-\overline{\bm Y}_t$};
  \node[block] (encr) at (5.2,2.2) {残差编码\\量化};
  \node[block] (symr) at (9.0,2.2) {$\widehat{\bm z}_t$};
  \draw[arr] (srcm) -- (encm);
  \draw[arr] (encm) -- (symm);
  \draw[arr] (srcr) -- (encr);
  \draw[arr] (encr) -- (symr);
  \draw[BookTeal,fill=BookTeal!8,line width=0.8pt] (12.7,1.65) rectangle (15.6,4.15);
  \node at (14.15,3.6) {运动码流};
  \node at (14.15,2.2) {残差码流};
  \draw[arr] (symm.east) -- node[above] {熵编码} (12.7,3.6);
  \draw[arr] (symr.east) -- node[above] {熵编码} (12.7,2.2);
  \draw[BookRule,line width=0.6pt] (-0.2,1.2) -- (15.6,1.2);
  \node[anchor=west] at (-0.2,0.6) {解码端：先从码流恢复符号};
  \node[block] (dm) at (1.4,-0.2) {$\widehat{\bm m}_t$};
  \node[block] (decm) at (5.2,-0.2) {运动解码};
  \node[block] (comp) at (9.0,-0.2) {运动补偿};
  \node[circle,draw=BookTeal,fill=BookTeal!8,line width=0.8pt,
    minimum size=6mm,inner sep=0pt] (sum) at (12.2,-0.2) {$+$};
  \node[block] (out) at (14.4,-0.2) {$\widehat{\bm Y}_t$};
  \draw[arr] (dm) -- (decm);
  \draw[arr] (decm) -- node[above] {$\widehat{\bm u}_t$} (comp);
  \draw[arr] (comp) -- node[above] {$\overline{\bm Y}_t$} (sum);
  \draw[arr] (sum) -- (out);
  \node (ref) at (9.0,0.7) {$\widehat{\bm Y}_{t-1}$};
  \draw[arr] (ref) -- (comp);
  \node[block] (dr) at (1.4,-1.7) {$\widehat{\bm z}_t$};
  \node[block] (decr) at (5.2,-1.7) {残差解码};
  \draw[arr] (dr) -- (decr);
  \draw[arr] (decr.east) -- node[above,pos=0.4] {$\widehat{\bm E}_t$}
    (12.2,-1.7) -- (sum.south);
  \node[align=center] (buffer) at (14.4,-1.7) {保存为\\下一帧参考};
  \draw[arr] (out.south) -- (buffer.north);
\end{tikzpicture}

  \caption{DVC式单参考帧间编码的主要依赖，根据原论文的计算职责重绘。上路在编码端形成量化运动与残差符号并写入码流，下路说明独立解码器如何从码流重建。两端运动补偿均使用解码后的运动和上一重建帧，编码器还须执行与下路相同的本地解码以更新参考；原始当前帧只在编码端可见。}
  \label{fig:vision-dvc-reference}
\end{figure*}
\FloatBarrier

用四个像素表示横向移动的方块。上一重建为$(0,1,1,0)$，向右移动一个像素、左边以零填充的预测为$(0,0,1,1)$。设当前源帧为$(0.2,0,1,1)$，左侧新进入的背景并非零，运动预测无法解释它，残差为$(0.2,0,0,0)$。若残差解码得到$(0.18,0,0,0)$，当前重建就是$(0.18,0,1,1)$，四像素均方误差为$0.02^2/4=10^{-4}$。这是省略了补偿网络修正的教学简化。

下一帧要使用含0.18的重建参考，而不能使用编码端掌握的原值0.2。否则，两端即使拿到相同的运动和残差，也会从不同起点产生不同图像，误差沿参考链继续传播。这就是编码器内部必须运行一遍本地解码的原因。首帧需要另行用帧内方式编码，DVC的帧间网络没有消除这个起点成本。

\subsection{运动与残差的编码}
\label{sec:vision-video-rate-allocation}

确定了参考，第二步要把运动与残差两类信息变成比特，并决定两者各花多少。输入是当前帧、参考帧和由此得到的运动与残差，输出是两段码流。难点在于运动越粗，运动比特越少，但预测越差，残差比特越多；只优化其中一方会把成本推给另一方。新显露区域等帧间预测失效的地方，还要决定是否改用帧内预测。

混合编码先确定在什么粒度上描述运动。H.264以\term{宏块}（macroblock）为基本单位，一个宏块覆盖$16\times16$的亮度样本和两个$8\times8$的色度块（4:2:0采样，见\ref{sec:vision-videofound-formats}节）\scite{vision-wiegand2003h264}。P宏块可以划分为$16\times16$、$16\times8$、$8\times16$或$8\times8$的预测块，$8\times8$块还可以再分成$8\times4$、$4\times8$或$4\times4$，于是一个宏块最多携带16个运动矢量。HEVC把宏块换成大小由编码器选择的\term{编码树单元}（coding tree unit, CTU），亮度块边长取16、32或64，再按四叉树递归划分成编码单元；预测块从$64\times64$到$4\times4$，变换块从$4\times4$到$32\times32$\scite{vision-sullivan2012hevc}。大块适合平坦背景，小块适合运动边界，划分本身也要花比特，这正是比特分配的第一个层面。

运动矢量由编码器搜索得到，常用\ref{sec:vision-videofound-block-matching}节的块匹配。H.264的运动矢量精度为四分之一像素：半像素位置用六抽头滤波器在水平和垂直方向插值，四分之一位置再由整数和半像素位置的样本平均得到；HEVC改用七抽头或八抽头滤波器。运动矢量本身先由相邻块的矢量预测，只传差值，HEVC还提供合并模式，直接继承相邻块的运动。新显露或运动复杂的区域可以改用帧内预测，用同一帧中已解码的相邻样本沿某个方向外推：H.264的$4\times4$帧内预测有9种模式、$16\times16$有4种，HEVC有33种方向模式加平面和直流两种模式。

预测残差经二维变换去掉空间相关。这一步与JPEG的$8\times8$离散余弦变换和量化表（\ref{sec:vision-restoration-jpeg-transform}节）作用相同，变换本身的定义见\ref{sec:vision-transforms-dct}节；视频标准的不同在于变换作用于预测残差而非像素，并改用整数近似。H.264主要使用$4\times4$的可分离整数变换，变换矩阵为
\begin{equation}
  \bm{T}=\begin{pmatrix}
    1&1&1&1\\
    2&1&-1&-2\\
    1&-1&-1&1\\
    1&-2&2&-1
  \end{pmatrix},\qquad
  \bm{C}=\bm{T}\bm{R}\bm{T}^{\top}\text{，}
  \label{eq:vision-video-h264-transform}
\end{equation}
其中$\bm{R}$是$4\times4$的残差块，$\bm{C}$是变换系数。$\bm{T}$的各行与$4\times4$离散余弦变换的基向量形状相近，但元素都是小整数，正反变换只需加法和移位；各行范数不同（2或$\sqrt{10}$），归一化并入量化时的缩放。一维教学例：一行线性增加的残差$(10,12,14,16)^\top$，乘以$\bm{T}$得到$(52,-14,0,-2)^\top$，能量几乎全集中在前两个系数，后两个取整后可为零。HEVC把整数变换扩展到$32\times32$，并对帧内亮度$4\times4$残差另行规定一种由离散正弦变换导出的整数变换。

变换系数按\term{量化参数}（quantization parameter, QP）均匀量化。H.264的量化参数有52个取值，每增加1，量化步长约增大12\%，每增加6，步长恰好加倍；按概述的经验，步长增大约12\%时码率约降低12\%\scite{vision-wiegand2003h264}。量化后的系数按之字形顺序扫描（与JPEG相同，见\ref{sec:vision-restoration-symbolization}节），再熵编码。H.264提供两种熵编码：\term{上下文自适应变长编码}（context-adaptive variable-length coding, CAVLC）按已传输的语法元素在多张变长码表之间切换，与JPEG用固定Huffman码表（\ref{sec:vision-restoration-huffman}节）相比多了上下文自适应；\term{上下文自适应二进制算术编码}（context-adaptive binary arithmetic coding, CABAC）先把每个语法元素二值化，再用由已编码元素构成的上下文选择概率模型，以算术编码输出比特\scite{vision-cabac2003}。算术编码的区间划分原理见\ref{sec:vision-restoration-entropy-coding}节；它能为一个符号分配非整数比特，概率还能随统计变化自适应。概述报告，CABAC相对CAVLC通常节省5\%到15\%的码率；HEVC只保留了改进后的CABAC。

这些模块都提供了多种选择：块怎样划分，用帧间还是帧内预测，参考哪一帧，运动矢量取哪个。编码器用\term{率失真优化}（rate-distortion optimization）在它们之间选择：对每个候选模式计算失真$D$和码率$R$，选使拉格朗日代价
\begin{equation}
  J=D+\lambda R
  \label{eq:vision-video-lagrangian}
\end{equation}
最小的那个\scite{vision-rdo-video1998}。$\lambda>0$是拉格朗日乘子，表示多花一比特愿意换回多少失真下降。这一做法的依据是约束优化中的乘子法：在码率预算下最小化失真，可以化为对某个$\lambda$求无约束的$D+\lambda R$的最小值，乘子的一般理论可参考\BookChapterRef{chap:optimization}{卷一的最优化章}。条件是所求的工作点位于全部可选（码率，失真）点的下凸包上；凸包以内的点无论取什么$\lambda$都选不到。由于各块的$D$和$R$可以相加，整帧的代价分解为逐块代价之和，于是可以逐块独立选择。率失真函数本身的信息论刻画卷一和卷二都没有展开，本章只使用这种可操作的形式。算法\ref{alg:vision-video-mode-decision}写出逐块选择的步骤。

\begin{algorithm}[基于拉格朗日代价的块模式选择]
  \label{alg:vision-video-mode-decision}
  \begin{algorithmic}[1]
    \Require 当前块$\bm{B}$；参考帧与已解码的相邻样本；候选模式集合$\mathcal{M}$；乘子$\lambda>0$
    \Ensure 所选模式$m^\star$
    \For{每个候选模式$m\in\mathcal{M}$}
      \State 按模式$m$形成预测$\overline{\bm{B}}_m$（帧内外推或按运动矢量补偿）
      \State 对残差$\bm{B}-\overline{\bm{B}}_m$做变换、量化和反变换，得到重建$\widehat{\bm{B}}_m$
      \State $D_m\gets\lVert\bm{B}-\widehat{\bm{B}}_m\rVert_2^2$；$R_m\gets$模式、运动矢量与残差系数的比特数
      \State $J_m\gets D_m+\lambda R_m$
    \EndFor
    \State $m^\star\gets\argmin_{m\in\mathcal{M}}J_m$
  \end{algorithmic}
\end{algorithm}

用一个$16\times16$块复算，数值为教学构造。四个候选的（失真，比特）分别为：跳过（直接复制参考，400，1）、单个$16\times16$运动矢量（120，30）、四个$8\times8$运动矢量（60，90）、帧内预测（50，150）。取$\lambda=4$，代价依次为404、240、420、750，选单个运动矢量；取$\lambda=0.5$，代价依次为400.5、135、105、125，改选四个$8\times8$运动矢量。$\lambda$变小意味着更看重失真，编码器愿意为更细的运动多花比特。这个$\lambda$与量化参数需要配套。在最优工作点附近，$\lambda$等于率失真曲线斜率的绝对值，即多一比特能减少的失真；码率越低，曲线越陡，所以粗量化应配较大的$\lambda$。

DVC把上述手工模块逐个换成网络。光流网络是\ref{sec:vision-video-warping}节的SPyNet，以预训练权重初始化，在率失真目标下一起微调。运动编码器借用学习式图像压缩的自编码器结构：四层$3\times3$、步长2、128通道的卷积，前三层后各接广义除法归一化（定义见\ref{sec:vision-restoration-quantization}节），把$H\times W\times2$的光流变成$H/16\times W/16\times128$的$\bm{m}_t$；运动解码器以三层带逆归一化的转置卷积和一层输出2通道的转置卷积对称地恢复光流。运动补偿网络把取样后的帧、参考帧和解码光流拼成8通道输入，先经$3\times3$、64通道的卷积，再经两次池化和两次上采样组成的编码—解码结构，其中共有6个预激活残差块，下采样一侧的特征经跳跃连接送到上采样一侧，最后由两层$3\times3$卷积输出3通道的预测帧。残差编码器、解码器和估计码率用的概率模型沿用尺度超先验方法（\ref{sec:vision-restoration-hyperprior}节）。表\ref{tab:vision-video-dvc-layers}汇总这些结构；原文报告整个模型约有1100万个参数。

\begin{table*}[tbp]
  \centering
  \small
  \begin{tabularx}{\linewidth}{@{}>{\raggedright\arraybackslash}p{22mm}>{\raggedright\arraybackslash}X>{\raggedright\arraybackslash}X>{\raggedright\arraybackslash}p{17mm}>{\raggedright\arraybackslash}X>{\raggedright\arraybackslash}X@{}}
    \toprule
    阶段或层组 & 层类型 & 卷积核与步长 & 输出通道 & 重复次数 & 输出尺寸\\
    \midrule
    输入 & 当前帧与上一重建帧 & — & 3 & 2 & $H\times W$\\
    光流估计（SPyNet） & 金字塔逐级残差卷积 & $7\times7$，步长1 & 32, 64, 32, 16, 2 & 每级5层，共5级 & $H\times W\times2$\\
    运动编码器 & 卷积，前三层后接广义除法归一化 & $3\times3$，步长2 & 128 & 4 & $H/16\times W/16\times128$\\
    量化 & 取整；训练时加均匀噪声 & — & 128 & 1 & $H/16\times W/16\times128$\\
    运动解码器 & 转置卷积，前三层后接逆归一化 & $3\times3$，步长2 & 128，末层2 & 4 & $H\times W\times2$\\
    运动补偿：入口 & 拼接与卷积 & $3\times3$，步长1 & 64 & 1 & $H\times W\times64$\\
    运动补偿：主体 & 预激活残差块；池化与上采样 & 论文未给出块内配置 & 64 & 残差块6个，池化与上采样各2次 & 回到$H\times W$\\
    运动补偿：出口 & 卷积 & $3\times3$，步长1 & 64，3 & 2 & $H\times W\times3$\\
    残差编码器与解码器 & 尺度超先验的非线性变换 & 论文未给出 & 论文未给出 & 论文未给出 & 与原图同尺寸的残差\\
    码率估计 & 尺度超先验的概率模型 & 论文未给出 & — & — & 每个符号的概率\\
    \bottomrule
  \end{tabularx}
  \caption{DVC的结构，按原论文正文、运动编码器图和补充材料中的运动补偿网络图整理。光流网络的层配置取自SPyNet原论文，DVC只说明以其预训练权重初始化。残差编码器和概率模型沿用尺度超先验方法，DVC未给出所用的通道数；池化的类型与跳跃连接的合并方式，图中也没有标明。原文报告全模型约1100万个参数。}
  \label{tab:vision-video-dvc-layers}
\end{table*}

量化与熵编码沿用\ref{sec:vision-restoration-quantization}节和\ref{sec:vision-restoration-entropy-coding}节的做法：训练时以加均匀噪声代替取整以便求梯度，使用时真正取整为整数符号，再按共享的概率模型熵编码。DVC用学习的概率模型估计两类符号的负对数概率，作为码率$R_{\mathrm{motion},t}$与$R_{\mathrm{residual},t}$；残差的概率模型带有尺度超先验，其辅助信息的比特也计入残差码率，计算方式见式\eqref{eq:vision-hyperprior-rate}。接收端先把比特还原成相同符号，再运行运动和残差解码网络。潜在张量的元素数不等于码流大小，实际长度还受概率分布和编码格式影响。

与混合编码逐块用式\eqref{eq:vision-video-lagrangian}选择模式不同，DVC把率失真权衡放进训练目标。对一帧使用
\begin{equation}
  L_t=\lambda D(\bm{Y}_t,\widehat{\bm{Y}}_t)
       +R_{\mathrm{motion},t}+R_{\mathrm{residual},t}\text{，}
  \label{eq:vision-video-rate-distortion}
\end{equation}
其中$D$在原实现中取均方误差，$\lambda>0$控制质量的权重。把式\eqref{eq:vision-video-lagrangian}两边同除以$\lambda$可见，两者是同一个拉格朗日代价，只是乘子写在失真一侧，所以这里$\lambda$越大，失真受到的惩罚越重，网络倾向于多花比特。原文分别取$\lambda=256,512,1024,2048$训练四个模型，得到四个码率工作点；编码时不再逐块搜索，权衡已经固化在网络参数中。若$D$按像素平均、$R$按每帧比特，$\lambda$的尺度与改用比特/像素时不同，比较时应固定归一化。光流网络也受这个目标约束，因而它学到的运动要兼顾可压缩性与预测用途，不应仅以光流端点误差评价它。

以教学数值说明联合目标的必要性。假设同一帧的方案A用20比特描述运动、80比特描述残差，方案B用10比特描述运动、105比特描述残差，总数分别为100与115。运动分支省了10比特，整个系统却多花15比特，因为粗运动把更多误差留给了残差。这个对照不假设两方案失真相同；若失真不同，还要一起报告。

原文的消融实验给出了联合训练的实际效果。若在整个训练中固定预训练的光流网络，编码光流需要0.044比特/像素，取样所得预测帧的PSNR为27.33 dB；联合训练后只需0.029比特/像素，PSNR反而升到28.17 dB，即学到的运动既更省比特，又更利于预测。去掉运动自编码器、直接量化原始光流，光流要花0.20比特/像素，预测帧PSNR降到约24.3 dB，整体码率—质量曲线在同等码率下下降超过2 dB；去掉运动补偿网络、只用取样结果作预测，同等码率下PSNR约降1.0 dB。在其报告的$\lambda=512,1024,2048$三个工作点上，运动信息约占总码率的36\%、35\%和27\%，质量要求越高，占比越低\scite{vision-dvc2019}。

训练还要处理参考依赖：第$t$帧的输入依赖第$t-1$帧的重建，而后者又依赖更早的帧。DVC采用在线更新的缓存：把每个训练片段上一步的重建帧存入缓存，下一步直接读取作参考，避免每次从序列首帧回算整条链；原文报告，用原始前一帧代替缓存中的重建帧训练，同等码率下约损失0.2 dB。训练数据是Vimeo-90K的短片段，$256\times256$裁剪，批量4，Adam优化，学习率从$10^{-4}$开始在损失停滞时降为十分之一，在两块Titan X上约训练7天。缓存只是训练组织方式；真正编码时仍要顺着当前参数形成的重建链运行，独立解码器不能依赖训练缓存或原视频。

两条路线在比特分配上的差别可以归结为三点。第一，运动的表示：混合编码用块和运动矢量，划分本身要花比特，但描述紧凑；DVC用逐像素光流再压缩成潜在表示，没有块边界，比特由网络自行分配。第二，权衡发生的时刻：混合编码在编码时对每块搜索模式，可以随内容和码率目标调整；DVC在训练时固化权衡，一个$\lambda$对应一个模型，编码时只做一次前向计算。第三，变换的适应性：混合编码的整数变换固定，DVC的非线性变换从数据学到。DVC因此可以联合优化各模块，代价是缺少逐块自适应的模式选择，码率点也只有训练过的几个。

完整视频应统计首帧或关键帧、各帧运动、残差以及实际传输的尺寸、模式等辅助信息。估计码率用于训练，真实码流字节数用于核验传输成本。模型若预先部署在两端，需说明模型版本；若传输中更新模型或概率参数，那些字节也属于必要成本。

\subsection{视频帧的解码重建}
\label{sec:vision-video-decoding}

第三步在接收端按依赖顺序还原每一帧。输入只有码流和事先共享的模型，输出是重建序列。难点在于任何一帧的解码都依赖它的参考，参考出错或缺失会影响后续所有帧，而用户可能想从中间开始观看。

混合编码的解码器按码流顺序解析每块的模式、运动矢量和系数，反量化、反变换后加回预测，再做环路滤波。块状处理在块边界上留下的不连续是最显眼的压缩伪影。H.264规定了自适应的环路\term{去块滤波}（deblocking filter）\scite{vision-wiegand2003h264}：在块边界两侧比较样本差异，若边界处差异明显而又小于随量化参数变化的阈值，就认为它是量化造成的块效应并加以平滑；若差异大到量化的粗细解释不了，就认为是图像中的真实边缘而不处理。因为滤波发生在预测环路之内，滤波后的帧也是后续帧的参考，改善的质量可以传给后续预测；概述报告，在相同客观质量下它通常节省5\%到10\%的码率。HEVC在去块之后又加入\term{样点自适应偏移}（sample adaptive offset, SAO），用少量编码器统计出的参数，以查表方式修正样本幅值\scite{vision-sullivan2012hevc}。

DVC的独立解码器先恢复图像组的首帧，再按参考关系读取后续帧的运动与残差符号，用解码网络得到$\widehat{\bm{v}}_t$和$\widehat{\bm{E}}_t$，执行运动补偿与式\eqref{eq:vision-video-dvc-reconstruction}的相加，然后把当前结果写回参考缓存。原文认为逐像素的运动补偿不产生块效应，因而不需要手工设计的环路滤波和样点自适应偏移；按原文的设计，取样结果中的伪影由运动补偿网络去除，它承担了类似的修正职责。所有运算的尺寸、裁剪、量化和数值约定都应与编码器的本地解码一致。这里出现了混合编码已经解决、学习式编码重新面对的问题：神经网络在不同硬件上的浮点运算结果可能有末位差异，两端的参考于是不完全相同，正是H.264用精确整数反变换避免的漂移，需要在部署时检查。只有从保存的码流独立重建成功，才能说明输出是一套可用的编码结果。

单参考链不能从任意中间帧直接解码，随机访问要从所在图像组的首帧开始恢复依赖。图像组越短，独立帧越多，码率越高，但随机访问等待和错误传播范围越小。DVC的原实验在不同测试集上分别使用长度12和10的图像组。传输丢包还可能破坏后续参考，不能只检查每帧单独的解码函数。多参考或双向预测会改变依赖、缓冲、解码顺序和延迟条件，应作为不同编码配置比较。

DVC原文报告，在其测试设置下，按PSNR和MS-SSIM计算的码率—质量曲线在多数数据集上优于其所配置的H.264编码器，按MS-SSIM与H.265相近或更好\scite{vision-dvc2019}。这一比较对应原文的具体配置：H.264和H.265由FFmpeg的very fast预设生成，图像组长度如上；补充材料按序列给出了相对H.264的BD码率，例如在BasketballDrive序列上，按PSNR计算H.265为$-44.37\%$，DVC为$-23.17\%$。编码标准名不能替代具体编码器、预设、图像组长度和速率控制配置，更不能据此给出当前持续有效的排名。速度也取决于实现：在$352\times288$的视频上，原文测得DVC编码约24.5帧每秒、解码约41帧每秒（单块1080Ti），官方参考软件JM和HM的编码速度分别为2.4和0.35帧每秒，而高度优化的x264和x265分别达到250和42帧每秒。

DVC与BasicVSR都依赖光流和反向取样，比较二者能看清编码问题的特殊之处。在BasicVSR中，光流在使用端由输入帧随时计算，不需付出比特，可以尽量精确；在DVC中，光流必须编码传输，精度要用比特换取，运动网络因此学习可压缩的运动。BasicVSR的参考是低质量输入，误差不会沿参考链放大；DVC的参考是自己的重建，每帧的编码误差都会成为下一帧的起点。前者以质量为唯一目标，后者以率失真为目标，评价时也必须同时报告码率。

\section{处理效果的核验}
\label{sec:vision-video-restoration-evaluation}

一个结果可能逐帧更清楚，却在播放时闪烁；另一个可能节省比特，却要等待整段视频或改变字符。本节要解决的问题是：在什么条件下，两份视频处理结果的比较能说明方法本身的差别，以及“处理成功”需要哪些证据。\ref{sec:vision-video-restoration-challenges}节已经给出了指标的定义，这里关心的是怎样使用它们。核验有三个维度，分别针对不同的失败来源：比较协议排除输入和信息条件不同造成的假差异；内容与时间连续性检查单帧指标看不到的时序错误；质量、码率与延迟则把收益与成本放在一起，按用途选择。

\subsection{比较协议的核验}
\label{sec:vision-video-restoration-protocol}

比较前要确认双方面对的是同一个问题。超分需要固定倍率、下采样与模糊方式、裁边和色彩空间，BasicVSR的两种退化就要分开报告；插帧需要固定两端间隔、目标时刻以及是否以先前生成帧作输入；压缩需要固定输入分辨率、帧率、图像组长度、参考结构和码流统计方式。可见帧范围也属于协议：一个方法看整段，另一个只看过去，单帧质量差异就同时包含信息条件差异。训练与测试划分要按源视频进行，相邻帧或近重复片段落入不同划分会泄漏信息。

指标的汇总方式同样属于协议。PSNR是逐帧计算再平均，还是先平均均方误差再换算，按\ref{sec:vision-video-restoration-challenges}节的算例可以相差数个分贝；只在亮度通道上计算还是在RGB上计算，也会改变数值，BasicVSR的原文表格就对不同数据集采用了不同通道。码率按式\eqref{eq:vision-video-bitrate}由完整码流字节数计算，不能用潜在张量的元素数或训练时的估计码率代替。

计算延迟、等待后帧的延迟和显示缓冲要按式\eqref{eq:vision-video-latency}分别测量。吞吐量描述每秒处理多少帧，延迟描述某帧从拍摄到输出等待多久；报告时还需给出硬件、数值精度、批量、输入长度和峰值内存，避免只比较网络某个局部模块的耗时。DVC与x264的速度比较就说明了这一点：同一种编码思路，参考软件与高度优化的商业实现可以相差两个数量级，速度数字只对具体实现有效。

\subsection{内容与时间连续性的核验}
\label{sec:vision-video-temporal-quality}

协议统一后，要检查输出内容是否正确、播放是否连续。单帧指标PSNR、SSIM和LPIPS仍然适用，但它们逐帧计算再平均，看不到闪烁、重影、拖尾和字符随时间变化。直接惩罚相邻输出之差$\widehat{\bm{X}}_t-\widehat{\bm{X}}_{t-1}$也不行，因为正常运动时相邻帧本来就不同，这样会鼓励画面停滞；时间比较应补偿真实运动，并排除遮挡、新显露和越界区域，这正是式\eqref{eq:vision-video-temporal-error}的扭曲误差的设计。

使用扭曲误差时要说明三件事：对应和掩码来自哪里，掩码怎样得到（例如由式\eqref{eq:vision-video-fb-consistency}的前后向一致性），以及按什么范数和归一化计算。若对应来自被评价方法自己的光流，误差会被方法本身掩盖；光流不准、曝光变化或掩码漏判也会影响它。扭曲误差和tOF回答的问题不同：前者问相邻输出在运动补偿后是否一致，后者问输出的运动是否与真值的运动一致，需要配对真值。两者应与逐帧指标一起报告，因为任何一个单独都可以被投机：模糊输出的扭曲误差低；逐帧补出彼此不一致细节的输出，逐帧感知距离可能较低，播放时却闪烁。

时间差小也不保证内容正确。每帧都输出同一张模糊灰图，可以很稳定却丢失全部细节；每帧都把车牌同一字符改错，也能保持时间一致。因此要同时检查与参考内容的关系及实际用途：逐帧读取车牌，回查原观测，比较识别置信度和字符串随时间的变化，能识别“单帧好看、播放时字符反复变换”的失败。有配对清晰视频时，可比较输出与真值的运动补偿差异；没有真值时，主观连续性、下游任务和原始证据检查更重要，但不能把这些替代检查说成已证明像素恢复正确。

发现失败后，应追溯到具体环节：取样位置错误（对齐）、错误内容被赋予高权重（聚合）、错误沿状态扩散（传播）、两侧都未见的区域由先验补出（插帧合成），或参考链失配（编码）。本章各方法在这些环节上的设计不同，同一种表面现象可能来自不同原因：同样是残影，在运动补偿时域滤波中多半来自过大的时间权重，在RVRT中可能来自错误特征被逐片段传播，在DVC中则可能来自两端参考不一致造成的漂移。

\subsection{质量、码率与延迟的核验}
\label{sec:vision-video-quality-rate-latency}

最后要在收益与成本之间作出选择。单个工作点无法比较两种方法：一个码率更低、质量也更低的结果，可能只是选了不同的$\lambda$或量化参数。正确做法是固定源序列，用多个配置分别形成质量—码率曲线和质量—等待曲线，在相同码率下比较质量，或在相同质量下比较码率；式\eqref{eq:vision-bd-rate}的BD码率把这种比较概括为一个数，但它只在两条曲线重叠的质量区间内有效，曲线交叉时一个平均数会掩盖“低码率时更好、高码率时更差”的情况，应同时给出曲线。若一个配置更省码率却必须等待整段视频，另一个略差但能即时输出，选择取决于用途：观看容忍适量细节损失，车牌读取更在意字符，测量还要求原始观测和时间坐标可追溯。

表\ref{tab:vision-video-tradeoff}汇总不同处理职责最容易遗漏的条件，说明同一“清晰视频”目标为何需要多项核验。比较质量、码率与延迟时，应同时保留对应、可见性、真实细节和解码依赖的检查。

\begin{table}[htbp]
  \centering
  \begin{tabularx}{\linewidth}{@{}lXXX@{}}
    \toprule
    处理职责 & 主要新增信息 & 必须同时检查 & 主要成本\\
    \midrule
    多帧复原 & 邻帧互补内容 & 对应、可见性与错误传播 & 等待范围与缓存\\
    空间超分 & 位移形成的采样差异 & 真实细节与先验生成细节 & 计算量与等待\\
    时间插帧 & 两端的时间约束 & 后帧等待、遮挡与运动歧义 & 后帧等待\\
    帧间压缩 & 已解码参考的可预测内容 & 码流、参考一致与随机访问 & 码率、重排等待与解码依赖\\
    \bottomrule
  \end{tabularx}
  \caption{视频处理的不同收益依赖不同信息条件。比较应固定输入与用途，把右侧两列与单帧质量一起检查，不能据某一项改善推断整个系统已满足要求。}
  \label{tab:vision-video-tradeoff}
\end{table}
\FloatBarrier

\section{本章小结}
\label{sec:vision-video-restoration-summary}

开篇的问题是怎样利用其他时刻的路牌观测得到可读、可放大、可补帧、可传输的视频。本章先把这个问题拆成五项挑战：对齐误差带来伪影，遮挡使邻帧缺少所需内容，时间一致性与逐帧质量相互牵制，读取范围受等待约束，编码还要在运动与残差之间分配比特。回答都围绕帧间对应展开。邻帧只有在对应正确、内容可见时才有用：反向取样决定到哪里读取，前后向一致性筛掉不可见位置，权重决定读来的信息怎样参与，状态传播决定哪些历史仍能影响后续输出。经典的运动补偿时域滤波用按局部差异调整的时间权重做递归平均，其方差公式把降噪强度与残影持续时间联系在一起；VBM4D沿轨迹把相似图块分组协同收缩；RVRT以片段内联合处理和光流引导的多位置注意力，在逐帧循环与整段并行之间取得折中，并把手工的相似度和阈值换成学到的注意力。

输出网格决定了结果中有多少是观测、多少是估计。原网格恢复每个位置都有退化观测；BasicVSR用双向传播取得整段的亚像素观测，再由像素重排扩展空间网格，新增像素依赖邻帧的间接观测；FILM从未观测的中间时刻出发向两侧取样，输出完全是估计。编码则要求共享信息能在解码端重现：H.264/HEVC的混合编码以块运动补偿、整数变换、量化参数和上下文自适应熵编码组成，编码器用拉格朗日代价逐块选择模式，编码器内部复制解码环路，精确整数反变换消除了两端漂移；DVC保留这一数据流，把运动估计、运动压缩、运动补偿和残差变换换成网络，用率失真损失联合训练，使学到的运动更便于压缩，代价是权衡固化在模型中，也重新面对浮点运算的两端一致问题。

本章没有解决的部分同样来自这些依赖。双向方法需要等待整段视频，改成在线要重新设计和验证；大运动、遮挡、切镜和真实时变退化仍会使对应失效，错误还可能沿状态和参考链扩散；插帧和超分中由先验补出的细节无法从输出本身区分出来；极低码率下，运动—残差框架能保留的信息有限，学习式编码的模型部署与跨硬件一致也还需逐项检查。因此，画面更锐利、更平滑或码流更小，都不足以单独证明任务完成；改进需要在相同输入条件、可见范围和用途下，同时核验逐帧质量、时间一致性、码率与延迟。

\paragraph{扩展阅读}
跨帧对齐与传播方向，可以对照本章以光流取样为主的路线阅读几项工作。TDAN用可变形卷积在特征层面直接学习采样偏移，不先显式估计光流；EDVR在滑动窗口内用金字塔级联的可变形对齐和时空注意力融合处理邻帧，是循环方法普及前窗口路线的代表；BasicVSR++在BasicVSR的双向传播上加入二阶网格传播和光流引导的可变形对齐，也就是本章与RVRT对比时提到的后续版本\scite{vision-tdan2020,vision-edvr2019,vision-basicvsrpp2022}VRT是RVRT作者此前提出的整段并行模型，用时间互注意力在帧间交换信息，可用来具体检查并行路线的融合收益与长序列内存代价；FastDVDnet则用两级级联的U形网络直接读取五帧，不做显式运动估计，与依赖对应的经典滤波和VBM4D形成对照\scite{vision-vrt2024,vision-fastdvdnet2020}

插帧方向，Super SloMo把两帧间的双向光流线性组合成中间时刻的近似光流，并预测可见性图处理遮挡，可以一次插出任意多个中间帧；DAIN在光流投影中加入深度，让较近的物体在遮挡处优先；RIFE用一个网络直接估计中间时刻的光流，不再从两帧间光流缩放，换取实时速度\scite{vision-superslomo2018,vision-dain2019,vision-rife2022}三者都可与FILM比较对大运动和去遮挡区域的处理方式。编码方向，H.266/VVC的综述介绍了HEVC之后的混合编码标准怎样加入更灵活的块划分和更多预测工具，并扩展到屏幕内容、全景视频等应用；学习式编码中，FVC把运动估计、压缩与补偿从像素空间移到特征空间，DCVC则不再从预测帧中减出残差，而把时间上下文特征作为编码、解码与熵模型的条件，二者分别从运动表示和残差形式上改变了DVC的数据流\scite{vision-vvc2021,vision-fvc2021,vision-dcvc2021}

